\documentclass[11pt,a4paper]{report}

\usepackage[margin=2.4cm]{geometry}
\usepackage{amsmath,amssymb,amsthm,mathtools,bm}
\usepackage{braket}
\usepackage{microtype}
\usepackage[hidelinks,hypertexnames=false]{hyperref}
\usepackage{algorithm}
\usepackage{algpseudocode}
% \LState{...}: a pseudocode statement whose wrapped lines keep the statement's indentation
\makeatletter
\algnewcommand{\LState}[1]{\State\parbox[t]{\dimexpr\linewidth-\ALG@tlm\relax}{\raggedright\strut#1\strut}}
\makeatother
\usepackage{booktabs}
\usepackage{array}
\usepackage{longtable}
\usepackage{enumitem}
\usepackage[most]{tcolorbox}
\usepackage{xcolor}
\usepackage{pgfplots}
\pgfplotsset{compat=1.17}
\usepackage{float}

\definecolor{boxbg}{RGB}{245,247,250}
\definecolor{boxframe}{RGB}{70,90,120}
\newtcolorbox{keybox}[1]{colback=boxbg,colframe=boxframe,title=#1,fonttitle=\bfseries,
  boxrule=0.6pt,arc=1.5mm,left=2mm,right=2mm,top=1.5mm,bottom=1.5mm}
\newtcolorbox{gpubox}[1]{colback=green!4,colframe=green!35!black,title=#1,fonttitle=\bfseries,
  boxrule=0.6pt,arc=1.5mm,left=2mm,right=2mm,top=1.5mm,bottom=1.5mm}

\theoremstyle{plain}
\newtheorem{theorem}{Theorem}[chapter]
\newtheorem{lemma}[theorem]{Lemma}
\theoremstyle{definition}
\newtheorem{definition}[theorem]{Definition}
\theoremstyle{remark}
\newtheorem{remark}[theorem]{Remark}

\newcommand{\rr}{\mathbf{r}}
\newcommand{\RR}{\mathbf{R}}
\newcommand{\kk}{\mathbf{k}}
\newcommand{\GG}{\mathbf{G}}
\newcommand{\qq}{\mathbf{q}}
\newcommand{\dd}{\mathrm{d}}
\newcommand{\ii}{\mathrm{i}}
\newcommand{\ee}{\mathrm{e}}
\newcommand{\Ha}{\hat{H}}
\newcommand{\Exc}{E_{\mathrm{xc}}}
\newcommand{\vxc}{v_{\mathrm{xc}}}
\newcommand{\veff}{v_{\mathrm{eff}}}
\newcommand{\vext}{v_{\mathrm{ext}}}
\newcommand{\vH}{v_{\mathrm{H}}}
\newcommand{\fdv}[2]{\frac{\delta #1}{\delta #2}}
\newcommand{\pdv}[2]{\frac{\partial #1}{\partial #2}}

\allowdisplaybreaks
\makeatletter\setlength{\@fptop}{0pt}\makeatother % float pages start at the top
\hyphenation{pseudo-po-ten-tial pseudo-po-ten-tials}
\setlength{\parskip}{0.35em}

\title{\textbf{Density Functional Theory from Scratch}\\[0.6em]
\Large A complete theoretical and algorithmic guide to a plane-wave,
spinor, spin--orbit-coupled DFT code for FCC gold --- its bands, its
topology, its conductivity and its colour\\[0.6em]
\large Designed for a JAX / GPU implementation\\[0.6em]
\large \textbf{Written by Claude AI}}
\author{}
\date{July 2026; revised and numerically verified November 2026}

\begin{document}
\maketitle
{\setlength{\parskip}{0pt}\tableofcontents}

\chapter{Introduction: What We Are Building and Why Gold}
\label{ch:intro}

\section{The goal}

You are going to write, from scratch, a program that solves the quantum
mechanics of a real crystalline solid. The input will be almost embarrassingly
small: the chemical species (gold), the crystal structure (face-centered
cubic), and a lattice constant. From this the program will predict, with no
empirical fitting whatsoever, the electronic band structure, the total energy,
the equilibrium lattice constant, the bulk modulus, the density of states,
a topological invariant of bulk gold, the ingredients of its electrical
conductivity, and its reflectivity spectrum and colour --- quantities you
can check against experiment and against
mature codes such as Quantum ESPRESSO, VASP, or ABINIT.

The framework is \emph{Kohn--Sham density functional theory} (DFT), the
workhorse of modern computational materials science. This document contains
\emph{all} the theory and \emph{all} the algorithms you need, assuming only
undergraduate quantum mechanics: the time-independent Schr\"odinger equation,
Dirac notation, Hermitian operators, the variational principle, the
hydrogen atom, spin-$\tfrac12$ algebra with Pauli matrices, and first-order
perturbation theory --- plus enough Python and NumPy to write and debug a
few hundred lines of array code (the handful of JAX features used are
explained in Chapter~\ref{ch:gpu}). Everything
else --- functional calculus, the Hohenberg--Kohn theorems, exchange and
correlation, Bloch's theorem, pseudopotentials, relativistic corrections,
iterative eigensolvers, self-consistency acceleration --- is developed here.

No program listings appear in this document, by design: writing the code is
your job, and it is how you will learn the method. Every algorithm is instead
given as explicit mathematics plus pseudocode (the handful of JAX functions
you need are named and explained in Chapter~\ref{ch:gpu}), and each
computational chapter ends
with notes on how the algorithm maps onto a massively parallel accelerator
(GPU) under JAX: which loops become batched tensor contractions, which
operations are FFTs, where \texttt{vmap}-style batching applies,
and which numerical choices (block eigensolvers, Fourier-space mixing) exist
\emph{specifically because} they parallelize well.

\section{Why gold is the right ``hard'' first material}

A first DFT code is usually tested on silicon: a light-element semiconductor
where spin can be ignored and a coarse Brillouin-zone sampling suffices. This
guide deliberately aims at something harder, and gold is the ideal choice,
because it is
\emph{physically} demanding while remaining \emph{structurally} trivial:

\begin{enumerate}[leftmargin=2em]
\item \textbf{One atom per primitive cell.} FCC gold has a single atom in the
primitive cell. Your arrays stay small, debugging stays sane, and yet nothing
essential is simplified away.

\item \textbf{It is a metal.} There is no band gap. Occupation numbers are not
simply ``empty'' or ``full''; you must locate a Fermi level, introduce fractional
(smeared) occupations, sample the Brillouin zone densely, and stabilize the
self-consistency loop against ``charge sloshing.'' These are exactly the
algorithmic complications that separate a toy code from a real one
(Chapters~\ref{ch:solids} and \ref{ch:algorithms}).

\item \textbf{The $5d$ electrons make exchange--correlation matter.} Gold's
valence configuration is $5d^{10}6s^1$. The filled, spatially compact $d$
shell sits only $\sim$2--8\,eV below the Fermi level and hybridizes with the
$6s$ band. The position of these $d$ bands --- and hence the computed optical
absorption edge, lattice constant, and cohesive energy --- is sensitive to the
exchange--correlation functional. You will implement both the local density
approximation (LDA) and the PBE generalized gradient approximation and see
their differences directly (Chapter~\ref{ch:xc}).

\item \textbf{Spin--orbit coupling is enormous and observable.} With nuclear
charge $Z=79$, the $5d$ electrons of gold move at an appreciable fraction of
the speed of light near the nucleus. Spin--orbit coupling (SOC) splits the
$d$ bands at the $\Gamma$ point by more than $1$\,eV --- not a small correction
but a dominant feature of the band structure. Famously, relativity is why
gold is yellow while silver (same column, $Z=47$) is white: relativistic
stabilization of the $6s$ band and destabilization of the $5d$ band shift the
interband absorption edge from the ultraviolet (silver, $\approx 3.8$\,eV)
down into the visible ($\approx 2.4$\,eV, blue-green light), so gold absorbs
blue and reflects red and yellow. Your code
will treat SOC non-perturbatively with two-component spinor wavefunctions
(Chapter~\ref{ch:soc}).

\item \textbf{It hides a topological invariant.} Although gold is a metal,
its bands carry a nontrivial $\mathbb{Z}_2$ topological index, and the famous
surface state on Au(111) is its fingerprint. This is not an afterthought
bolted onto the project: the index is a property of the \emph{wavefunctions}
your code produces, and every tool needed to extract it is introduced where
it naturally belongs --- the parity of a zone-boundary state in the very
first band-gap model (Section~\ref{sec:nfe}), inversion and time reversal as
operations on plane-wave coefficients (Section~\ref{sec:symmetry}), and the
fixed-potential band solver (Section~\ref{sec:nscf}).
Chapter~\ref{ch:topology} then only has to assemble them.

\item \textbf{Excellent reference data.} The experimental lattice constant
($a = 4.0782$\,\AA), bulk modulus ($\approx 173$\,GPa; both are quoted at
room temperature, while DFT describes a static lattice at $0$\,K: removing the
zero-point and thermal expansion shrinks $a$ by $\approx 0.015$\,\AA{}, to
$\approx 4.063$\,\AA, and stiffens the lattice: $B \approx 180$\,GPa from the
$0$\,K elastic constants, up to $\approx 196$\,GPa in the static-lattice
estimates used by density-functional benchmarks), the photoemission
$d$-band positions, and published DFT band structures at every level of
theory give you unambiguous targets for verification
(Chapter~\ref{ch:verify}).
\end{enumerate}

\begin{keybox}{The physical system}
FCC gold. Primitive lattice vectors (with $a = 4.0782$\,\AA{}
$= 7.706681$ Bohr; convert with the full $a_0 = 0.529177211$\,\AA{} ---
rounding $a$ to four decimals already shifts the Ewald energy by
$10^{-4}$\,Ha and fails the tests of Chapter~\ref{ch:verify}):
\[
\mathbf{a}_1 = \tfrac{a}{2}(0,1,1),\qquad
\mathbf{a}_2 = \tfrac{a}{2}(1,0,1),\qquad
\mathbf{a}_3 = \tfrac{a}{2}(1,1,0),
\]
one Au atom at the origin. With a scalar-relativistic + SOC pseudopotential
treating $5d^{10}6s^1$ (11 electrons) or $5s^25p^65d^{10}6s^1$ (19 electrons,
``semicore,'' more accurate) explicitly.
\end{keybox}

\section{Units}

We use \textbf{Hartree atomic units} throughout unless stated otherwise:
$\hbar = m_e = e = 4\pi\varepsilon_0 = 1$. Then lengths are in Bohr radii
($a_0 = 0.529177$\,\AA), energies in Hartree ($1\,\mathrm{Ha} = 27.2114$\,eV
$= 2$\,Ry), and the speed of light is $c = 1/\alpha \approx 137.036$, where
$\alpha$ is the fine-structure constant. The one-electron kinetic operator is
$-\tfrac12\nabla^2$ and the Coulomb interaction between electrons is
$1/|\rr_1-\rr_2|$. Atomic units eliminate an entire class of bugs; convert to
eV and \AA{} only at the input/output boundary of the program.

\section{Bird's-eye view of the whole machine}

It helps to see the destination before the journey. The final program is a
fixed-point iteration --- the \emph{self-consistent field} (SCF) loop ---
wrapped around a large eigenvalue problem:

\begin{enumerate}[leftmargin=2em]
\item \textbf{Guess} an electron density $n(\rr)$ (e.g.\ a normalized Gaussian
blob of valence charge on each atom).
\item \textbf{Build the effective potential} $\veff[n](\rr)$: the sum of the
ionic (pseudo)potential, the classical electrostatic (Hartree) potential of
the density itself, and the exchange--correlation potential, which encodes
all quantum many-body effects (Chapters~\ref{ch:hkks}--\ref{ch:xc}).
\item \textbf{Solve} the single-particle Schr\"odinger-like (Kohn--Sham)
equation in this potential at many points $\kk$ of the Brillouin zone, using
a plane-wave basis, obtaining bands $\varepsilon_{n\kk}$ and two-component
spinor wavefunctions $\psi_{n\kk}$ (Chapters~\ref{ch:solids},
\ref{ch:soc}, \ref{ch:hamiltonian}). This is done with an iterative,
GPU-friendly block eigensolver (Chapter~\ref{ch:algorithms}).
\item \textbf{Occupy} the bands according to a smeared Fermi--Dirac-like
distribution, locating the Fermi level so the electron count is correct.
\item \textbf{Recompute the density} from the occupied wavefunctions.
\item \textbf{Mix} the new density with the old ones (Pulay/Anderson mixing
with Kerker preconditioning --- what keeps metallic cells stable) and go to step 2 until
the density stops changing.
\item \textbf{Evaluate} the total energy, band structure, density of
states, and wavefunction-based quantities (parities, Berry phases,
velocity matrix elements for transport and optics). (Forces
on the atoms vanish by symmetry with one atom per cell, so they are not
needed here; the equilibrium geometry is found from the energy as a function
of the lattice constant.)
\end{enumerate}

Almost all floating-point work lives in step 3, and there it lives inside two
kernels: \emph{batched 3D FFTs} and \emph{batched dense linear algebra}. Both
are the native food of GPUs. That is not a coincidence --- the plane-wave
formulation of DFT was co-evolved with vector and parallel hardware, and the
algorithm choices in this document (LOBPCG over sequential Davidson,
Fourier-space Kerker mixing, Monkhorst--Pack batching over $\kk$-points) are
made with your JAX target explicitly in mind.

\section{How to use this document}

Read Chapters~\ref{ch:manybody}--\ref{ch:soc} for the physics, with pen and
paper; they contain no programming. Chapters~\ref{ch:hamiltonian} and
\ref{ch:algorithms} are the specification of the program, and
Chapter~\ref{ch:gpu} translates it into JAX. Then \emph{build in the order
of Chapter~\ref{ch:verify}}: every stage there has a number that your code
must reproduce before you move on, so that at no point are you debugging
more than one new ingredient. Three conventions will save you days:
\begin{itemize}[leftmargin=2em]
\item Boxes with a blue-grey frame summarize what must be remembered; boxes
with a green frame contain implementation advice for JAX.
\item Every formula that you will type into the computer comes with a
\emph{numerical anchor} --- a value to compare against --- either next to it
or in Chapter~\ref{ch:verify}. Use them all; a wrong constant is invisible
in a band plot.
\item The folder \texttt{verification/} distributed with this document
holds the scripts that produced every reference number, including an
independent JAX implementation written from this text alone. Treat it as
the answer key: consult it only after your own attempt.
\end{itemize}

\section{Chapter map}

\begin{itemize}[leftmargin=2em]
\item \textbf{Ch.~\ref{ch:manybody}} --- the exact many-electron problem and
why it is intractable; the Born--Oppenheimer approximation; the density as
the central variable.
\item \textbf{Ch.~\ref{ch:hkks}} --- functional calculus from zero; the
Hohenberg--Kohn theorems with proofs; the Kohn--Sham construction and a full
derivation of the Kohn--Sham equations.
\item \textbf{Ch.~\ref{ch:xc}} --- exchange and correlation: the homogeneous
electron gas, derivation of LDA exchange, the PW92 correlation
parametrization, the PBE GGA, and the (non)collinear spin generalization
needed once SOC is on.
\item \textbf{Ch.~\ref{ch:solids}} --- crystals: Bravais and reciprocal
lattices, Bloch's theorem (proved), plane-wave basis sets and cutoffs, the
two-plane-wave model of a band gap, Brillouin-zone integration,
Monkhorst--Pack meshes, metallic occupations and smearing, the Fermi level.
\item \textbf{Ch.~\ref{ch:pseudo}} --- pseudopotentials: why all-electron
plane-wave calculations are hopeless, norm conservation, the
Kleinman--Bylander separable form, and the analytic relativistic HGH
pseudopotentials you will use for gold.
\item \textbf{Ch.~\ref{ch:soc}} --- relativity: from the Dirac equation to
the Pauli Hamiltonian, the physical origin of SOC, $j$-dependent
pseudopotentials, two-component spinor Kohn--Sham theory, time-reversal and
inversion symmetry as explicit operations on your coefficient arrays, and
what SOC does to gold's bands.
\item \textbf{Ch.~\ref{ch:hamiltonian}} --- the full Hamiltonian in the
plane-wave basis, term by term: kinetic, local via FFT, nonlocal via
projectors, Hartree in reciprocal space, Ewald summation for the ion--ion
energy, and the total-energy expression.
\item \textbf{Ch.~\ref{ch:algorithms}} --- the algorithms: basis
construction, matrix-free Hamiltonian application, preconditioned block
eigensolvers (LOBPCG and blocked Davidson), subspace diagonalization,
Fermi-level bisection, density construction, Kerker-preconditioned Pulay
mixing, the complete SCF driver, and the fixed-potential solver used for
all post-processing --- all as pseudocode.
\item \textbf{Ch.~\ref{ch:gpu}} --- mapping everything onto JAX/GPU: tensor
layouts, batching over $\kk$ and bands, FFT strategy, mixed precision,
JIT-compatible control flow, and memory budgeting.
\item \textbf{Ch.~\ref{ch:topology}} --- the topology of gold's bands:
Berry phases and Wannier charge centers from scratch, the $\mathbb{Z}_2$
invariant as a statement about how those centers flow, its evaluation from
wavefunction parities and by Wilson loops, with pseudocode.
\item \textbf{Ch.~\ref{ch:conductivity}} --- electrical conductivity: the
velocity operator, the Drude plasma frequency from the Fermi surface, and
an honest account of what a ground-state calculation cannot give (the
scattering time).
\item \textbf{Ch.~\ref{ch:colour}} --- the colour of gold: the dielectric
function from interband transitions, the reflectivity, and the conversion
of a reflectance spectrum into the colour a human eye sees.
\item \textbf{Ch.~\ref{ch:verify}} --- a staged verification plan with
quantitative targets, from the free-electron gas to the SOC-split $d$ bands
of gold and its topological index, including reference numbers from an
independent implementation of exactly the equations in this document.
\item \textbf{Appendix~\ref{app:q19}} --- the semicore (19-electron)
pseudopotential parameters.
\end{itemize}

\section{Notation at a glance}
\label{sec:notation}

The same few indices appear in every chapter; fix them now.
\begin{center}
\begin{tabular}{@{}lp{0.66\textwidth}@{}}
\toprule
Symbol & Meaning\\
\midrule
$\rr$, $\RR$ & position in space; lattice translation (or atomic position
$\RR_a$)\\
$\kk$, $\GG$ & wavevector in the Brillouin zone; reciprocal lattice vector\\
$n$ (as an index), $\kk$ & band index and wavevector of a Bloch state
$\psi_{n\kk}$\\
$n(\rr)$, $\mathbf m(\rr)$ & electron density; magnetization density\\
$s$, $\alpha$, $\beta$ $\in\{\uparrow,\downarrow\}$ & spinor component\\
$c^{s}_{n\kk}(\GG)$ & plane-wave coefficients of a spinor Bloch state
(the program's central array)\\
$f_{n\kk} \in [0,1]$ & occupation of one spinor state (one electron at
most)\\
$w_{\kk}$ & Brillouin-zone weight of a $\kk$-point, $\sum_{\kk} w_{\kk} = 1$\\
$\varepsilon_{n\kk}$, $\mu$ & Kohn--Sham eigenvalue; Fermi level (chemical
potential)\\
$\Omega$, $N_e$ & volume of the primitive cell; number of valence electrons
per cell\\
$\ell$, $m$, $i$ & angular momentum, its projection, and radial index of a
pseudopotential projector\\
$E_{\mathrm{cut}}$, $\sigma$ & plane-wave cutoff; smearing width\\
\bottomrule
\end{tabular}
\end{center}
A hat ($\hat H$) marks an operator; a tilde ($\tilde p$, $\tilde V$) marks
a Fourier- or Bessel-transformed function. ``Pair'' always means a
Kramers pair of spinor bands (two degenerate bands; Section~\ref{sec:symmetry}).
A few letters are unavoidably reused, each time with the meaning standard
in that field, and the context always says which is meant: $\sigma$ is the
smearing width, but $\bm\sigma$ the Pauli matrices and, in
Chapter~\ref{ch:conductivity}, the conductivity; $s$ is a spinor index, but
also the smearing entropy $s(x)$, the reduced density gradient of PBE and a
Gaussian width; $\mu$ is the Fermi level, but also a PBE constant and the
index of a real spherical harmonic; $\alpha,\beta$ are spinor indices, but
also the fine-structure constant, the PW92 and PBE constants and, in
Chapter~\ref{ch:conductivity}, Cartesian directions; and $A$, $B$, $K$, $P$
each have several meanings (lattice matrices, PW92/PBE constants, mixing
amplitude; preconditioner, complex conjugation, Kerker operator, the
$K$ point; parity operator, projections, an eigensolver block).

\chapter{The Many-Electron Problem}
\label{ch:manybody}

\section{The full Hamiltonian}

A piece of gold is a collection of nuclei (charge $Z$, mass $M$, positions
$\RR_I$) and electrons (positions $\rr_i$, spins $\sigma_i$). Ignoring
relativity for now (it returns in force in Chapter~\ref{ch:soc}), the
non-relativistic Hamiltonian in atomic units is
\begin{equation}
\Ha =
\underbrace{-\sum_i \tfrac12 \nabla_i^2}_{\hat T_e}
\underbrace{-\sum_I \tfrac{1}{2M_I}\nabla_I^2}_{\hat T_N}
\underbrace{-\sum_{i,I} \frac{Z_I}{|\rr_i - \RR_I|}}_{\hat V_{eN}}
+\underbrace{\sum_{i<j} \frac{1}{|\rr_i - \rr_j|}}_{\hat V_{ee}}
+\underbrace{\sum_{I<J} \frac{Z_I Z_J}{|\RR_I - \RR_J|}}_{\hat V_{NN}}.
\label{eq:fullH}
\end{equation}
Every property of gold --- its color, conductivity, malleability, melting
point --- is in principle encoded in the eigenstates of
Eq.~\eqref{eq:fullH}. The entire difficulty of condensed-matter theory is
that we cannot solve it.

\section{Born--Oppenheimer: freezing the nuclei}

A gold nucleus is $\sim 3.6\times 10^5$ times heavier than an electron.
Electrons therefore adjust essentially instantaneously to any nuclear
arrangement. Formally, one writes the total wavefunction as
$\Psi(\{\rr\},\{\RR\}) \approx \chi(\{\RR\})\,\Phi_{\RR}(\{\rr\})$, where
$\Phi_{\RR}$ solves the \emph{electronic} problem at clamped nuclei:
\begin{equation}
\Big[\hat T_e + \hat V_{ee} + \hat V_{eN}(\{\RR\})\Big]\,
\Phi_{\RR}(\{\rr\}) = E_{\mathrm{el}}(\{\RR\})\,\Phi_{\RR}(\{\rr\}).
\label{eq:electronicH}
\end{equation}
The electronic energy $E_{\mathrm{el}}(\{\RR\}) + V_{NN}(\{\RR\})$ then acts
as a potential-energy surface on which the nuclei move; its minimum over the
lattice constant $a$ gives the equilibrium structure, and its curvature gives
the bulk modulus. The small parameter of the Born--Oppenheimer expansion is
$\kappa = (m_e/M)^{1/4} \approx 0.04$; the corrections to the clamped-nucleus
picture appear at higher orders in $\kappa$ and are irrelevant for the
ground-state properties we target. From now on
``the problem'' means Eq.~\eqref{eq:electronicH}: $N$ interacting electrons
in the \emph{external potential}
\begin{equation}
\vext(\rr) = -\sum_I \frac{Z_I}{|\rr - \RR_I|}
\end{equation}
created by the fixed nuclei (later: by pseudopotentials).

\section{Why the wavefunction is a dead end}

The electronic wavefunction $\Phi(\rr_1\sigma_1,\dots,\rr_N\sigma_N)$ is a
function of $3N$ continuous coordinates and $N$ spins, antisymmetric under
exchange of any two electrons (Pauli principle). Suppose you discretized
each coordinate axis with a modest 10 points. For a single primitive cell of
gold with 11 valence electrons, storing $\Phi$ would need $10^{33}$ complex
numbers times $2^{11}$ spin components --- more numbers than atoms in a
kilogram of gold. This ``exponential wall'' (Walter Kohn's phrase) is not
a matter of waiting for better GPUs: the
object itself is too large. Any tractable theory must replace the
wavefunction with a compact descriptor.

\section{The electron density}

The natural candidate is the \textbf{electron density}
\begin{equation}
n(\rr) = N \sum_{\sigma_1\cdots\sigma_N}\int
\big|\Phi(\rr\sigma_1, \rr_2\sigma_2, \dots, \rr_N\sigma_N)\big|^2\,
\dd\rr_2 \cdots \dd\rr_N ,
\qquad \int n(\rr)\,\dd\rr = N,
\label{eq:densitydef}
\end{equation}
the probability per unit volume of finding \emph{any} electron at $\rr$,
regardless of what the others are doing. It is a function of just three
variables no matter how many electrons there are: on a $64^3$ real-space
grid it is $2.6\times 10^5$ numbers --- about two megabytes.

Two contributions to the total energy at clamped nuclei are already exact
functionals of $n$ alone:
\begin{equation}
\braket{\Phi|\hat V_{eN}|\Phi} = \int \vext(\rr)\, n(\rr)\, \dd\rr,
\qquad
V_{NN} = \text{constant for fixed nuclei}.
\end{equation}
The kinetic energy and the electron--electron repulsion are not obviously
functionals of the density --- they depend on correlations between electrons,
i.e.\ on the full $\Phi$. The astonishing content of the Hohenberg--Kohn
theorems (next chapter) is that they are functionals of $n$ \emph{anyway}:
the ground-state density secretly determines everything.

\section{Spin densities and magnetization}

Once spin--orbit coupling enters, spin is no longer a spectator. The complete
one-body descriptor is the $2\times 2$ Hermitian \textbf{spin-density
matrix}
\begin{equation}
n_{\alpha\beta}(\rr) = N\!\!\sum_{\sigma_2\cdots\sigma_N}\!\int
\Phi^*(\rr\beta, \rr_2\sigma_2, \dots)\,
\Phi(\rr\alpha, \rr_2\sigma_2, \dots)\, \dd\rr_2\cdots \dd\rr_N ,
\qquad \alpha,\beta\in\{\uparrow,\downarrow\},
\label{eq:spindensmat}
\end{equation}
which can be decomposed on the identity and the Pauli matrices
$\bm\sigma = (\sigma_x,\sigma_y,\sigma_z)$:
\begin{equation}
n_{\alpha\beta}(\rr) = \tfrac12\Big[n(\rr)\,\delta_{\alpha\beta}
+ \mathbf{m}(\rr)\cdot\bm\sigma_{\alpha\beta}\Big],
\qquad
n = \mathrm{tr}\,\hat n, \quad
\mathbf{m}(\rr) = \mathrm{tr}\,[\bm\sigma\, \hat n(\rr)].
\end{equation}
Here $n(\rr)$ is the total density and $\mathbf{m}(\rr)$ the local
\textbf{magnetization density}, a vector field that may point in different
directions at different points (``noncollinear magnetism''). Bulk gold is
non-magnetic --- time-reversal symmetry will force $\mathbf{m}=0$
(Chapter~\ref{ch:soc}) --- but your code should carry the full
$n_{\alpha\beta}$ machinery, because with spinor wavefunctions it costs
little and keeps the theory honest.

\begin{keybox}{Where we stand}
Problem: $N \sim 10^{23}$ interacting electrons; wavefunction exponentially
large. Strategy: reformulate the ground-state problem exactly in terms of the
density $n(\rr)$ (3 variables) or spin-density matrix $n_{\alpha\beta}(\rr)$
(4 real fields: $n$ and $\mathbf m$), then approximate the one piece that has no known closed
form --- exchange and correlation.
\end{keybox}

\chapter{Hohenberg--Kohn and Kohn--Sham Theory}
\label{ch:hkks}

\section{Functional calculus in twenty minutes}

DFT is written in the language of \emph{functionals} --- machines that eat a
whole function and return a number. The total energy $E[n]$ is a functional
of the density. You need exactly two tools.

\paragraph{Functional derivative.} For a functional $F[n]$, the functional
derivative $\fdv{F}{n(\rr)}$ is defined by the first-order response to an
arbitrary small perturbation $\delta n(\rr)$:
\begin{equation}
F[n + \delta n] - F[n]
= \int \fdv{F}{n(\rr)}\,\delta n(\rr)\, \dd\rr + O(\delta n^2).
\label{eq:funcderivdef}
\end{equation}
It is the continuum analogue of the gradient: if $n$ were a vector with
components $n_i$ (which it literally becomes on your real-space grid!), then
$\delta F = \sum_i (\partial F/\partial n_i)\,\delta n_i$, and
$\fdv{F}{n(\rr)}$ is just $\partial F/\partial n_i$ divided by the grid-cell
volume. Three examples you will use constantly:
\begin{align}
F[n] &= \int \vext(\rr)\,n(\rr)\,\dd\rr
&&\Longrightarrow&
\fdv{F}{n(\rr)} &= \vext(\rr),
\\[2pt]
E_{\mathrm H}[n] &= \tfrac12 \iint
\frac{n(\rr)\,n(\rr')}{|\rr - \rr'|}\, \dd\rr\, \dd\rr'
&&\Longrightarrow&
\fdv{E_{\mathrm H}}{n(\rr)} &= \int \frac{n(\rr')}{|\rr-\rr'|}\,\dd\rr'
\equiv \vH(\rr),
\label{eq:hartreederiv}
\\[2pt]
F[n] &= \int f\big(n(\rr), \nabla n(\rr)\big)\, \dd\rr
&&\Longrightarrow&
\fdv{F}{n(\rr)} &= \pdv{f}{n} - \nabla\cdot\pdv{f}{(\nabla n)},
\label{eq:gradientfd}
\end{align}
the last being the Euler--Lagrange rule (integrate by parts once), needed for
gradient-corrected functionals in Chapter~\ref{ch:xc}.

\paragraph{Constrained minimization.} To minimize $E[n]$ subject to
$\int n = N$, introduce a Lagrange multiplier $\mu$ and demand
stationarity of $E[n] - \mu\big(\int n\,\dd\rr - N\big)$:
\begin{equation}
\fdv{E}{n(\rr)} = \mu \qquad \text{for all } \rr.
\label{eq:eulerlagrange}
\end{equation}
The multiplier $\mu$ turns out to be the chemical potential (Fermi level).

\section{The Hohenberg--Kohn theorems}

Consider $N$ electrons with the universal pieces $\hat T_e + \hat V_{ee}$
plus an external potential $\vext$. Different materials differ \emph{only}
through $\vext$ (and $N$). The 1964 theorems of Hohenberg and Kohn establish
that the ground-state density is a legitimate fundamental variable.

\begin{theorem}[HK1: the density determines the potential]
The ground-state density $n_0(\rr)$ of an interacting electron system
determines the external potential $\vext(\rr)$ uniquely, up to an additive
constant. Consequently $n_0$ determines the Hamiltonian, hence \emph{all}
ground- and excited-state properties of the system.
\end{theorem}

\begin{proof}
By contradiction, in two steps. Assume for simplicity non-degenerate ground
states. Suppose two potentials $\vext^{(1)} \neq \vext^{(2)} + \mathrm{const}$
give ground states $\Phi_1 \neq \Phi_2$ with the \emph{same} density $n_0$.
(That $\Phi_1\neq\Phi_2$ follows because they satisfy Schr\"odinger equations
with potentials differing by more than a constant.) By the variational
principle, using $\Phi_2$ as a trial state for Hamiltonian $\Ha_1$:
\begin{equation}
E_1 < \braket{\Phi_2|\Ha_1|\Phi_2}
= \braket{\Phi_2|\Ha_2|\Phi_2}
+ \braket{\Phi_2|\hat V_{\!1} - \hat V_{\!2}|\Phi_2}
= E_2 + \int \big[\vext^{(1)} - \vext^{(2)}\big]\, n_0(\rr)\, \dd\rr .
\end{equation}
Swapping labels $1\leftrightarrow 2$:
\begin{equation}
E_2 < E_1 + \int \big[\vext^{(2)} - \vext^{(1)}\big]\, n_0(\rr)\, \dd\rr .
\end{equation}
Adding the two inequalities gives $E_1 + E_2 < E_1 + E_2$, a contradiction.
Hence two essentially different potentials cannot share a ground-state
density: the map $n_0 \mapsto \vext$ is single-valued.
\end{proof}

\begin{theorem}[HK2: variational principle for the density]
There exists a universal functional $F_{\mathrm{HK}}[n]$ (the same for every
material) such that
\begin{equation}
E_v[n] \equiv F_{\mathrm{HK}}[n] + \int \vext(\rr)\, n(\rr)\, \dd\rr
\end{equation}
satisfies $E_v[n] \geq E_0$ for every admissible density with $\int n = N$,
with equality if and only if $n = n_0$ (for a non-degenerate ground state;
with degeneracy, equality holds for any ground-state density).
\end{theorem}

\begin{proof}
The cleanest construction is Levy's constrained search. Define
\begin{equation}
F_{\mathrm{HK}}[n] \equiv
\min_{\Phi \to n} \braket{\Phi | \hat T_e + \hat V_{ee} | \Phi},
\label{eq:levysearch}
\end{equation}
the minimum over all antisymmetric $N$-electron wavefunctions that yield the
prescribed density $n$. This is manifestly universal: $\vext$ does not appear.
Then for any $n$,
\begin{equation}
E_v[n]
= \min_{\Phi\to n}\braket{\Phi|\hat T_e + \hat V_{ee}|\Phi}
+ \int \vext\, n
= \min_{\Phi\to n}\braket{\Phi|\Ha_{\mathrm{el}}|\Phi}
\;\geq\; \min_{\Phi}\braket{\Phi|\Ha_{\mathrm{el}}|\Phi} = E_0 ,
\end{equation}
because searching within a fixed density is more restrictive than searching
over all wavefunctions. Minimizing over $n$ removes the restriction, so the
minimum of $E_v[n]$ is $E_0$, attained at $n_0$.
\end{proof}

\begin{remark}
These theorems are an existence statement, not a recipe: nobody knows
$F_{\mathrm{HK}}[n]$ in closed form --- knowing it exactly would be
equivalent to solving the many-body problem. The genius of Kohn and Sham
(1965) was to isolate the \emph{large, known} parts of $F_{\mathrm{HK}}$ and
approximate only the small remainder.
\end{remark}

\section{The Kohn--Sham construction}
\label{sec:ks}

The dominant contribution to the kinetic energy is well approximated by that
of \emph{non-interacting} electrons with the same density. So: introduce an
auxiliary fictitious system of $N$ non-interacting electrons moving in some
effective potential $\veff(\rr)$, engineered so that its ground-state
density equals the true interacting density. Non-interacting electrons are
easy: they occupy single-particle orbitals $\psi_i(\rr)$ (for now spinless;
spinors come in Chapter~\ref{ch:soc}) obeying a one-body Schr\"odinger
equation, and
\begin{equation}
n(\rr) = \sum_{i} f_i\, |\psi_i(\rr)|^2 ,
\qquad
T_s[n] = \sum_i f_i \int \psi_i^*(\rr)
\Big(-\tfrac{1}{2}\nabla^2\Big) \psi_i(\rr)\,\dd\rr,
\label{eq:ksdensity}
\end{equation}
with occupations $f_i \in [0,1]$ and $\sum_i f_i = N$. (Here $i$ counts
\emph{spin-orbitals}: one state, one electron at most. If spin is a pure
spectator one can instead let each spatial orbital hold two electrons; we
will not, because with spin--orbit coupling every state is its own spinor.)
The many-electron wavefunction of non-interacting fermions is a single
\emph{Slater determinant} of the occupied orbitals,
$\Phi = (N!)^{-1/2}\det[\psi_i(\rr_j)]$ --- the antisymmetrized product
``one electron in each orbital''; Eq.~\eqref{eq:ksdensity} is
what Eq.~\eqref{eq:densitydef} gives for a determinant.
$T_s$ is the \emph{non-interacting} kinetic energy --- not the true one, but
its largest part, and crucially one we can evaluate exactly through the
orbitals.

Now \emph{define} the exchange--correlation functional as everything left
over:
\begin{equation}
\boxed{\;
F_{\mathrm{HK}}[n] = T_s[n] + E_{\mathrm H}[n] + \Exc[n]
\;}
\qquad\Longleftrightarrow\qquad
\Exc[n] \equiv \big(T - T_s\big) + \big(V_{ee} - E_{\mathrm H}\big).
\label{eq:excdef}
\end{equation}
$\Exc$ contains (i) the difference between true and non-interacting kinetic
energies, (ii) the exchange energy --- the energy lowering from Pauli
antisymmetry keeping same-spin electrons apart, (iii) the correlation energy
--- the further lowering from electrons dynamically dodging one another
beyond mean field, and (iv) the correction for the self-interaction hiding
inside $E_{\mathrm H}$ (in $E_{\mathrm H}$ each electron feels its own
charge). $\Exc$ is a small fraction of the total energy ($\sim$1--10\%), yet
it contains all the chemistry; approximating it is Chapter~\ref{ch:xc}.

\section{Deriving the Kohn--Sham equations}

Apply the Euler--Lagrange condition \eqref{eq:eulerlagrange} to
$E_v[n] = T_s[n] + E_{\mathrm H}[n] + \Exc[n] + \int \vext n$:
\begin{equation}
\fdv{T_s}{n(\rr)} + \vH(\rr) + \vxc(\rr) + \vext(\rr) = \mu ,
\qquad
\vxc(\rr) \equiv \fdv{\Exc}{n(\rr)} .
\label{eq:eltrue}
\end{equation}
Now write the same stationarity condition for the fictitious non-interacting
system in potential $\veff$:
\begin{equation}
\fdv{T_s}{n(\rr)} + \veff(\rr) = \mu' .
\label{eq:elks}
\end{equation}
Equations \eqref{eq:eltrue} and \eqref{eq:elks} produce the same density if
and only if
\begin{equation}
\boxed{\;
\veff(\rr) = \vext(\rr) + \vH[n](\rr) + \vxc[n](\rr)
\;}
\label{eq:veff}
\end{equation}
(up to a constant). But the ground state of non-interacting electrons in
$\veff$ is found by solving the one-body eigenproblem. Minimizing
$E$ with respect to the orbitals themselves, subject to orthonormality
$\braket{\psi_i|\psi_j} = \delta_{ij}$ (Lagrange multipliers
$\varepsilon_{ij}$, diagonalized to $\varepsilon_i$), gives the same result
directly. A complex function has two independent real degrees of freedom,
its real and imaginary parts; it is equivalent, and shorter, to treat
$\psi_i$ and $\psi_i^*$ as the two independent variables and set the
derivative with respect to $\psi_i^*$ to zero (the derivative with respect
to $\psi_i$ gives the complex-conjugate equation). Varying $\psi_i^*$,
\begin{equation}
\fdv{}{\psi_i^*(\rr)}\Big[E - \textstyle\sum_{jk}
\varepsilon_{jk}\big(\braket{\psi_j|\psi_k} - \delta_{jk}\big)\Big] = 0
\;\Longrightarrow\;
\Big[-\tfrac12 \nabla^2 + \veff(\rr)\Big] \psi_i(\rr)
= \varepsilon_i\, \psi_i(\rr).
\label{eq:kseq}
\end{equation}
(Strictly, varying $\psi_i^*$ gives
$f_i\,\hat h\,\psi_i = \sum_k \varepsilon_{ik}\psi_k$ with
$\hat h = -\tfrac12\nabla^2+\veff$, because $\psi_i$ enters $E$ with weight
$f_i$. Project onto $\psi_j$: $f_i\braket{\psi_j|\hat h|\psi_i} =
\varepsilon_{ij}$. The multiplier matrix can be taken Hermitian, and so is
$\hat h$; combining the $(i,j)$ and $(j,i)$ equations gives
$(f_i - f_j)\braket{\psi_j|\hat h|\psi_i} = 0$. So $\hat h$ has no matrix
elements between orbitals of different occupation, and within a set of
equally occupied orbitals a unitary rotation --- which changes neither $n$
nor $E$ --- diagonalizes it. Hence the form above.)
Equations \eqref{eq:ksdensity}, \eqref{eq:veff} and \eqref{eq:kseq} are the
\textbf{Kohn--Sham equations}. Note the nonlinearity: $\veff$ depends on $n$,
which depends on the $\psi_i$, which depend on $\veff$. The equations must be
solved \emph{self-consistently} --- the fixed-point iteration at the heart of
your program (Chapter~\ref{ch:algorithms}).

\section{The total energy}

Given a self-consistent solution, the ground-state energy is
\begin{equation}
E = T_s + \int \vext\, n\, \dd\rr + E_{\mathrm H}[n] + \Exc[n] + E_{NN}.
\label{eq:etot1}
\end{equation}
A form better suited to computation eliminates $T_s$ using the eigenvalues.
Multiply \eqref{eq:kseq} by $f_i\psi_i^*$, integrate, and sum:
$\sum_i f_i \varepsilon_i = T_s + \int \veff\, n$. Hence
\begin{equation}
\boxed{\;
E = \sum_i f_i\,\varepsilon_i
\;-\; E_{\mathrm H}[n]
\;-\; \int \vxc(\rr)\, n(\rr)\, \dd\rr
\;+\; \Exc[n]
\;+\; E_{NN}.
\;}
\label{eq:etot2}
\end{equation}
The corrections after the band-structure sum $\sum f_i\varepsilon_i$ are
called \emph{double-counting} terms: the eigenvalue sum counts the Hartree
energy twice (each electron feels the full potential of all electrons) and
counts $\int \vxc n$ instead of $\Exc$. For a metal with smeared occupations
an entropy term is added (Section~\ref{sec:smearing}). Both forms are
worth coding: Eq.~\eqref{eq:etot1} (``direct'') and Eq.~\eqref{eq:etot2}
(``eigenvalue sum'') must agree to rounding error at self-consistency, and
a disagreement pinpoints a bookkeeping bug.

\begin{keybox}{Interpretation warnings}
The Kohn--Sham orbitals and eigenvalues belong to the \emph{fictitious}
system. Their only rigorous roles are to build $n(\rr)$ and $T_s$.
Nevertheless, the KS band structure of metals like gold is a good
approximation to the measured quasiparticle bands (photoemission; for gold
the $d$ bands come out several tenths of an eV too shallow --- $0.3$--$0.9$\,eV
with the functionals and pseudopotential of this document,
Chapters~\ref{ch:soc} and~\ref{ch:colour}), and
(for the exact functional) the highest occupied eigenvalue is exact: it
equals minus the ionization energy of a finite system, and the Fermi level
of a metal. Band
\emph{gaps} of semiconductors are systematically underestimated --- a famous
limitation that does not affect our metallic target.
\end{keybox}

\section{Spin-density functional theory}

When a magnetic field couples to spin, or when SOC entangles spin and orbit,
the theory is generalized (von Barth--Hedin 1972): the energy becomes a
functional of the full spin-density matrix
$n_{\alpha\beta}(\rr)$ of Eq.~\eqref{eq:spindensmat},
$E = E[n_{\alpha\beta}]$, and the XC potential becomes a $2\times2$ matrix
field,
\begin{equation}
v^{\mathrm{xc}}_{\alpha\beta}(\rr)
= \fdv{\Exc}{n_{\beta\alpha}(\rr)}
= \vxc(\rr)\,\delta_{\alpha\beta}
+ \mathbf{B}_{\mathrm{xc}}(\rr)\cdot\bm\sigma_{\alpha\beta},
\label{eq:vxcmatrix}
\end{equation}
with an \emph{exchange--correlation magnetic field}
$\mathbf{B}_{\mathrm{xc}} = \delta \Exc / \delta \mathbf{m}$. The Kohn--Sham
orbitals become two-component spinors
$\psi_i = (\psi_i^{\uparrow}, \psi_i^{\downarrow})^{\mathsf T}$ and the KS
equation a $2\times 2$ matrix equation --- exactly the structure SOC requires,
developed fully in Chapter~\ref{ch:soc}. For non-magnetic gold,
$\mathbf{m} = 0$ and $\mathbf{B}_{\mathrm{xc}} = 0$ at self-consistency, but
your data structures should support the general case.

\chapter{Exchange and Correlation}
\label{ch:xc}

This chapter builds the two functionals you will implement --- LDA (in the
PW92 parametrization) and the PBE GGA --- from the ground up, then extends
them to spinor/noncollinear densities as required by SOC. This is the only
approximated physics in the whole code; everything else is numerics.

\section{What exchange and correlation physically are}

Take one electron sitting at $\rr$. The conditional density of the other
$N-1$ electrons is depleted around it: same-spin electrons avoid it because
of Pauli antisymmetry (\emph{exchange}), and all electrons avoid it because
of Coulomb repulsion (\emph{correlation}). The missing charge --- exactly one
electron's worth --- is called the \textbf{exchange--correlation hole}
$n_{\mathrm{xc}}(\rr,\rr')$, satisfying the sum rule
$\int n_{\mathrm{xc}}(\rr,\rr')\,\dd\rr' = -1$. One can show exactly that
\begin{equation}
\Exc[n] = \tfrac12 \iint
\frac{n(\rr)\,\bar n_{\mathrm{xc}}(\rr,\rr')}{|\rr-\rr'|}\,
\dd\rr\,\dd\rr' ,
\end{equation}
the electrostatic interaction of each electron with its own
(coupling-constant-averaged) hole. Two consequences guide all approximations:
$\Exc$ is negative and mostly \emph{local} --- the hole is centered on the
electron and its long-range details matter less than its depth and sum rule.
That is why approximating $\Exc$ by local information succeeds far beyond
what one might expect.

\section{The homogeneous electron gas}
\label{sec:heg}

The reference system for all local approximations: $N$ electrons in volume
$V$ with a uniform neutralizing positive background, $n = N/V$ constant.
Conventionally parametrized by the \textbf{Wigner--Seitz radius}
\begin{equation}
\frac{4}{3}\pi r_s^3 = \frac{1}{n}
\qquad\Longleftrightarrow\qquad
r_s = \left(\frac{3}{4\pi n}\right)^{1/3},
\end{equation}
the radius (in Bohr) of the sphere containing one electron on average.
Valence densities of metals correspond to $r_s \approx 2$--$6$; gold's $6s$
electrons have $r_s \approx 3$.

\paragraph{Kinetic energy.} Non-interacting electrons in a box fill plane
waves $\ee^{\ii\kk\cdot\rr}/\sqrt{V}$ up to the Fermi wavevector $k_F$.
Counting states (2 spins, density of $\kk$-states $V/(2\pi)^3$):
\begin{equation}
N = 2\,\frac{V}{(2\pi)^3}\,\frac{4}{3}\pi k_F^3
\;\Longrightarrow\;
k_F = (3\pi^2 n)^{1/3},
\qquad
\frac{T_s}{N} = \frac{3}{5}\,\frac{k_F^2}{2}
= \frac{3}{10}(3\pi^2 n)^{2/3}.
\label{eq:heg_ts}
\end{equation}

\paragraph{Exchange energy --- full derivation.}
Evaluate $\braket{\hat V_{ee}}$ in a Slater determinant of orbitals
$\phi_i$. Antisymmetry produces two terms: the classical Hartree energy,
and a negative term in which the two electrons swap orbitals --- the
\emph{exchange energy}. Swapping is only possible between orbitals of the
same spin, hence the sum over $\sigma$ (this is the Hartree--Fock
expression):
\begin{equation}
E_x = -\frac{1}{2} \sum_{\sigma} \sum_{\kk,\kk'}^{\mathrm{occ}}
\iint \frac{\phi_{\kk}^*(\rr)\phi_{\kk'}(\rr)\,
\phi_{\kk'}^*(\rr')\phi_{\kk}(\rr')}{|\rr-\rr'|}\,\dd\rr\,\dd\rr'.
\end{equation}
With plane waves, the double integral gives the Fourier transform of the
Coulomb interaction, $4\pi/|\kk - \kk'|^2$, times $1/V$:
\begin{equation}
E_x = -\frac{1}{2}\cdot 2 \cdot \frac{1}{V}
\sum_{|\kk|<k_F}\sum_{|\kk'|<k_F} \frac{4\pi}{|\kk-\kk'|^2}
= -\frac{V}{(2\pi)^6}
\iint_{|\kk|,|\kk'|<k_F} \frac{4\pi}{|\kk-\kk'|^2}\,\dd\kk\,\dd\kk'.
\end{equation}
The inner integral is elementary (it is the electrostatic potential of a
uniformly filled Fermi sphere in $\kk$-space):
\begin{equation}
\int_{|\kk'|<k_F} \frac{\dd\kk'}{|\kk - \kk'|^2}
= 2\pi k_F\, S\!\left(\frac{k}{k_F}\right),
\qquad
S(x) = 1 + \frac{1-x^2}{2x}\ln\left|\frac{1+x}{1-x}\right| .
\end{equation}
Performing the remaining radial integral
$\int_0^1 x^2 S(x)\,\dd x = \tfrac12$ yields the classic result
\begin{equation}
\frac{E_x}{N}
= -\frac{3 k_F}{4\pi}
= -\frac{3}{4}\left(\frac{3}{\pi}\right)^{1/3} n^{1/3}
\qquad\Longrightarrow\qquad
\varepsilon_x(n) = -\frac{3}{4}\left(\frac{3}{\pi}\right)^{1/3} n^{1/3}
\approx -\frac{0.4582}{r_s}\ \mathrm{Ha}.
\label{eq:ldax}
\end{equation}
Exchange per electron grows as $n^{1/3}$: denser gas, deeper exchange hole.

\paragraph{Correlation energy.} No closed form exists. Its limits are known
analytically --- the high-density (RPA) expansion
$\varepsilon_c \to A \ln r_s + B + \dots$ and the low-density (Wigner
crystal) expansion $\varepsilon_c \to -d_0/r_s + d_1/r_s^{3/2} + \dots$ ---
and quantum Monte Carlo (Ceperley--Alder 1980) supplies essentially exact
values in between. \textbf{Perdew--Wang 1992 (PW92)} interpolates all of this
with the analytic form (for spin-unpolarized and fully polarized cases,
$\zeta = 0, 1$):
\begin{equation}
\varepsilon_c(r_s)
= -2 A (1 + \alpha_1 r_s)\,
\ln\!\left[1 + \frac{1}{2A\big(\beta_1 r_s^{1/2} + \beta_2 r_s
+ \beta_3 r_s^{3/2} + \beta_4 r_s^{2}\big)}\right],
\label{eq:pw92}
\end{equation}
with the constants of Table~\ref{tab:pw92}. Nonmagnetic gold only ever
uses the first row. This smooth, infinitely differentiable expression is ideal
for automatic differentiation: \emph{in JAX you will get $\vxc =
\dd(n\varepsilon_{\mathrm{xc}})/\dd n$ for free via \texttt{grad}}, but you
should still know the analytic derivative to verify it: since
$n \propto r_s^{-3}$, for any function $\varepsilon(r_s)$
\begin{equation}
\frac{\dd (n\varepsilon)}{\dd n}
= \varepsilon - \frac{r_s}{3}\,\frac{\dd \varepsilon}{\dd r_s},
\label{eq:rsderiv}
\end{equation}
which gives $v_x = \tfrac43\varepsilon_x$ at once and $v_c$ after one
differentiation of Eq.~\eqref{eq:pw92}. Writing $Q(r_s) = 2A\big(\beta_1
r_s^{1/2} + \beta_2 r_s + \beta_3 r_s^{3/2} + \beta_4 r_s^2\big)$ for the
denominator inside the logarithm,
\begin{equation}
\begin{gathered}
\frac{\dd\varepsilon_c}{\dd r_s}
= -2A\alpha_1 \ln\!\Big(1 + \frac{1}{Q}\Big)
+ \frac{2A(1+\alpha_1 r_s)}{Q^2 + Q}\,\frac{\dd Q}{\dd r_s},
\\
\frac{\dd Q}{\dd r_s} = 2A\Big(\frac{\beta_1}{2 r_s^{1/2}} + \beta_2
+ \tfrac32\beta_3 r_s^{1/2} + 2\beta_4 r_s\Big).
\end{gathered}
\label{eq:pw92deriv}
\end{equation}

\begin{table}[htbp]
\centering
\begin{tabular}{@{}lcccccc@{}}
\toprule
 & $A$ & $\alpha_1$ & $\beta_1$ & $\beta_2$ & $\beta_3$ & $\beta_4$\\
\midrule
$\varepsilon_c(r_s, \zeta{=}0)$ & 0.031091 & 0.21370 & 7.5957 & 3.5876 & 1.6382 & 0.49294\\
$\varepsilon_c(r_s, \zeta{=}1)$ & 0.015545 & 0.20548 & 14.1189 & 6.1977 & 3.3662 & 0.62517\\
$-\alpha_c(r_s)$ & 0.016887 & 0.11125 & 10.357 & 3.6231 & 0.88026 & 0.49671\\
\bottomrule
\end{tabular}
\caption{PW92 constants for the form \eqref{eq:pw92}. The third row gives
\emph{minus} the spin stiffness $\alpha_c$. These are the rounded constants
of the original paper; some libraries use more digits for $A$
($0.0310907$, $0.01554535$, $0.0168869$; the ``modified PW92'' used
inside most PBE implementations), which changes $\varepsilon_c$ in the
seventh decimal place. Either is acceptable --- but know which one you
have when comparing total energies between codes to better than
$10^{-6}$\,Ha.}
\label{tab:pw92}
\end{table}

\section{The Local Density Approximation}

The LDA declares: at each point of an inhomogeneous system, use the XC energy
density of a homogeneous gas with the local density,
\begin{equation}
\boxed{\;
\Exc^{\mathrm{LDA}}[n] = \int n(\rr)\,
\varepsilon_{\mathrm{xc}}\big(n(\rr)\big)\, \dd\rr,
\qquad
\varepsilon_{\mathrm{xc}} = \varepsilon_x + \varepsilon_c .
\;}
\end{equation}
The XC potential follows from Eq.~\eqref{eq:funcderivdef}:
\begin{equation}
\vxc^{\mathrm{LDA}}(\rr)
= \varepsilon_{\mathrm{xc}}\big(n(\rr)\big)
+ n(\rr)\, \frac{\dd \varepsilon_{\mathrm{xc}}}{\dd n}\Big|_{n(\rr)} .
\label{eq:vxclda}
\end{equation}
For the exchange part explicitly:
$v_x = -\big(3n/\pi\big)^{1/3}$. On the computer, LDA is a pure pointwise
(embarrassingly parallel) map over the real-space grid --- one fused GPU
kernel.

\paragraph{Spin polarization (LSDA).} With collinear spin densities
$n_\uparrow, n_\downarrow$, define $n = n_\uparrow + n_\downarrow$ and the
polarization $\zeta = (n_\uparrow - n_\downarrow)/n$. Exchange obeys the
exact \emph{spin-scaling relation}
$E_x[n_\uparrow,n_\downarrow] = \tfrac12 E_x[2n_\uparrow]
+ \tfrac12 E_x[2n_\downarrow]$, giving
\begin{equation}
\varepsilon_x(n,\zeta) = \varepsilon_x(n)\,
\frac{(1+\zeta)^{4/3} + (1-\zeta)^{4/3}}{2}.
\end{equation}
Correlation uses the PW92 interpolation between $\zeta=0$ and $\zeta=1$,
\begin{gather}
\varepsilon_c(r_s,\zeta) = \varepsilon_c(r_s,0)
+ \alpha_c(r_s)\,\frac{f(\zeta)}{f''(0)}\,(1-\zeta^4)
+ \big[\varepsilon_c(r_s,1) - \varepsilon_c(r_s,0)\big] f(\zeta)\,\zeta^4,
\\
f(\zeta) = \frac{(1+\zeta)^{4/3} + (1-\zeta)^{4/3} - 2}{2^{4/3} - 2},
\notag
\end{gather}
with $f''(0) = 1.709921$ and the three functions of
Table~\ref{tab:pw92}; the two potentials are
$v_{\mathrm{xc}}^{\sigma} = \partial (n\varepsilon_{\mathrm{xc}})/\partial
n_\sigma$.

\section{The PBE Generalized Gradient Approximation}
\label{sec:pbe}

LDA overbinds: for gold, all-electron-quality LDA calculations give a
lattice constant $0.6\%$ too small ($4.053$ against a measured
$4.078$\,\AA) and a bulk modulus $13\%$ too large ($195$ against
$173$\,GPa). (Both comparisons are with room-temperature experiment.
Against the static-lattice value of Chapter~\ref{ch:intro},
$a \approx 4.063$\,\AA, the LDA error shrinks to about $0.25\%$ but keeps its
sign. For the bulk modulus the verdict depends on how the zero-point motion
is removed: against the $0$\,K elastic constants, $\approx 180$\,GPa, LDA
is still $8\%$ too stiff; against the $\approx 196$\,GPa used in
density-functional benchmarks it is about right.) The next rung feeds the local
\emph{gradient} to the functional. The natural dimensionless inhomogeneity
measure is the \textbf{reduced gradient}
\begin{equation}
s(\rr) = \frac{|\nabla n|}{2 k_F n}
= \frac{|\nabla n|}{2 (3\pi^2)^{1/3} n^{4/3}},
\end{equation}
the fractional change of the density over a Fermi wavelength (up to a
factor of order one). PBE
(Perdew--Burke--Ernzerhof 1996) writes exchange as LDA times an
\emph{enhancement factor}:
\begin{equation}
E_x^{\mathrm{PBE}} = \int n\, \varepsilon_x(n)\, F_x(s)\, \dd\rr,
\qquad
F_x(s) = 1 + \kappa - \frac{\kappa}{1 + \mu s^2/\kappa},
\label{eq:pbex}
\end{equation}
with $\kappa = 0.804$ and $\mu = \pi^2\beta/3 = 0.2195149727645171$.
Every constant is fixed by an exact condition, not by fitting: $\mu$ is
chosen so that, for slowly varying densities, the gradient correction to
exchange cancels that to correlation (whose coefficient is the $\beta$
given below), which preserves the LDA's excellent linear response of the
homogeneous gas; and $\kappa$ saturates $F_x \le 1+\kappa$ to satisfy
the Lieb--Oxford lower bound on $\Exc$. PBE correlation is
\begin{equation}
E_c^{\mathrm{PBE}} = \int n \Big[\varepsilon_c(r_s,\zeta)
+ H(r_s, \zeta, t)\Big] \dd\rr,
\qquad
H = \gamma \phi^3
\ln\!\left[1 + \frac{\beta}{\gamma} t^2
\frac{1 + A t^2}{1 + A t^2 + A^2 t^4}\right],
\label{eq:pbec}
\end{equation}
where $t = |\nabla n| / (2\phi k_s n)$ is the gradient measured against the
Thomas--Fermi screening wavevector $k_s = \sqrt{4 k_F/\pi}$,
$\phi(\zeta) = [(1+\zeta)^{2/3} + (1-\zeta)^{2/3}]/2$,
$\gamma = (1 - \ln 2)/\pi^2 = 0.0310906908696549$,
$\beta = 0.06672455060314922$, and
$A = (\beta/\gamma)\,[\exp(-\varepsilon_c/(\gamma\phi^3)) - 1]^{-1}$
(not to be confused with the PW92 constant $A$; here $\varepsilon_c =
\varepsilon_c(r_s,\zeta)$ is the LDA correlation energy per electron at
that grid point). $H$ is built so that correlation vanishes in the
rapid-variation limit and recovers the exact gradient expansion in the
slow-variation limit. For nonmagnetic gold $\zeta = 0$ and $\phi = 1$
throughout. (The PBE paper prints the rounded values $\beta = 0.066725$ and
$\mu = 0.21951$; the libxc library and most production codes use the
digits above, and the difference is visible --- $5\times10^{-6}$\,Ha in the total energy of gold.
Type them in full.)

For spin-polarized densities PBE exchange needs no new formula: apply the
same spin-scaling relation as in the LSDA,
$E_x^{\mathrm{PBE}}[n_\uparrow, n_\downarrow] = \tfrac12
E_x^{\mathrm{PBE}}[2n_\uparrow] + \tfrac12 E_x^{\mathrm{PBE}}[2n_\downarrow]$,
with Eq.~\eqref{eq:pbex} evaluated for each doubled spin density and its
gradient.

\paragraph{The GGA potential.} Because the integrand depends on $\nabla n$,
Eq.~\eqref{eq:gradientfd} gives
\begin{equation}
\vxc^{\mathrm{GGA}}(\rr)
= \pdv{f}{n}
- \nabla \cdot \pdv{f}{(\nabla n)}
= \pdv{f}{n}
- \nabla \cdot \left[\pdv{f}{|\nabla n|}\,
\frac{\nabla n}{|\nabla n|}\right],
\qquad f \equiv n\,\varepsilon_{\mathrm{xc}}^{\mathrm{GGA}} .
\label{eq:vxcgga}
\end{equation}
In practice write $f$ as a function of $n$ and of the squared gradient
$\varsigma \equiv |\nabla n|^2$ (no square roots, no division by
$|\nabla n|$). Since $\partial\varsigma/\partial(\nabla n) = 2\nabla n$, the
chain rule turns the same formula into
\begin{equation}
\vxc^{\mathrm{GGA}} = \pdv{f}{n}
- \nabla\cdot\Big[2\,\pdv{f}{\varsigma}\,\nabla n\Big].
\label{eq:vxcgga2}
\end{equation}
In terms of $\varsigma$ the two gradient measures of PBE are simply
$s^2 = \varsigma/(2k_Fn)^2$ and $t^2 = \varsigma/(2\phi k_s n)^2$, and only
$s^2$ and $t^2$ enter Eqs.~\eqref{eq:pbex}--\eqref{eq:pbec}.
The recipe is therefore: (1) $\nabla n$ by FFT; (2) pointwise $f$,
$\partial f/\partial n$, $\partial f/\partial\varsigma$; (3) form the vector
field $2(\partial f/\partial\varsigma)\nabla n$; (4) its divergence by FFT;
(5) subtract. The energy is $\Exc = \int f\,\dd\rr$.
\begin{gpubox}{GPU/JAX notes for XC}
(i) Evaluate $\nabla n$ \emph{spectrally}: FFT the density,
multiply by $\ii \GG$, inverse-FFT --- exact on the plane-wave grid, three
batched FFTs, no finite-difference stencils. The divergence in
\eqref{eq:vxcgga} is likewise done as $\mathrm{IFFT}[\,\ii\GG\cdot
\mathrm{FFT}[\cdot]\,]$ (the ``White--Bird'' route). Keep the real part of
each inverse transform: on a grid with an even number of points the highest
(Nyquist) frequency has no partner of opposite sign, so $\ii\GG$ times a
real function's transform is not exactly the transform of a real function;
a complex $\vxc$ would make $\hat H$ non-Hermitian. (ii) The pointwise parts
of LDA/PBE are pure elementwise math, perfectly fused by XLA. (iii) You can
implement \emph{only the energy density} $f(n, |\nabla n|^2)$ and obtain
$\partial f/\partial n$ and $\partial f/\partial |\nabla n|^2$ by automatic
differentiation --- one of the places where JAX genuinely shortens a DFT
code. (Without automatic differentiation, central differences with a
relative step of $10^{-5}$ --- one-sided where $|\nabla n|^2 = 0$ ---
reproduce the anchors at the end of this chapter to all quoted digits.) Guard against $n \to 0$ (clamp the density from below at
$\sim10^{-12}$ before evaluating the functional) because negative powers
of $n$ appear in the derivatives; with $f$ written in terms of
$|\nabla n|^2$ there is no division by the gradient and no further guard
is needed.
\end{gpubox}

\paragraph{Which functional for gold?} Implement both. LDA overbinds; PBE
underbinds (lattice constant
$2.3\%$ too large for gold in all-electron-quality calculations) but
improves energetics. (With the 11-electron pseudopotential used here both
lattice constants come out larger still, by $0.03$\,\AA{} (PBE) and
$0.04$\,\AA{} (LDA); the
computed values are in Chapter~\ref{ch:verify}, Stage~3.) Both place the $d$ bands several tenths of an eV too
close to the Fermi level compared with photoemission (Table~\ref{tab:dbandexp}). Comparing them
against experiment is itself a classic exercise: it teaches you that the XC
functional is a controlled physical approximation, not a fudge.

\section{Exchange--correlation with spinors: noncollinear DFT}
\label{sec:noncollinear_xc}

With SOC, orbitals are two-component spinors and the fundamental variable is
the $2\times2$ matrix $n_{\alpha\beta}(\rr)$, i.e.\ $\{n(\rr),
\mathbf{m}(\rr)\}$. But LSDA/PBE only know two scalar inputs
$(n_\uparrow, n_\downarrow)$. The standard resolution (K\"ubler et al.) is
the \textbf{locally collinear approximation}: at each grid point,
diagonalize the spin-density matrix --- i.e.\ quantize spin along the local
magnetization direction $\hat{\mathbf m}(\rr)$ --- and feed the eigenvalues
\begin{equation}
n_{\pm}(\rr) = \tfrac12\big[n(\rr) \pm |\mathbf{m}(\rr)|\big]
\end{equation}
to the ordinary collinear functional. The resulting potential matrix is
\begin{equation}
v^{\mathrm{xc}}_{\alpha\beta}(\rr)
= \vxc\big(n, |\mathbf m|\big)\,\delta_{\alpha\beta}
+ B_{\mathrm{xc}}\big(n, |\mathbf m|\big)\,
\hat{\mathbf m}(\rr)\cdot \bm\sigma_{\alpha\beta},
\qquad
\begin{aligned}
\vxc &= \tfrac12\big(v^{+}_{\mathrm{xc}} + v^{-}_{\mathrm{xc}}\big),\\
B_{\mathrm{xc}} &= \tfrac12\big(v^{+}_{\mathrm{xc}} - v^{-}_{\mathrm{xc}}\big),
\end{aligned}
\label{eq:noncoll_vxc}
\end{equation}
with $v^{\pm}_{\mathrm{xc}} = \partial (n\varepsilon_{\mathrm{xc}})/\partial
n_\pm$: the XC magnetic field is parallel to the local magnetization. For
nonmagnetic gold $\mathbf m = 0$ everywhere and only the scalar $\vxc(n)$
survives --- but implement Eq.~\eqref{eq:noncoll_vxc} in general form; it
costs one pointwise $2\times2$ decomposition and lets the same code later
handle magnetic materials (e.g.\ FePt, where SOC + magnetism give magnetic
anisotropy). Two practical notes: where $|\mathbf m|$ is below a small
threshold the direction $\hat{\mathbf m}$ is undefined --- set
$B_{\mathrm{xc}}\hat{\mathbf m} = 0$ there; and for a GGA the gradients
fed to the functional are those of $n_\pm$, i.e.\ of $n$ and
$|\mathbf m|$. (The gradient of $|\mathbf m|$ is ill-defined where
$\mathbf m$ passes through zero; this is harmless for gold, where
$\mathbf m \equiv 0$, but noncollinear GGA for magnets needs more care
than this simple prescription.)

\begin{keybox}{Sanity anchors for your implementation}
For the unpolarized HEG at $r_s = 3$ (close to gold's $6s$ density):
$\varepsilon_x = -0.15272$\,Ha, $\varepsilon_c^{\mathrm{PW92}} =
-0.03694$\,Ha per electron, and $\vxc = -0.24668$\,Ha. At $r_s=2$:
$\varepsilon_x = -0.22908$, $\varepsilon_c = -0.04476$. Fully polarized
correlation at $r_s = 3$: $\varepsilon_c(\zeta{=}1) = -0.01991$. Reproduce
these to $10^{-5}$ before going further. Also check $F_x(s\!=\!0)=1$ and
$F_x(s\to\infty)=1.804$.

One full PBE point, unpolarized, at $n = 0.05$\,Bohr$^{-3}$ and
$\varsigma = |\nabla n|^2 = 0.01$\,Bohr$^{-8}$ (checked against the libxc library):
$\varepsilon_{\mathrm{xc}} = f/n = -0.332094$\,Ha,
$\partial f/\partial n = -0.392041$,
$\partial f/\partial\varsigma = -0.083683$; the LDA values at the same
density are $\varepsilon_{\mathrm{xc}} = -0.320387$ and
$\vxc = -0.418065$. Agreement to six digits validates every constant in
Eqs.~\eqref{eq:pbex}--\eqref{eq:pbec} at once.
\end{keybox}

\chapter{Electrons in Crystals: Plane Waves, k-Points, and Metallic Occupations}
\label{ch:solids}

\section{Direct and reciprocal lattices}

A crystal is defined by three primitive vectors $\mathbf a_1, \mathbf a_2,
\mathbf a_3$; the set of translations $\RR = n_1\mathbf a_1 + n_2\mathbf a_2
+ n_3 \mathbf a_3$ ($n_i \in \mathbb Z$) is the Bravais lattice. For FCC gold
the primitive vectors were given in Chapter~\ref{ch:intro}; the primitive
cell volume is $\Omega = |\mathbf a_1 \cdot (\mathbf a_2 \times \mathbf
a_3)| = a^3/4$.

The \textbf{reciprocal lattice} consists of all vectors $\GG$ such that
$\ee^{\ii \GG\cdot\RR} = 1$ for every lattice translation $\RR$. Its
primitive vectors are
\begin{equation}
\mathbf b_1 = \frac{2\pi\, (\mathbf a_2 \times \mathbf a_3)}{\Omega},
\qquad
\mathbf b_2 = \frac{2\pi\, (\mathbf a_3 \times \mathbf a_1)}{\Omega},
\qquad
\mathbf b_3 = \frac{2\pi\, (\mathbf a_1 \times \mathbf a_2)}{\Omega},
\qquad
\mathbf a_i \cdot \mathbf b_j = 2\pi\, \delta_{ij}.
\end{equation}
(In matrix language: if $A$ has rows $\mathbf a_i$ then $B = 2\pi
(A^{-1})^{\mathsf T}$ has rows $\mathbf b_i$ --- implement it that way.) The
reciprocal of FCC is BCC. The Wigner--Seitz cell of the reciprocal lattice is
the \textbf{first Brillouin zone} (BZ); for FCC it is the truncated
octahedron with the high-symmetry points $\Gamma = (0,0,0)$,
$X = \tfrac{2\pi}{a}(1,0,0)$, $L = \tfrac{\pi}{a}(1,1,1)$,
$W = \tfrac{2\pi}{a}(1,\tfrac12,0)$, $K = \tfrac{2\pi}{a}(\tfrac34,
\tfrac34, 0)$, along which you will plot gold's band structure.

\section{Bloch's theorem}

The Kohn--Sham potential of a crystal is lattice-periodic,
$\veff(\rr + \RR) = \veff(\rr)$.

\begin{theorem}[Bloch]
The eigenstates of a Hamiltonian with a lattice-periodic potential can be
chosen as
\begin{equation}
\psi_{n\kk}(\rr) = \ee^{\ii \kk\cdot\rr}\, u_{n\kk}(\rr),
\qquad
u_{n\kk}(\rr + \RR) = u_{n\kk}(\rr),
\label{eq:bloch}
\end{equation}
labeled by a wavevector $\kk$ in the first Brillouin zone and a discrete band
index $n$.
\end{theorem}

\begin{proof}
Let $\hat T_{\RR}$ be the translation operator, $(\hat T_{\RR}\psi)(\rr) =
\psi(\rr + \RR)$. Periodicity of the potential means $[\Ha, \hat T_{\RR}] =
0$ for all $\RR$, and the $\hat T_{\RR}$ commute among themselves, so all can
be simultaneously diagonalized with $\Ha$. Since $\hat T_{\RR}$ is unitary
its eigenvalues have modulus one, and the group law $\hat T_{\RR}\hat
T_{\RR'} = \hat T_{\RR + \RR'}$ forces the eigenvalue to be an exponential of
a linear function of $\RR$: $\hat T_{\RR} \psi = \ee^{\ii\kk\cdot\RR}\psi$
for some real vector $\kk$. Defining $u(\rr) \equiv
\ee^{-\ii\kk\cdot\rr}\psi(\rr)$, one checks $u(\rr+\RR) = u(\rr)$. Finally,
$\kk$ and $\kk + \GG$ give the same eigenvalue of every $\hat T_{\RR}$, so
$\kk$ may be restricted to the first BZ.
\end{proof}

Physically: in a perfect crystal, momentum is conserved modulo reciprocal
lattice vectors; $\kk$ (``crystal momentum'') is the corresponding quantum
number. The spectrum at each $\kk$ is discrete ($\varepsilon_{n\kk}$), and
plotted along BZ paths gives the \emph{band structure}. Bloch's theorem is
also a factorization of the computation: \textbf{the giant eigenproblem of
the infinite crystal decouples into independent moderate-size eigenproblems,
one per $\kk$-point} --- the single most important source of parallelism in
your program.

\section{The plane-wave basis}

Since $u_{n\kk}$ is cell-periodic, expand it in the discrete Fourier series
of the lattice:
\begin{equation}
\psi_{n\kk}(\rr) = \frac{1}{\sqrt{\Omega}} \sum_{\GG}
c_{n\kk}(\GG)\; \ee^{\ii(\kk + \GG)\cdot \rr}.
\label{eq:pwexpansion}
\end{equation}
The basis functions $\ket{\kk+\GG}$ are orthonormal over the cell. The
expansion is truncated by keeping every $\GG$ with kinetic energy below a
\textbf{cutoff}:
\begin{equation}
\tfrac12 |\kk + \GG|^2 \leq E_{\mathrm{cut}} .
\label{eq:ecut}
\end{equation}
This is the \emph{only} convergence parameter of the basis: increasing
$E_{\mathrm{cut}}$ enlarges the variational space, so the total energy
decreases monotonically toward its converged value (other quantities ---
eigenvalue differences, lattice constants --- also converge, though not
necessarily monotonically).
The number of plane waves is, by counting $\GG$-points inside a sphere,
\begin{equation}
N_{\mathrm{pw}} \approx \frac{\Omega}{6\pi^2}\,
\big(2E_{\mathrm{cut}}\big)^{3/2}.
\end{equation}
For gold with an HGH pseudopotential, $E_{\mathrm{cut}} \approx
40$--$60$\,Ha; with $\Omega \approx 114$\,Bohr$^3$ this gives $N_{\mathrm{pw}}
\sim 1400$--$2500$ coefficients per band per spin component --- small dense
problems, ideal for GPU batching.

Plane waves have three cardinal virtues: (i) they are unbiased --- no atom-
centered basis-set superposition errors; (ii) matrix elements of the kinetic
operator and of any multiplicative potential are trivial (diagonal, and
convolution, respectively); (iii) the map between the coefficient
representation $c(\GG)$ and the real-space representation $u(\rr)$ on a
uniform grid is the FFT, costing $O(N \log N)$. Their vice --- inability to
resolve the violent oscillations of wavefunctions near nuclei --- is cured by
pseudopotentials (Chapter~\ref{ch:pseudo}).

In this basis the Kohn--Sham equation \eqref{eq:kseq} becomes an ordinary
matrix eigenproblem for the coefficient vector at each $\kk$,
\begin{equation}
\sum_{\GG'} H_{\GG\GG'}(\kk)\, c_{n\kk}(\GG') = \varepsilon_{n\kk}\,
c_{n\kk}(\GG),
\qquad
H_{\GG\GG'}(\kk) = \tfrac12|\kk+\GG|^2\,\delta_{\GG\GG'}
+ V(\GG - \GG') + \dots,
\label{eq:pwmatrix}
\end{equation}
where $V(\GG)$ are the Fourier components of the local potential
(convention \eqref{eq:fourierconv} below) and the dots stand for the
nonlocal pseudopotential of Chapter~\ref{ch:pseudo}. The matrix is never
stored; Chapter~\ref{ch:hamiltonian} shows how to apply it.

\paragraph{Fourier conventions (fix these once, use them everywhere).}
Every cell-periodic function $g(\rr)$ --- density, potential, $u_{n\kk}$ ---
is expanded as
\begin{equation}
g(\rr) = \sum_{\GG} g(\GG)\, \ee^{\ii\GG\cdot\rr},
\qquad
g(\GG) = \frac{1}{\Omega}\int_{\Omega} g(\rr)\, \ee^{-\ii\GG\cdot\rr}\,\dd\rr .
\label{eq:fourierconv}
\end{equation}
With this convention $g(\GG{=}0)$ is the cell average of $g$, and
$\int_\Omega g^* h\,\dd\rr = \Omega\sum_{\GG} g^*(\GG) h(\GG)$. On an
$N_1{\times}N_2{\times}N_3$ grid with points
$\rr_{j} = \sum_i (j_i/N_i)\,\mathbf a_i$, Eq.~\eqref{eq:fourierconv} is
exactly the discrete Fourier transform:
$g(\GG) = \tfrac{1}{N}\,\mathrm{FFT}[g(\rr_j)]$ and
$g(\rr_j) = N\cdot\mathrm{IFFT}[g(\GG)]$ with $N = N_1N_2N_3$ and the usual
library normalization (forward unnormalized, inverse divided by $N$).
Integer FFT frequency $(m_1,m_2,m_3)$ corresponds to
$\GG = \sum_i m_i \mathbf b_i$, with $m_i$ running over
$-N_i/2 \dots N_i/2 - 1$ in the library's wrapped order. For wavefunctions
the coefficients in Eq.~\eqref{eq:pwexpansion} carry an extra
$1/\sqrt\Omega$ so that normalization is simply
$\sum_{\GG,s} |c^s_{n\kk}(\GG)|^2 = 1$; then
$u_{n\kk}(\rr_j) = \tfrac{N}{\sqrt\Omega}\,\mathrm{IFFT}[c]$ and the density
contribution of the state is $|u_{n\kk}(\rr_j)|^2$ in electrons per
Bohr$^3$.

\paragraph{Densities need a doubled sphere.} The density $n(\rr) \sim
|\psi|^2$ contains products of plane waves, hence Fourier components up to
$2\sqrt{2E_{\mathrm{cut}}}$: represent $n$, $\veff$, and all potentials on a
grid/sphere with cutoff $4E_{\mathrm{cut}}$ (``dense grid''). Choosing the
FFT grid at least this large avoids aliasing errors when forming
$n$ and when multiplying $\veff \times \psi$ in real space.

\section{A first band gap: the two-plane-wave model}
\label{sec:nfe}

Before any machinery, solve Eq.~\eqref{eq:pwmatrix} by hand in the one
situation where two plane waves suffice. It explains where band gaps come
from, it is the first nontrivial test of your code (Chapter~\ref{ch:verify},
Stage~1), and --- with one extra observation about symmetry --- it is the
seed of the topological property computed in Chapter~\ref{ch:topology}.

Take a weak local potential and a $\kk$ on the zone boundary, halfway to a
reciprocal lattice vector: $\kk = \tfrac12\GG_1$ (in FCC, $X$ and $L$ are
such points). The two plane waves $\ket{\kk}$ and $\ket{\kk - \GG_1}$ have
wavevectors $+\tfrac12\GG_1$ and $-\tfrac12\GG_1$: the same kinetic energy
$T = \tfrac18|\GG_1|^2$. Degenerate states are mixed by any perturbation
that connects them, however weak, and the potential does so through its
Fourier component $V_1 \equiv V(\GG_1)$. Keeping just these two basis
functions, Eq.~\eqref{eq:pwmatrix} is
\begin{equation}
\begin{pmatrix} T & V_1 \\ V_1^* & T \end{pmatrix}
\begin{pmatrix} c(\mathbf 0) \\ c(-\GG_1)\end{pmatrix}
= \varepsilon \begin{pmatrix} c(\mathbf 0) \\ c(-\GG_1)\end{pmatrix}
\qquad\Longrightarrow\qquad
\varepsilon_\pm = T \pm |V_1| .
\end{equation}
A gap $2|V_1|$ opens: this is the origin of every band gap at a zone
boundary. Now the observation. If the crystal has an atom at the origin and
inversion symmetry (as gold does), $V(\rr) = V(-\rr)$ makes $V_1$ real, and
the eigenvectors are $(1, \pm1)/\sqrt2$ --- in real space
\begin{equation}
\begin{aligned}
\psi_{\mathrm{even}} &\propto \ee^{\ii\GG_1\cdot\rr/2} + \ee^{-\ii\GG_1\cdot\rr/2}
\propto \cos\!\big(\tfrac12\GG_1\cdot\rr\big),
&\qquad \varepsilon_{\mathrm{even}} &= T + V_1,\\
\psi_{\mathrm{odd}} &\propto \ee^{\ii\GG_1\cdot\rr/2} - \ee^{-\ii\GG_1\cdot\rr/2}
\propto \sin\!\big(\tfrac12\GG_1\cdot\rr\big),
&\qquad \varepsilon_{\mathrm{odd}} &= T - V_1 .
\end{aligned}
\label{eq:nfeparity}
\end{equation}
The cosine state is \emph{even} under $\rr \to -\rr$ and piles its charge
on the atoms, like an $s$ orbital; the sine state is \emph{odd}, has a node
on every atom and its charge between them, like a $p$ orbital. Which one
is lower is decided by the sign of a single number:
\begin{itemize}[leftmargin=2em]
\item $V_1 < 0$ (the electron is attracted to the atomic site at this length
scale): the even, $s$-like state is below the odd, $p$-like one --- the
order you would guess from atomic physics, where $s$ levels lie below $p$
levels.
\item $V_1 > 0$: the order is \emph{inverted}; the odd state is lower.
\end{itemize}
Both signs occur in real materials, because the effective potential felt by
valence electrons is not simply attractive (Chapter~\ref{ch:pseudo}), and
because other bands of the same parity can push one of the two states.
Remember this picture: an integer-valued property of gold's band structure
will turn out to be nothing more than a record of which of these two
orderings is realized at $X$ and at $L$.

\section{Brillouin-zone integration}
\label{sec:bzintegration}

Every observable of the crystal is a BZ average; e.g.\ the density and
band energy are
\begin{equation}
n(\rr) = \sum_n \int_{\mathrm{BZ}} \frac{\dd\kk}{\Omega_{\mathrm{BZ}}}\,
f_{n\kk}\, |\psi_{n\kk}(\rr)|^2,
\qquad
E_{\mathrm{band}} = \sum_n \int_{\mathrm{BZ}}
\frac{\dd\kk}{\Omega_{\mathrm{BZ}}}\, f_{n\kk}\, \varepsilon_{n\kk}.
\end{equation}
Numerically the integral becomes a weighted sum over a finite mesh. The
standard \textbf{Monkhorst--Pack} mesh is the uniform grid
\begin{equation}
\kk_{(m_1 m_2 m_3)} = \sum_{i=1}^{3} \frac{2 m_i - N_i - 1}{2 N_i}\,
\mathbf b_i,
\qquad m_i = 1, \dots, N_i,
\label{eq:mpmesh}
\end{equation}
i.e.\ a uniform $N_1{\times}N_2{\times}N_3$ grid centered on the origin. For
odd $N_i$ it contains $\Gamma$; for even $N_i$ it straddles $\Gamma$
symmetrically without containing it (the fractional coordinates are
$\pm\tfrac{1}{2N}, \pm\tfrac{3}{2N}, \dots$). The alternative
``$\Gamma$-centered'' grid uses fractional coordinates $m_i/N_i$,
$m_i = 0,\dots,N_i{-}1$, and always contains $\Gamma$. Both converge to
the same answer as $N_i\to\infty$, but for the FCC lattice they are
\emph{not} equally good at finite $N_i$: the $\Gamma$-centered grid is
mapped onto itself by all 48 operations of the cubic point group, whereas
the even-$N$ grid of Eq.~\eqref{eq:mpmesh}, built along the primitive
$\mathbf b_i$, is not (its offset
$\tfrac{1}{2N}(\mathbf b_1{+}\mathbf b_2{+}\mathbf b_3)$ is not invariant
modulo the grid spacing). A symmetry-breaking mesh produces a density of
slightly lower symmetry, and the damage is visible: in a $4^3$
calculation for gold the fourfold-degenerate $d$ levels at $\Gamma$ split
by $8$\,meV and the four equivalent $L$ points come out with energies
differing by $22$\,meV; with an $8^3$ mesh the same defects are still
about $1$ and $2$\,meV. With a $\Gamma$-centered mesh they vanish to
solver accuracy. \textbf{Use $\Gamma$-centered meshes for
gold.} They have a second advantage: for even $N_i$ they contain all eight
points $\tfrac12(n_1,n_2,n_3)$ that Chapter~\ref{ch:topology} needs.
Point-group symmetry can fold the mesh to the irreducible wedge with integer
weights; \emph{for a first code, skip symmetry reduction} --- use the full
mesh with equal weights $w_{\kk} = 1/N_{\kk}$. It is simpler, impossible to
get subtly wrong, and on a GPU the extra $\kk$-points are nearly free because
they batch perfectly, and symmetry-related points \emph{must} then come out
with identical eigenvalues, which is a useful test. (One reduction is free
and safe: time reversal maps
$\kk \to -\kk$ with identical eigenvalues and density,
Section~\ref{sec:symmetry}, so you may keep one point of each $\pm\kk$ pair
with doubled weight.) Gold needs a dense mesh ($12^3$ to $16^3$ with
smearing) for well-converged total energies, because the integrand jumps
at the Fermi surface; an $8^3$ mesh is already adequate for band
structures and for debugging. Coordinates:
store each $\kk$ as its three fractional components along
$\mathbf b_{1,2,3}$ and convert to Cartesian with $\kk_{\mathrm{cart}} =
\kk_{\mathrm{frac}} B$ (rows of $B$ are the $\mathbf b_i$). The
high-symmetry points quoted above are, in fractional form,
$\Gamma = (0,0,0)$, $X = (0, \tfrac12, \tfrac12)$,
$L = (\tfrac12,\tfrac12,\tfrac12)$,
$W = (\tfrac14,\tfrac12,\tfrac34)$,
$K = (\tfrac38,\tfrac38,\tfrac34)$ for the $\mathbf a_i$ of
Chapter~\ref{ch:intro}.

\section{Metals: fractional occupations and smearing}
\label{sec:smearing}

In an insulator, bands are either full or empty and $f_{n\kk} \in \{0,1\}$
(per spinor state; no factor 2 once we use spinors). In a metal the
occupation at $T=0$ is a step function $f = \theta(\mu -
\varepsilon_{n\kk})$: the integrand of every BZ integral is discontinuous
across the Fermi surface, and a discrete $\kk$-sum converges miserably
($\sim 1/N_{\kk}$, with wild oscillations). Worse, states at the Fermi level
flip between occupied and empty from one SCF iteration to the next, making
self-consistency unstable.

The cure is \textbf{smearing}: replace the step by a smooth function of width
$\sigma$,
\begin{equation}
f_{n\kk} = f\!\left(\frac{\varepsilon_{n\kk} - \mu}{\sigma}\right).
\end{equation}
Two standard choices to implement:
\begin{enumerate}[leftmargin=2em]
\item \textbf{Fermi--Dirac}: $f(x) = 1/(\ee^{x} + 1)$. Physically a finite
electronic temperature $T = \sigma/k_B$. The associated entropy per state is
$s = -[f \ln f + (1-f)\ln(1-f)]$.
\item \textbf{Gaussian / Methfessel--Paxton}: $f_0(x) = \tfrac12
\mathrm{erfc}(x)$, and the MP family adds Hermite-polynomial corrections,
$f_1(x) = f_0(x) - \tfrac{x}{2\sqrt\pi} \ee^{-x^2}$ (order 1), which makes the
smeared free energy insensitive to $\sigma$ to higher order --- the usual
choice for metals like gold, with $\sigma \approx 0.01$\,Ha.
\end{enumerate}

\paragraph{The free energy.} With smearing, the variational quantity is not
$E$ but the \textbf{Mermin free energy}
\begin{equation}
F = E - \sum_{n\kk} w_{\kk}\, \sigma\, s(f_{n\kk}),
\label{eq:mermin}
\end{equation}
which is what must decrease during SCF and what enters total-energy
comparisons at fixed $\sigma$. Here $s$ is evaluated at $x = (\varepsilon_{n\kk}-\mu)/\sigma$. For the
non-thermal schemes $s$ is a ``generalized entropy'' defined so that $F$ is
still variational:
$s_0(x) = \ee^{-x^2}/(2\sqrt{\pi})$ for Gaussian and
$s_1(x) = (1 - 2x^2)\,\ee^{-x^2}/(4\sqrt{\pi})$ for Methfessel--Paxton
order~1 (which can be negative --- it is not a physical entropy). The best
estimate of the $\sigma\to0$ energy is $E_0 \approx \tfrac12 (E + F)$ for
Fermi--Dirac and Gaussian smearing; for Methfessel--Paxton $F$ itself is
already the estimate, which is the point of the scheme.

\paragraph{Finding the Fermi level.} $\mu$ is fixed by charge conservation:
\begin{equation}
N_e = \sum_{n\kk} w_{\kk}\,
f\!\left(\frac{\varepsilon_{n\kk} - \mu}{\sigma}\right).
\label{eq:fermibisect}
\end{equation}
The right-hand side is monotonically increasing in $\mu$ for Fermi--Dirac
and Gaussian smearing (for MP orders $\ge 1$ the occupation function
overshoots and $N(\mu)$ can decrease slightly across gapped regions of the
spectrum; bisection between the brackets below still converges to a valid
root, but do not assume uniqueness when debugging); solve by bisection
(or Newton with bisection fallback) between $\min \varepsilon - 10\sigma$ and
$\max\varepsilon + 10\sigma$ to $\sim 10^{-12}$. Include enough
\emph{empty bands} in the eigensolver: 11-electron gold fills, on average,
11 spinor bands per $\kk$-point, so solve for $\geq 20$ bands so that
the smeared occupations have room to spill and the eigensolver has a buffer
(Chapter~\ref{ch:algorithms}).

\begin{gpubox}{GPU/JAX notes for this chapter}
The $(\kk, n)$ pair is your fundamental batch axis. Store all coefficients as
one array of shape $[N_k, N_{\mathrm{bands}}, 2, N_G^{\max}]$ (spinor
dimension 2), padding each $\kk$-point's sphere to a common $N_G^{\max}$ with
a boolean mask (JAX demands static shapes; masked padding wastes a few
percent and buys perfect batching). Fermi search
\eqref{eq:fermibisect} is a scalar bisection over a reduction of the whole
$\varepsilon$ array --- do it on device with \texttt{lax.while\_loop}; a
fixed 60--80 iteration \texttt{fori\_loop} bisection is even simpler and
JIT-friendly. Smearing functions must be implemented with numerically safe
forms ($\mathrm{erfc}$, \texttt{log1p}, clipping $x$ to $\pm 40$) to avoid
NaNs poisoning gradients.
\end{gpubox}

\chapter{Pseudopotentials}
\label{ch:pseudo}

\section{Why plane waves cannot see a gold nucleus}

Near a nucleus of charge $Z$, the wavefunction behaves like a hydrogenic
orbital with length scale $\sim 1/Z$ Bohr; moreover valence states must
oscillate rapidly there to remain orthogonal to the core states (gold has 68
core electrons, $1s$ through $5p$, below its 11 valence electrons). Resolving structure on the scale
$1/79$\,Bohr with plane waves would require $E_{\mathrm{cut}} \sim Z^2 \sim
6000$\,Ha and $\sim 3\times10^6$ plane waves per band --- hopeless. But this violent
region is chemically inert: bonding, band formation, and Fermi-surface
physics are governed by the smooth tails of the valence states between atoms.

The \textbf{pseudopotential} idea: replace the nucleus + frozen core by an
effective, angular-momentum-dependent potential that is (i) identical to the
true (screened) potential outside a core radius $r_c$, and (ii) weak and
smooth inside, such that its \emph{pseudo-wavefunctions} are nodeless, smooth,
and match the true valence wavefunctions exactly for $r > r_c$, with the same
eigenvalues.

\section{Norm conservation}

Modern norm-conserving pseudopotentials (Hamann--Schl\"uter--Chiang 1979) are
constructed, per angular momentum $\ell$, from an all-electron atomic
calculation by demanding, at the atomic reference energy
$\varepsilon_\ell$:
\begin{enumerate}[leftmargin=2em]
\item $\psi^{\mathrm{PS}}_\ell(r) = \psi^{\mathrm{AE}}_\ell(r)$ for $r > r_c$;
pseudo-eigenvalue $=$ all-electron eigenvalue.
\item \textbf{Norm conservation}: $\displaystyle\int_0^{r_c}
|\psi^{\mathrm{PS}}_\ell|^2 r^2 \dd r = \int_0^{r_c}
|\psi^{\mathrm{AE}}_\ell|^2 r^2 \dd r$ --- the pseudo-orbital carries the
correct amount of charge inside the core.
\end{enumerate}
Norm conservation is not cosmetic: by an identity relating the energy
derivative of the logarithmic derivative to the enclosed norm,
\begin{equation}
-\frac{\partial}{\partial \varepsilon}
\frac{\partial}{\partial r} \ln \psi_\ell(r, \varepsilon)\Big|_{r_c}
\;=\; \frac{2}{r_c^2\, |\psi_\ell(r_c)|^2} \int_0^{r_c}
|\psi_\ell|^2\, r^2\, \dd r,
\end{equation}
it guarantees that the pseudo-atom scatters electrons like the real atom not
only \emph{at} $\varepsilon_\ell$ but to first order \emph{around} it ---
the transferability that lets an atomic construction describe a solid.

\section{Semilocal form and the Kleinman--Bylander transformation}

The raw construction yields a \emph{semilocal} operator --- local in $r$,
different for each $\ell$:
\begin{equation}
\hat V_{\mathrm{PS}} = V_{\mathrm{loc}}(r)
+ \sum_{\ell m} \ket{\ell m}\, \Delta V_\ell(r)\, \bra{\ell m},
\qquad \Delta V_\ell = V_\ell - V_{\mathrm{loc}},
\end{equation}
where $\ket{\ell m}$ projects onto angular momentum about the atom. Matrix
elements of the semilocal part between $N_{\mathrm{pw}}$ plane waves cost
$O(N_{\mathrm{pw}}^2)$ --- too much. Kleinman and Bylander (1982) observed
that $\Delta V_\ell$ only needs to be correct when acting on wavefunctions
resembling the atomic reference $\phi_\ell$, and replaced it by a rank-one
\emph{fully separable} operator per channel:
\begin{equation}
\Delta \hat V_{\ell}^{\mathrm{KB}} =
\frac{\ket{\Delta V_\ell\, \phi_\ell} \bra{\phi_\ell\, \Delta V_\ell}}
{\braket{\phi_\ell | \Delta V_\ell | \phi_\ell}},
\end{equation}
which by construction reproduces the semilocal action on $\phi_\ell$
exactly. The general separable (multi-projector) form your code will
implement is
\begin{equation}
\boxed{\;
\hat V_{\mathrm{NL}} = \sum_{a}\sum_{\ell m}\sum_{ij}
\ket{\beta^{a}_{i\ell m}}\; h^{\ell}_{ij}\; \bra{\beta^{a}_{j \ell m}},
\;}
\label{eq:separable}
\end{equation}
a sum over atoms $a$, channels $\ell m$, and a small matrix $h^{\ell}_{ij}$
coupling a few radial projectors $\beta_{i\ell m}(\rr) = p_{i\ell}(r)\,
Y_{\ell m}(\hat \rr)$. Its action on a state costs only projections
$\braket{\beta | \psi}$ --- tall-skinny matrix multiplications, the second
GPU-native kernel of the code.

\section{HGH pseudopotentials: analytic, relativistic, perfect for from-scratch}
\label{sec:hgh}

For a self-contained implementation you want a pseudopotential defined by a
\emph{short table of numbers and closed-form functions}, not by files with
radial grids. The Goedecker--Teter--Hutter / \textbf{Hartwigsen--Goedecker--%
Hutter (HGH, 1998)} family is exactly that, is norm-conserving, is already in
separable form, has \emph{analytic Fourier transforms}, and --- decisively for
us --- the HGH paper tabulates \textbf{fully relativistic parameters
including spin--orbit projectors for every element from H to Rn}, gold
included.

\paragraph{Local part.} A smeared Coulomb tail plus a short-range Gaussian
polynomial:
\begin{equation}
V_{\mathrm{loc}}(r) = -\frac{Z_{\mathrm{ion}}}{r}\,
\mathrm{erf}\!\left(\frac{r}{\sqrt 2\, r_{\mathrm{loc}}}\right)
+ \ee^{-\frac{r^2}{2 r_{\mathrm{loc}}^2}}
\Big[C_1 + C_2 \tilde r^{\,2} + C_3 \tilde r^{\,4} + C_4 \tilde r^{\,6}\Big],
\qquad \tilde r = r/r_{\mathrm{loc}} .
\label{eq:hghloc}
\end{equation}
Its Fourier transform (what you actually need, Chapter~\ref{ch:hamiltonian})
is analytic:
\begin{equation}
\tilde V_{\mathrm{loc}}(G) =
-\frac{4\pi Z_{\mathrm{ion}}}{G^2}\, \ee^{-\frac{(G r_{\mathrm{loc}})^2}{2}}
+ (2\pi)^{3/2} r_{\mathrm{loc}}^3\, \ee^{-\frac{(G r_{\mathrm{loc}})^2}{2}}
\times \big[\text{polynomial in } (G r_{\mathrm{loc}})^2
\text{ with } C_1..C_4\big],
\label{eq:hghlocG}
\end{equation}
where, with $x = G r_{\mathrm{loc}}$, the polynomial is
\begin{equation}
C_1 + C_2\,(3 - x^2) + C_3\,(15 - 10x^2 + x^4)
+ C_4\,(105 - 105x^2 + 21x^4 - x^6).
\label{eq:hghlocpoly}
\end{equation}
Here $\tilde V_{\mathrm{loc}}(G) \equiv \int V_{\mathrm{loc}}(r)\,
\ee^{-\ii\GG\cdot\rr}\,\dd\rr$ is the transform of \emph{one atom's}
potential over all space (no $1/\Omega$ yet; that enters in
Eq.~\eqref{eq:structfact}).
The $-4\pi Z/G^2$ piece diverges at $G=0$ like the bare Coulomb potential;
it cancels analytically against Hartree and Ewald divergences
(Section~\ref{sec:g0}).

\paragraph{Nonlocal part.} Gaussian-type radial projectors
\begin{equation}
p_{i\ell}(r) = \frac{\sqrt 2\; r^{\ell + 2(i-1)}\,
\ee^{-r^2/(2 r_\ell^2)}}
{r_\ell^{\,\ell + (4i-1)/2}\,
\sqrt{\Gamma\big(\ell + \tfrac{4i-1}{2}\big)}},
\qquad i = 1,2,3,
\label{eq:hghproj}
\end{equation}
with coupling matrices $h^{\ell}_{ij}$ --- and, in the relativistic version,
additional \textbf{spin--orbit coupling matrices} $k^{\ell}_{ij}$ entering as
\begin{equation}
\hat V_{\mathrm{NL}} = \sum_{\ell m, ij} \ket{\beta_{i\ell m}}
h^{\ell}_{ij} \bra{\beta_{j\ell m}}
\;+\; \sum_{\ell m m', ij} \ket{\beta_{i\ell m}}\,
k^{\ell}_{ij}\; \big(\mathbf{L}\cdot\mathbf{S}\big)^{\ell}_{mm'}
\,\bra{\beta_{j\ell m'}} ,
\label{eq:hghso}
\end{equation}
where $\mathbf L \cdot \mathbf S$ acts on the product of orbital ($m$) and
spinor indices --- the full derivation and matrix elements are in
Chapter~\ref{ch:soc}. The radial functions are normalized,
$\int_0^\infty p_{i\ell}^2\, r^2\,\dd r = 1$.

\paragraph{Projectors in reciprocal space.} A plane wave expands in
spherical waves (the Rayleigh formula),
\begin{equation}
\ee^{\ii\qq\cdot\rr} = 4\pi \sum_{\ell=0}^{\infty}\sum_{m=-\ell}^{\ell}
\ii^{\ell}\, j_\ell(qr)\, Y_{\ell m}(\hat\qq)\, Y_{\ell m}(\hat\rr),
\label{eq:rayleigh}
\end{equation}
with $j_\ell$ the spherical Bessel functions and $Y_{\ell m}$ \emph{real}
spherical harmonics (Table~\ref{tab:realylm}). Insert this into the overlap
of a projector $\beta_{i\ell m}(\rr) = p_{i\ell}(r) Y_{\ell m}(\hat\rr)$,
centered on an atom at $\RR_a$, with the normalized plane wave
$\ee^{\ii\qq\cdot\rr}/\sqrt\Omega$, $\qq = \kk+\GG$. Orthonormality of the
$Y_{\ell m}$ kills every term but one:
\begin{equation}
\tilde\beta^{a}_{i\ell m}(\qq) \equiv \braket{\qq|\beta^{a}_{i\ell m}}
= \frac{1}{\sqrt\Omega}\int \ee^{-\ii\qq\cdot\rr}\,
\beta_{i\ell m}(\rr - \RR_a)\,\dd\rr
= \frac{4\pi}{\sqrt\Omega}\,(-\ii)^{\ell}\;
\tilde p_{i\ell}(q)\; Y_{\ell m}(\hat\qq)\;\ee^{-\ii\qq\cdot\RR_a} .
\label{eq:betaG}
\end{equation}
So the overlap reduces to a one-dimensional \emph{spherical Bessel
transform} of the radial part,
\begin{equation}
\tilde p_{i\ell}(q) \equiv \int_0^\infty p_{i\ell}(r)\, j_\ell(qr)\, r^2\,
\dd r ,
\label{eq:besseltransform}
\end{equation}
and for Gaussians these integrals are elementary. With $x = q r_\ell$ and
the common factor $\mathcal E = \pi^{1/4}\, \ee^{-x^2/2}$:
\begin{equation}
\begin{aligned}
\tilde p_{1,s} &= \mathcal E \sqrt{2 r_\ell^3},
&
\tilde p_{2,s} &= \mathcal E\, 2\sqrt{\tfrac{2 r_\ell^3}{15}}\,(3 - x^2),
&
\tilde p_{3,s} &= \mathcal E\, \tfrac43\sqrt{\tfrac{2 r_\ell^3}{105}}\,
(15 - 10x^2 + x^4),
\\
\tilde p_{1,p} &= \mathcal E\, 2\sqrt{\tfrac{r_\ell^5}{3}}\; q,
&
\tilde p_{2,p} &= \mathcal E\, 4\sqrt{\tfrac{r_\ell^5}{105}}\; q\,(5 - x^2),
&
\tilde p_{3,p} &= \mathcal E\, \tfrac83\sqrt{\tfrac{r_\ell^5}{1155}}\; q\,
(35 - 14x^2 + x^4),
\\
\tilde p_{1,d} &= \mathcal E\, 2\sqrt{\tfrac{2 r_\ell^7}{15}}\; q^2,
&
\tilde p_{2,d} &= \mathcal E\, \tfrac43\sqrt{\tfrac{2 r_\ell^7}{105}}\;
q^2\,(7 - x^2),
&
\tilde p_{3,d} &= \mathcal E\, \tfrac83\sqrt{\tfrac{2 r_\ell^7}{15015}}\;
q^2\,(63 - 18x^2 + x^4).
\end{aligned}
\label{eq:hghprojG}
\end{equation}
(Each line was checked against direct numerical integration of
Eq.~\eqref{eq:besseltransform}; do the same in your unit tests.) Gaussians
times polynomials in $q$ --- evaluate them directly on your $|\kk+\GG|$
grid, no radial integration, no splines, no file parsing. This is
why HGH is the right choice for a from-scratch code: \emph{the entire ionic
problem for gold is about twenty published numbers plus closed-form
functions}.

\paragraph{Off-diagonal couplings.} The radial functions $p_{1\ell}$ and
$p_{2\ell}$ are each normalized but \emph{not orthogonal to each other}
($\int p_{1\ell}\,p_{2\ell}\, r^2 \dd r = \sqrt{3/5}, \sqrt{5/7}, \sqrt{7/9}$
for $\ell = 0,1,2$). Consequently the matrices $h^\ell_{ij}$ and
$k^\ell_{ij}$, which are symmetric, have off-diagonal elements that are
\emph{not} zero. In the original HGH paper only the diagonal elements are
tabulated and the off-diagonal ones are fixed multiples of them, e.g.\ for
two projectors
\begin{equation}
h^{s}_{12} = -\tfrac12\sqrt{\tfrac35}\, h^{s}_{22},
\qquad
h^{p}_{12} = -\tfrac12\sqrt{\tfrac57}\, h^{p}_{22},
\qquad
h^{d}_{12} = -\tfrac12\sqrt{\tfrac79}\, h^{d}_{22},
\label{eq:hghoffdiag}
\end{equation}
and likewise for $k_{12}$. Later parameter sets in the same family (the PBE
set below) follow a different convention --- there the ratios are exactly
twice these (check: $-2.25067966/2.90561495 = -\sqrt{3/5}$). To make this
impossible to get wrong,
Table~\ref{tab:auhgh} lists \emph{every} matrix element explicitly: use
the tabulated matrices verbatim and derive nothing. Forgetting the
off-diagonal elements is the classic HGH bug: the code runs and gives
plausible but wrong bands.

\paragraph{Gold parameters.} Table~\ref{tab:auhgh} lists the complete
11-electron parameter sets. Gold uses two projectors ($i = 1,2$) in each of
the $s$, $p$, $d$ channels, so there are $2\,(1+3+5) = 18$ projector
functions per atom, and only $C_1$ is nonzero in the local part. Two
choices of valence: the 11-electron potential ($5d^{10}6s^1$, cheaper, fine for bands
near $E_F$) and the 19-electron semicore potential
($5s^2 5p^6 5d^{10} 6s^1$, stiffer, $E_{\mathrm{cut}}$ higher, better lattice
constant; parameters in Appendix~\ref{app:q19}, where a third $s$ and $p$
projector and a second local coefficient appear). Start with $q=11$;
graduate to $q=19$ when validating energetics.
HGH parameters exist in matched LDA and PBE flavors --- use the one matching
your XC functional.

\begin{table}[h]
\centering\small
\begin{tabular}{@{}llrrr@{}}
\toprule
\multicolumn{5}{@{}l}{\textbf{LDA} (HGH 1998; entry \texttt{Au GTH-PADE-q11}):\quad
$r_{\mathrm{loc}} = 0.59000000$,\ \ $C_1 = 11.60442790$}\\
\midrule
channel & & \multicolumn{1}{c}{$11$} & \multicolumn{1}{c}{$12 = 21$} & \multicolumn{1}{c}{$22$}\\
\midrule
$s$:\ $r_s = 0.52117975$ & $h_{ij}$ & $2.53861359$ & $-1.04613664$ & $2.70111320$\\
\addlinespace
$p$:\ $r_p = 0.63061306$ & $h_{ij}$ & $0.39485270$ & $-0.86959213$ & $2.05783056$\\
 & $k_{ij}$ & $-0.96005489$ & $-0.54790119$ & $1.29657086$\\
\addlinespace
$d$:\ $r_d = 0.44070643$ & $h_{ij}$ & $-4.71906966$ & $0.72777081$ & $-1.65042907$\\
 & $k_{ij}$ & $0.14848354$ & $0.07473921$ & $-0.16949260$\\
\midrule
\multicolumn{5}{@{}l}{\textbf{PBE} (Krack 2005; entry \texttt{Au GTH-PBE-q11}):\quad
$r_{\mathrm{loc}} = 0.59000000$,\ \ $C_1 = 10.51717881$}\\
\midrule
$s$:\ $r_s = 0.54384733$ & $h_{ij}$ & $5.78681103$ & $-2.25067966$ & $2.90561495$\\
\addlinespace
$p$:\ $r_p = 0.60917298$ & $h_{ij}$ & $4.28404873$ & $-4.07128520$ & $4.81720962$\\
 & $k_{ij}$ & $1.21911104$ & $-2.38187511$ & $2.81827264$\\
\addlinespace
$d$:\ $r_d = 0.43730922$ & $h_{ij}$ & $-7.40328124$ & $3.01019307$ & $-3.41323811$\\
 & $k_{ij}$ & $-0.12113041$ & $0.26246301$ & $-0.29760507$\\
\bottomrule
\end{tabular}
\caption{HGH/GTH pseudopotential parameters for Au with
$Z_{\mathrm{ion}} = 11$ (Hartree atomic units), transcribed with all
digits from the \texttt{GTH\_SOC\_POTENTIALS} database distributed with the
CP2K code (a copy is in \texttt{verification/reference\_data/}). Columns
$11$, $12$, $22$ are the elements of the symmetric $2\times2$ matrices
$h^\ell_{ij}$ and $k^\ell_{ij}$. $C_2 = C_3 = C_4 = 0$, there is no third
projector, and the $s$ channel has no $k$ (no spin--orbit coupling for
$\ell = 0$). The reference numbers of Chapter~\ref{ch:verify} were
produced with exactly these digits; rounding them to six decimals changes
total energies in the sixth decimal. A single wrong digit is invisible in
the code and obvious only in the lattice constant, so check your typing
against the source file.}
\label{tab:auhgh}
\end{table}

\begin{keybox}{What the pseudopotential changes in the formalism}
Everywhere ``$\vext$'' appeared in Chapters~\ref{ch:hkks}--\ref{ch:solids},
substitute: a \emph{local} part $\sum_a V_{\mathrm{loc}}(|\rr - \RR_a|)$
(handled in Fourier space with structure factors,
Chapter~\ref{ch:hamiltonian}) plus the \emph{nonlocal} separable operator
\eqref{eq:separable}--\eqref{eq:hghso}. The electron count $N_e$ becomes the
number of \emph{valence} electrons ($11$ or $19$ for Au). The ion--ion energy
$E_{NN}$ uses $Z_{\mathrm{ion}} = N_e$ point charges (Ewald,
Section~\ref{sec:ewald}). All theorems survive unchanged; the KS equations
merely acquire a nonlocal term.
\end{keybox}

\begin{gpubox}{GPU/JAX notes}
Projectors in reciprocal space factorize into radial, angular and
structure-phase parts, Eq.~\eqref{eq:betaG}; build them once per geometry as a dense
array $[N_k, N_{\mathrm{proj}}, N_G^{\max}]$. The nonlocal action on all
bands is then two batched GEMMs:
$P = B^{\dagger} C$ (shape $[N_k, N_{\mathrm{proj}}, N_{\mathrm{bands}}
\times 2]$) followed by $C \mathrel{+}= B\, (H P)$ where $H$ is the small
block-diagonal $h/k$ coupling matrix. Real spherical harmonics up to
$\ell = 2$ ($s,p,d$: 9 functions) are short closed-form polynomials of the
unit vector components --- write them explicitly, no recursion needed
(Table~\ref{tab:realylm}). At $\kk+\GG = 0$ the direction is undefined;
since $\tilde p_{i\ell}(0) = 0$ for $\ell > 0$ any finite value will do
(but make sure it \emph{is} finite: divide by $\max(|\qq|, \epsilon)$, not
by $|\qq|$, or a NaN at $\Gamma$ will poison the whole array).
\end{gpubox}

\begin{table}[h]
\centering
\begin{tabular}{@{}cll@{}}
\toprule
$\ell$ & index $\mu$ & $Y_{\ell\mu}(\hat{\mathbf q})$, with
$\hat{\mathbf q} = (x,y,z)$, $x^2+y^2+z^2=1$\\
\midrule
0 & $0$ & $\sqrt{1/4\pi}$\\
1 & $-1,\ 0,\ 1$ & $\sqrt{3/4\pi}\;\{\,y,\ z,\ x\,\}$\\
2 & $-2,\ -1,\ 0$ & $\sqrt{15/4\pi}\,xy$,\ \ $\sqrt{15/4\pi}\,yz$,\ \
$\sqrt{5/16\pi}\,(3z^2-1)$\\
  & $1,\ 2$ & $\sqrt{15/4\pi}\,xz$,\ \ $\sqrt{15/16\pi}\,(x^2-y^2)$\\
\bottomrule
\end{tabular}
\caption{Real spherical harmonics through $\ell=2$, in the order
$\mu = -\ell\dots\ell$ used throughout. They are orthonormal on the unit
sphere and satisfy $\sum_\mu Y_{\ell\mu}(\hat{\mathbf q})^2 =
(2\ell+1)/4\pi$ --- a one-line unit test.}
\label{tab:realylm}
\end{table}

\chapter{Relativity and Spin--Orbit Coupling}
\label{ch:soc}

Gold is the textbook relativistic solid. The $1s$ electrons of Au move with
$\langle v \rangle / c \sim Z\alpha \approx 0.58$; relativity contracts the
$s$ shells, which screens the nucleus more effectively, which expands and
destabilizes the $d$ shells. The net effect on the valence structure is a
$6s$ band pulled down and a $5d$ band pushed up, narrowing the $5d \to 6s$
interband gap from the UV (as in silver) to $\sim 2.4$\,eV --- absorbing blue
light and reflecting yellow. On top of these \emph{scalar} relativistic
shifts, the spin of the electron couples to its orbital motion and splits the
$d$ levels by $\sim 1.5$\,eV. This chapter derives where these effects come
from and exactly how they enter your Hamiltonian.

\section{From Dirac to Pauli}

The relativistic electron obeys the Dirac equation. For a particle in a
potential $V(\rr)$, the four-component Dirac Hamiltonian is
\begin{equation}
\Ha_D = c\, \bm\alpha \cdot \hat{\mathbf p} + \beta c^2 + V(\rr),
\end{equation}
with the standard $4\times4$ matrices $\bm\alpha, \beta$. Writing the
four-spinor as large and small components $(\varphi, \chi)$ and eliminating
$\chi$ in favor of $\varphi$ (valid when $E - V \ll c^2$, i.e.\ for valence
electrons), a systematic expansion in $1/c^2$ --- the Foldy--Wouthuysen /
Pauli reduction --- yields the \textbf{Pauli Hamiltonian} acting on
two-component spinors:
\begin{equation}
\Ha_{\mathrm{Pauli}} =
\underbrace{\frac{\hat p^2}{2} + V}_{\text{non-relativistic}}
\;\underbrace{-\; \frac{\hat p^4}{8 c^2}}_{\text{mass--velocity}}
\;+\; \underbrace{\frac{1}{8 c^2} \nabla^2 V}_{\text{Darwin}}
\;+\; \underbrace{\frac{1}{4 c^2}\,
\bm\sigma \cdot \big(\nabla V \times \hat{\mathbf p}\big)}_{\text{spin--orbit}} .
\label{eq:pauli}
\end{equation}
Interpretation of the three corrections:
\begin{itemize}[leftmargin=2em]
\item \textbf{Mass--velocity}: the relativistic increase of inertia,
$\sqrt{p^2c^2 + c^4} \approx c^2 + \tfrac{p^2}{2} - \tfrac{p^4}{8c^2}$.
Lowers high-kinetic-energy ($s$) states.
\item \textbf{Darwin}: a contact-like term from the electron's
\emph{Zitterbewegung} smearing it over a Compton wavelength; only $s$ states
(finite at the nucleus) feel the $\nabla^2 V \propto \delta(\rr)$ of a
Coulomb potential strongly.
\item \textbf{Spin--orbit}: in the electron's rest frame, the electric field
$-\nabla V$ of the nucleus appears partly as a magnetic field
$\mathbf B' \approx -(\mathbf v \times \mathbf E)/c$ (in our Gaussian-based
atomic units; in SI units the denominator is $c^2$) that acts on the spin
moment (a factor $\tfrac12$ from Thomas precession is included in
\eqref{eq:pauli}).
\end{itemize}
For a central potential, $\nabla V = (\dd V/\dd r)\, \hat{\mathbf r}$, and
the SOC term becomes the familiar
\begin{equation}
\Ha_{\mathrm{SO}} = \frac{1}{2c^2} \frac{1}{r} \frac{\dd V}{\dd r}\;
\mathbf{L}\cdot\mathbf{S},
\qquad \mathbf S = \tfrac12 \bm\sigma .
\label{eq:hso_central}
\end{equation}
Because $\dd V/\dd r \sim Z/r^2$ near the nucleus, $\Ha_{\mathrm{SO}}$ is
utterly dominated by the core region --- SOC scales roughly like $Z^2$
(for valence $d$ states of the heavy elements, between $Z^2$ and $Z^4$
depending on screening), which is why it is negligible in aluminum, small in
copper, and huge in gold.

The two spin-diagonal corrections (mass--velocity + Darwin) are called
\emph{scalar relativistic}; they shift bands without splitting spins. The SOC
term couples the two spin components: eigenstates are no longer spin-up or
spin-down but two-component \textbf{spinors}, and the good atomic quantum
number becomes the total angular momentum $j = \ell \pm \tfrac12$, since
\begin{equation}
\mathbf L \cdot \mathbf S
= \tfrac12\big(J^2 - L^2 - S^2\big)
= \tfrac12\Big[j(j{+}1) - \ell(\ell{+}1) - \tfrac34\Big]
= \begin{cases}
+\,\ell/2, & j = \ell + \tfrac12,\\[2pt]
-\,(\ell+1)/2, & j = \ell - \tfrac12 .
\end{cases}
\label{eq:ls_eigen}
\end{equation}
For a $d$ shell ($\ell=2$): the $j = 5/2$ level rises by $+\xi$ and
$j = 3/2$ falls by $-\tfrac32 \xi$, where $\xi = \tfrac{1}{2c^2}\braket{
r^{-1}\, \dd V/\dd r}$; the splitting is $\tfrac52 \xi \approx 1.5$\,eV in
the gold \emph{atom}, i.e.\ $\xi \approx 0.6$\,eV.

\paragraph{What the crystal does to this picture.} In the solid the atom
sits in a cubic environment, and $j$ is no longer a good quantum number. At
$\Gamma$ the cubic crystal field alone (no SOC) splits the five $d$
orbitals into a triplet $t_{2g}$ ($xy, yz, xz$; band label
$\Gamma_{25'}$) and, in gold about $1.5$\,eV higher, a doublet $e_g$
($z^2, x^2{-}y^2$; label $\Gamma_{12}$). Switching on SOC then splits the
$t_{2g}$ triplet (six states with spin) into a quartet and a doublet, while
the $e_g$ doublet (four states with spin) stays together. So at $\Gamma$
you should find \textbf{three} $d$ levels with degeneracies
$\mathbf{4, 2, 4}$ (labels $\Gamma_8^+, \Gamma_7^+, \Gamma_8^+$), not the
$4+6$ pattern of the free atom. A ten-line model reproduces the pattern:
diagonalize the $10\times10$ matrix $\Delta\,\hat P_{e_g} + \xi\,\mathbf
L\cdot\mathbf S$ with the matrices of Section~\ref{sec:lsmatrix}, where
$\hat P_{e_g}$ projects on the two $e_g$ orbitals (in the real-harmonic
order of Table~\ref{tab:realylm}, $\mu = 0$ and $\mu = 2$, for both
spins). For $\Delta = 2\xi$ the eigenvalues are exactly $-\xi$ (4-fold),
$+\xi$ (2-fold), $+2.5\,\xi$ (4-fold). For gold the reference calculation
of Chapter~\ref{ch:verify} (LDA) gives levels at $-5.45$, $-4.24$ and
$-3.02$\,eV relative to $E_F$, which the model reproduces to within $0.02$\,eV
with $\xi = 0.62$\,eV and $\Delta = 1.54$\,eV --- the latter being
precisely the $t_{2g}$--$e_g$ splitting found with SOC switched off.
Angle-resolved photoemission places the three levels at $-5.9$ to
$-6.0$, $-4.5$ to $-4.7$ and $-3.6$ to $-3.7$\,eV (three independent
measurements, collected in Table~\ref{tab:dbandexp}): the same pattern,
$0.5$--$0.7$\,eV deeper (a known shortcoming of LDA and GGA for the $d$
bands of noble metals).

\section{Relativity through the pseudopotential}

Here is the beautiful practical point: since all three relativistic terms
live near the nucleus --- exactly the region the pseudopotential replaces ---
\textbf{relativity can be injected into the calculation entirely through the
pseudopotential}, while the valence electrons in the interstitial region are
treated with ordinary non-relativistic kinetic energy $-\tfrac12\nabla^2$.
The construction: solve the \emph{radial Dirac equation} for the isolated
atom (all electrons, fully relativistic), obtaining valence levels
$\varepsilon_{\ell j}$ and wavefunctions labeled by $j = \ell \pm \tfrac12$;
then build norm-conserving pseudopotentials $V_{\ell j}(r)$ per $(\ell, j)$
channel. Regroup the two $j$ channels of each $\ell$ into a
\textbf{scalar-relativistic average} and a \textbf{spin--orbit difference}:
\begin{equation}
V^{\mathrm{SR}}_{\ell}
= \frac{(\ell{+}1)\, V_{\ell, \ell+\frac12} + \ell\, V_{\ell, \ell-\frac12}}
{2\ell + 1},
\qquad
\Delta V^{\mathrm{SO}}_{\ell}
= \frac{2\,\big(V_{\ell,\ell+\frac12} - V_{\ell,\ell-\frac12}\big)}{2\ell+1},
\end{equation}
in terms of which the ionic potential acts on spinors as
\begin{equation}
\hat V_{\mathrm{ion}} = V_{\mathrm{loc}} + \sum_{\ell}
\Big[V^{\mathrm{SR}}_{\ell}
+ \Delta V^{\mathrm{SO}}_{\ell}\; \mathbf L \cdot \mathbf S\Big]
\hat P_{\ell},
\label{eq:vsplit}
\end{equation}
$\hat P_\ell$ projecting on angular momentum $\ell$ around the atom. (Check
the algebra: acting on a $j = \ell\pm\tfrac12$ state and using
\eqref{eq:ls_eigen} recovers $V_{\ell j}$ exactly.) The HGH relativistic
parametrization of Chapter~\ref{ch:pseudo} delivers precisely
\eqref{eq:vsplit} in separable form: the matrices $h^{\ell}_{ij}$ encode
$V^{\mathrm{SR}}_\ell$ and the matrices $k^{\ell}_{ij}$ encode
$\Delta V^{\mathrm{SO}}_\ell$, both with the same analytic projectors
\eqref{eq:hghproj}. Switching SOC off ($k \to 0$) gives you the
scalar-relativistic reference calculation --- a one-line experiment that
isolates the SOC physics in your band plots.

\section{Matrix elements of $\mathbf{L}\cdot\mathbf{S}$}
\label{sec:lsmatrix}

You need $(\mathbf L\cdot\mathbf S)$ in the basis
$\ket{\ell m}\otimes\ket{s}$, $s\in\{\uparrow,\downarrow\}$, used by the
projectors. Write
\begin{equation}
\mathbf L \cdot \mathbf S
= L_z S_z + \tfrac12\big(L_+ S_- + L_- S_+\big),
\qquad
L_\pm \ket{\ell m} = \sqrt{\ell(\ell{+}1) - m(m{\pm}1)}\;\ket{\ell, m{\pm}1}.
\end{equation}
In block form over spin (each block an operator on $m = -\ell..\ell$):
\begin{equation}
\mathbf L\cdot\mathbf S =
\frac12
\begin{pmatrix}
L_z & L_- \\
L_+ & -L_z
\end{pmatrix},
\label{eq:lsblocks}
\end{equation}
a Hermitian $2(2\ell{+}1)\times 2(2\ell{+}1)$ matrix you can hard-code for
$\ell = 1, 2$ (SOC vanishes identically for $\ell = 0$). \emph{Caution}: the
ladder-operator formulas hold in the basis of \emph{complex} spherical
harmonics $Y_{\ell m}$. If your projectors use \emph{real} spherical
harmonics (Table~\ref{tab:realylm}, recommended for everything else),
transform each orbital matrix first: $L^{\mathrm{real}}_{z,\pm} =
U^\dagger L_{z,\pm} U$, then assemble Eq.~\eqref{eq:lsblocks} from the
transformed blocks. $U$ is the $(2\ell{+}1)\times(2\ell{+}1)$ unitary
whose column $\mu$ holds the expansion of the real harmonic $Y_{\ell\mu}$
in complex ones (Condon--Shortley phases):
\begin{equation}
Y_{\ell\mu} =
\begin{cases}
\dfrac{1}{\sqrt2}\Big[Y_{\ell}^{-\mu} + (-1)^{\mu}\, Y_{\ell}^{\mu}\Big],
 & \mu > 0,\\[8pt]
Y_{\ell}^{0}, & \mu = 0,\\[4pt]
\dfrac{\ii}{\sqrt2}\Big[Y_{\ell}^{-|\mu|} - (-1)^{|\mu|}\, Y_{\ell}^{|\mu|}\Big],
 & \mu < 0 .
\end{cases}
\label{eq:realtocomplex}
\end{equation}
In the real basis $L_x, L_y, L_z$ are purely imaginary antisymmetric
matrices --- SOC is where your Hamiltonian becomes irreducibly
complex-valued. Three checks that catch every sign and ordering error:
(i) $U^\dagger U = \mathbf 1$; (ii) for $\ell = 1$ in the order
$(y, z, x)$ the result must be $(L_k)_{ab} = -\ii\,\epsilon_{kab}$, e.g.\
$(L_z)_{xy} = -\ii$; (iii) the eigenvalues of the assembled
$\mathbf L\cdot\mathbf S$ must be those of Eq.~\eqref{eq:ls_eigen}:
$\{+\tfrac12 (\times4), -1 (\times 2)\}$ for $\ell=1$ and
$\{+1 (\times 6), -\tfrac32 (\times 4)\}$ for $\ell = 2$.

The complete nonlocal operator in your code is then
Eq.~\eqref{eq:hghso}: block matrices
$D^{\ell}_{(im s),(j m' s')} = h^{\ell}_{ij}\,\delta_{mm'}\delta_{ss'}
+ k^{\ell}_{ij}\, (\mathbf L\cdot\mathbf S)_{ms,\,m's'}$,
contracted with projections $\braket{\beta_{j\ell m}|\psi^{s'}}$.

\section{Spinor Kohn--Sham equations}

Assembling everything, the object your eigensolver diagonalizes at each
$\kk$-point is the $2\times2$-block operator acting on spinors
$\Psi_{n\kk} = \big(\psi^{\uparrow}_{n\kk}, \psi^{\downarrow}_{n\kk}\big)^{\mathsf T}$:
\begin{equation}
\boxed{\;
\sum_{\beta}\Big[
\Big(-\tfrac12 \nabla^2 + V_{\mathrm{loc}} + \vH + \vxc\Big)
\delta_{\alpha\beta}
+ \mathbf B_{\mathrm{xc}}\cdot \bm\sigma_{\alpha\beta}
+ \big(\hat V_{\mathrm{NL}}\big)_{\alpha\beta}
\Big]\, \psi^{\beta}_{n\kk}
= \varepsilon_{n\kk}\, \psi^{\alpha}_{n\kk},
\;}
\label{eq:spinorks}
\end{equation}
with density matrix and fields rebuilt from the spinors:
\begin{equation}
n_{\alpha\beta}(\rr) = \sum_{n\kk} w_{\kk} f_{n\kk}\,
\psi^{\alpha}_{n\kk}(\rr)\, \psi^{\beta *}_{n\kk}(\rr),
\qquad
n = n_{\uparrow\uparrow} + n_{\downarrow\downarrow},
\quad
\mathbf m = \mathrm{tr}[\bm\sigma\, \hat n].
\end{equation}
Only $\hat V_{\mathrm{NL}}$ (through $k^{\ell}_{ij}\,\mathbf L\cdot\mathbf S$)
and $\mathbf B_{\mathrm{xc}}$ couple the spin components. Note that each
band now holds \emph{one} electron ($f_{n\kk} \le 1$): spin degeneracy is no
longer factored out. Gold with $q{=}11$ needs $\ge 11$ occupied spinor bands
per $\kk$-point; solve for $\sim 20$--$24$.

\section{Symmetry: time reversal, inversion, Kramers}
\label{sec:symmetry}

Two symmetries of nonmagnetic gold organize everything about its spinor
bands. Both turn out to be simple, exact operations on your coefficient
arrays --- which makes them unit tests now, and working tools in
Chapter~\ref{ch:topology}.

\paragraph{Time reversal.} Reversing the direction of time reverses
momenta and spins. For a spin-$\tfrac12$ particle the operator is
$\Theta = \ii \sigma_y K$, with $K$ complex conjugation; on a spinor,
\begin{equation}
\Theta \begin{pmatrix}\psi^{\uparrow}\\ \psi^{\downarrow}\end{pmatrix}
= \begin{pmatrix}\ \ \psi^{\downarrow *}\\ -\psi^{\uparrow *}\end{pmatrix},
\qquad \Theta^2 = -1 .
\end{equation}
Complex conjugation turns $\ee^{\ii(\kk+\GG)\cdot\rr}$ into
$\ee^{\ii(-\kk-\GG)\cdot\rr}$, so $\Theta$ maps a Bloch state at $\kk$ to
one at $-\kk$ whose coefficients are
\begin{equation}
(\Theta c)^{\uparrow}_{-\kk}(\GG) = \big[c^{\downarrow}_{\kk}(-\GG)\big]^*,
\qquad
(\Theta c)^{\downarrow}_{-\kk}(\GG) = -\big[c^{\uparrow}_{\kk}(-\GG)\big]^* .
\label{eq:trcoeff}
\end{equation}
The Hamiltonian commutes with $\Theta$ whenever nothing in it singles out a
direction of spin or of circulation: kinetic energy and real local
potentials obviously qualify, and so does $\mathbf L\cdot\mathbf S$ (both
factors change sign). An exchange field $\mathbf B_{\mathrm{xc}}\cdot\bm\sigma$
does not. For non-magnetic gold, $[\Ha,\Theta]=0$, with two consequences.
\begin{itemize}[leftmargin=2em]
\item \textbf{Kramers theorem.} Every state at $\kk$ has a partner of the
same energy at $-\kk$, namely $\Theta\psi$, with reversed spin:
$\varepsilon_{n,-\kk} = \varepsilon_{n\kk}$. The partner is a genuinely
different state: if $\Theta\psi = \lambda\psi$ then
$\Theta^2\psi = |\lambda|^2\psi \ne -\psi$, a contradiction.
\item The charge density contributed by $\Theta\psi$ equals that of
$\psi$, while its spin density is reversed:
$\mathbf m_{\Theta\psi}(\rr) = -\mathbf m_{\psi}(\rr)$. Summed over the
whole mesh the magnetization contributions of $\kk$ and $-\kk$ therefore
cancel, $\mathbf m = 0$. For $n(\rr)$ the point $-\kk$ adds nothing new,
so you may sample one of each $\pm\kk$ pair with doubled weight --- but
then the cancellation of $\mathbf m$ is lost, and $\mathbf m$ must be set
to zero by hand (see ``Time reversal and the magnetization'' below).
\end{itemize}
At a few special points $-\kk$ and $\kk$ are the \emph{same} point of the
zone, because they differ by a reciprocal lattice vector: $2\kk = \GG_0$.
These are the \textbf{time-reversal-invariant momenta} (TRIM). With
fractional coordinates $\kk = \tfrac12(n_1,n_2,n_3)$, $n_i\in\{0,1\}$,
there are always eight; in FCC they are $\Gamma$, the three $X$ points and
the four $L$ points:
\begin{center}
\begin{tabular}{@{}lll@{}}
\toprule
Point & How many & Fractional coordinates $(n_1,n_2,n_3)/2$\\
\midrule
$\Gamma$ & 1 & $(0,0,0)$\\
$X$ & 3 & $(0,\tfrac12,\tfrac12)$, $(\tfrac12,0,\tfrac12)$,
$(\tfrac12,\tfrac12,0)$\\
$L$ & 4 & $(\tfrac12,0,0)$, $(0,\tfrac12,0)$, $(0,0,\tfrac12)$,
$(\tfrac12,\tfrac12,\tfrac12)$\\
\bottomrule
\end{tabular}
\end{center}
At a TRIM both Kramers partners live at the same $\kk$, so every level
there is at least doubly degenerate, whatever the crystal structure.

\paragraph{Inversion.} The gold crystal is unchanged by
$\rr\to-\rr$ about any atom. The operator $(\hat P\psi)(\rr) = \psi(-\rr)$
acts identically on both spinor components (spin is an angular momentum,
and angular momenta do not change sign under inversion), commutes with
$\Ha$ and with $\Theta$, and --- like $\Theta$ --- sends $\kk\to-\kk$:
\begin{equation}
(\hat P c)^{s}_{-\kk}(\GG) = c^{s}_{\kk}(-\GG) .
\label{eq:pcoeff}
\end{equation}
Hence $\varepsilon_{n,-\kk}=\varepsilon_{n\kk}$ a second time, for an
independent reason, and the combination $\hat P\Theta$ maps every $\kk$ to
\emph{itself}. It is anti-unitary and squares to $-1$, so the Kramers
argument applies at each $\kk$ separately: \textbf{in gold every band is
exactly two-fold degenerate at every} $\kk$. SOC in gold splits orbital
degeneracies (the $d$ manifold) but cannot spin-split individual bands.
Your computed eigenvalues must come in exact pairs --- a stringent, free
correctness test. We call such a pair of bands a \emph{Kramers pair}: pair
1 is spinor bands 1 and 2, pair 2 is bands 3 and 4, and so on. (In
crystals without inversion, e.g.\ GaAs or any surface, SOC does spin-split
the bands away from the TRIM --- Rashba/Dresselhaus physics the same code
captures.)

\paragraph{Time reversal and the magnetization.} A Kramers pair at one
$\kk$ is \emph{not} spin-neutral by itself. Its two members are $\psi$ and
$\hat P\Theta\psi$; $\Theta$ reverses the spin density and $\hat P$ maps
$\rr\to-\rr$, so the pair contributes
\begin{equation}
\mathbf m_{\mathrm{pair},\kk}(\rr) = \mathbf m_\psi(\rr) - \mathbf m_\psi(-\rr),
\label{eq:mpair}
\end{equation}
an odd function of $\rr$ that is in general \emph{not} zero once SOC is on
(the spin texture of an SOC eigenstate has no definite parity). This is
independent of how the eigensolver mixes the two members, because it is a
trace over the pair. The states at $-\kk$, which are $\Theta\psi$ and
$\hat P\psi$, contribute exactly $-\mathbf m_{\mathrm{pair},\kk}(\rr)$, so
$\mathbf m$ vanishes only when both $\kk$ and $-\kk$ are summed. On a mesh
reduced by time reversal (one point of each $\pm\kk$ pair, doubled weight)
the doubled weight reproduces $n(\rr)$ exactly but \emph{doubles}
Eq.~\eqref{eq:mpair} instead of cancelling it. Accumulated and fed to the
noncollinear XC functional, this spurious $\mathbf m$ produces a spurious
$\mathbf B_{\mathrm{xc}}$, which breaks time reversal and splits the
Kramers pairs. The rule is therefore: on a time-reversal-reduced mesh,
accumulate $n$ only and set $\mathbf m = 0$ (equivalently, symmetrize the
density matrix, $\hat n \to \tfrac12\big(\hat n + \sigma_y \hat n^{*}
\sigma_y\big)$, which removes $\mathbf m$ and keeps $n$); a magnetic
calculation must use the full $\pm\kk$ mesh. (Without SOC each state is a
spatial orbital times a fixed spinor, and inversion combined with complex
conjugation gives $|\phi_\kk(-\rr)|^2 = |\phi_\kk(\rr)|^2$, so
Eq.~\eqref{eq:mpair} vanishes identically --- the problem is specific to
SOC.)

\paragraph{Both operations at a TRIM: permutations of your array.} At a
TRIM the state at $-\kk$ is a state at $\kk$ with relabeled plane waves:
$-\kk+\GG' = \kk + (\GG'-\GG_0)$. Applying this relabeling to
Eqs.~\eqref{eq:trcoeff} and \eqref{eq:pcoeff} gives operators that act
\emph{within} the coefficient array of one $\kk$-point:
\begin{equation}
\boxed{\;
\begin{aligned}
(\hat P c)^{s}(\GG) &= c^{s}(-\GG-\GG_0),\\
(\Theta c)^{\uparrow}(\GG) &= \big[c^{\downarrow}(-\GG-\GG_0)\big]^*,\qquad
(\Theta c)^{\downarrow}(\GG) = -\big[c^{\uparrow}(-\GG-\GG_0)\big]^*,
\end{aligned}
\qquad \GG_0 = 2\kk .\;}
\label{eq:parityperm}
\end{equation}
The cutoff sphere is mapped onto itself because
$|\kk + (-\GG-\GG_0)| = |\kk+\GG|$. In integer (Miller-index) form with
$\kk = \tfrac12(n_1,n_2,n_3)$, the plane wave $(m_1,m_2,m_3)$ is sent to
$(-m_1-n_1,\,-m_2-n_2,\,-m_3-n_3)$: both operators are one fixed
permutation of the $\GG$ index (plus, for $\Theta$, a spin swap and a
conjugation). Three unit tests follow, each for a random spinor $c$ at
each of the eight TRIM:
\begin{equation}
\hat H(\hat P c) = \hat P(\hat H c),
\qquad
\hat H(\Theta c) = \Theta(\hat H c),
\qquad
\Theta(\Theta c) = -c ,
\end{equation}
to rounding error. The first fails if the projector phases
$(-\ii)^\ell Y_{\ell m}$ or the structure factors are wrong (and it
requires the atom at the origin, as we have it). The second fails if any
term breaks time reversal --- a stray $\mathbf B_{\mathrm{xc}}$, or an
$\mathbf L\cdot\mathbf S$ block whose orbital matrices are not purely
imaginary in the real-harmonic basis. (Neither test can detect a
\emph{consistent} relabeling of the harmonics or a sign change of the
spin-flip blocks, which are unitary transformations of a correct
Hamiltonian; the eigenvalue checks of Section~\ref{sec:lsmatrix} and the
degeneracy pattern at $\Gamma$ cover those.)

\paragraph{Parity as a band label.} Because $\hat P$ maps a TRIM to itself
and commutes with $\Ha$, the eigenstates there can be chosen as
eigenstates of $\hat P$, and since $\hat P^2 = 1$ the eigenvalue is
$\xi = +1$ (\emph{even}) or $\xi=-1$ (\emph{odd}). Since
$[\hat P,\Theta]=0$, both members of a Kramers pair have the same $\xi$:
\textbf{each Kramers pair at each TRIM carries one sign.} This is the
solid-state version of the parity you know from atoms, and the connection
is direct: a Bloch sum $\sum_{\RR}\ee^{\ii\kk\cdot\RR}\phi_\ell(\rr-\RR)$
of atomic orbitals of angular momentum $\ell$ has, at any TRIM, parity
$(-1)^\ell$ --- $s$ and $d$ bands are even, $p$ bands are odd. It is also
exactly the even/odd distinction between the $\cos$ and $\sin$ states of
the two-plane-wave model, Eq.~\eqref{eq:nfeparity}. The superscripts in
the group-theory labels used above are these parities: $\Gamma_8^{+}$ and
$\Gamma_7^{+}$ are even, as $d$ states must be. Computing $\xi$ takes
three lines: for a non-degenerate pair,
$\xi = \braket{\psi|\hat P|\psi} = \sum_{\GG s} c^{s*}(\GG)\,
c^{s}(-\GG-\GG_0)$; for a block of bands, the eigenvalues of the matrix
$\braket{\psi_m|\hat P|\psi_n}$. Chapter~\ref{ch:topology} multiplies
these signs together.

\begin{keybox}{What to expect in the gold band structure}
Without SOC (set $k^\ell = 0$): five scalar-relativistic $d$ bands
$\sim 1.7$--$7.5$\,eV below $E_F$, crossed by the broad $6s$ band. With SOC on:
at $\Gamma$ the $d$ manifold forms three levels with degeneracies $4, 2, 4$
(the $t_{2g}$-derived quartet and doublet, split by roughly
$1.2$--$1.5$\,eV, and the $e_g$-derived quartet above); degeneracies along $\Gamma$--$X$ and
$\Gamma$--$L$ are lifted; and all bands remain exactly doubly degenerate.
Compare against any published relativistic Au band structure --- the
agreement, from a code you wrote yourself, is the payoff of this whole
project. Figure~\ref{fig:bands} shows what it should look like.
\end{keybox}

\begin{figure}[H]
\centering
\pgfplotsset{bandplot/.style={width=0.52\textwidth, height=0.62\textwidth,
  xmin=0, xmax=3.37024, ymin=-11, ymax=8,
  xtick={0,0.81529,1.22294,1.79943,2.5055,3.37024},
  xticklabels={$\Gamma$,$X$,$W$,$L$,$\Gamma$,$K$},
  xmajorgrids=true, grid style={black!35},
  tick label style={font=\small}, label style={font=\small},
  title style={font=\small}, no markers}}
\begin{tikzpicture}
\begin{axis}[bandplot, ylabel={$\varepsilon - E_F$ (eV)}, title={SOC off (scalar relativistic)}]
\draw[black!50, dashed] (axis cs:0,0) -- (axis cs:3.37024,0);
\foreach \i in {1,...,9}{\addplot[blue!70!black, thick] table[x index=0, y index=\i] {figdata/bands_sr.dat};}
\end{axis}
\end{tikzpicture}\hfill
\begin{tikzpicture}
\begin{axis}[bandplot, yticklabels={}, title={SOC on (each line is a Kramers pair)}]
\draw[black!50, dashed] (axis cs:0,0) -- (axis cs:3.37024,0);
\foreach \i in {1,...,8}{\addplot[red!70!black, thick] table[x index=0, y index=\i] {figdata/bands_soc.dat};}
\end{axis}
\end{tikzpicture}
\caption{Band structure of FCC gold computed with the reference
implementation of Chapter~\ref{ch:verify} (HGH $q{=}11$ pseudopotential,
LDA, experimental lattice constant, $E_{\mathrm{cut}}=40$\,Ha). Left:
$k^\ell = 0$. Right: full spin--orbit coupling. The flat $d$ bands lie
between $-7.6$ and $-1$\,eV, crossed by the steep $sp$ band, which is the
only band cutting the Fermi level. SOC lifts degeneracies throughout the
$d$ complex --- at $\Gamma$ the threefold level at $-4.8$\,eV splits into
$-5.45$ (quartet) and $-4.24$\,eV (doublet) --- and separates the sixth
and seventh pairs where they nearly touch between $W$ and $L$. Your code
should reproduce this figure.}
\label{fig:bands}
\end{figure}

\begin{gpubox}{GPU/JAX notes}
The spinor dimension is a fixed axis of length 2. For the Hamiltonian
application and the eigensolver, fold it into the coefficient axis (treat a
spinor wavefunction as one vector of $2N_G$ complex coefficients, so Gram
and Rayleigh--Ritz matrices are ordinary GEMMs); keep it separate for the
density construction. Everything stays
\texttt{complex64/128}; SOC makes the Hamiltonian genuinely complex, so use
complex-Hermitian (not real-symmetric) linear algebra throughout
(\texttt{eigh} accepts complex Hermitian matrices
natively; remember that every transpose in a formula becomes a conjugate
transpose). The $\mathbf L\cdot\mathbf S$ blocks are tiny constant matrices:
precompute $D^{\ell}$ once and fuse into the nonlocal GEMM. Kramers pairing
of eigenvalues is a cheap on-device assertion (sorted eigenvalues should
match in pairs to the eigensolver tolerance). The parity and time-reversal
operations of Eq.~\eqref{eq:parityperm} are a single precomputed gather
index per TRIM.
\end{gpubox}

\chapter{The Hamiltonian in the Plane-Wave Basis}
\label{ch:hamiltonian}

This chapter assembles, term by term, exactly what your program computes.
Throughout, a spinor state at fixed $\kk$ is the coefficient array
$c^{s}(\GG)$, $s \in \{\uparrow,\downarrow\}$, over the sphere
$\tfrac12|\kk+\GG|^2 \le E_{\mathrm{cut}}$.

\section{Representations and the FFT bridge}

Two representations of $u_{n\kk}$:
\begin{itemize}[leftmargin=2em]
\item \textbf{Reciprocal space}: coefficients $c(\GG)$ on the cutoff sphere
(compact; kinetic and nonlocal terms are cheap here).
\item \textbf{Real space}: values $u(\rr_i)$ on a uniform $N_1{\times}N_2
{\times}N_3$ grid spanning the cell (local potential is cheap here).
\end{itemize}
The bridge is the discrete Fourier transform: zero-pad $c(\GG)$ from the
sphere into the full FFT box, inverse-FFT to get $u(\rr_i)$, and conversely
FFT then gather the sphere. How fine must the grid be? Let $G_{\max} = \sqrt{2E_{\mathrm{cut}}}$. The
density $|u|^2$ contains wavevectors up to $2G_{\max}$, so the grid must
represent all $|\GG| \le 2G_{\max}$:
$N_i \ge \tfrac{2}{\pi}\, G_{\max}\, |\mathbf a_i| + 1$,
rounded up to an FFT-friendly size ($2^a 3^b 5^c 7^d$). The same grid is
enough for the product $\veff\, u$: the product does contain wavevectors up
to $3G_{\max}$, but those alias back only onto $|\GG| \ge G_{\max}$,
outside the sphere you keep. For gold ($|\mathbf a_i| = a/\sqrt2 =
5.45$\,Bohr) at $E_{\mathrm{cut}} = 40$\,Ha this gives $N_i \ge 33$ (use
$35$ or $36$); at $60$\,Ha, $N_i = 40$. All conversions
follow the conventions of Eq.~\eqref{eq:fourierconv}.

\section{Term 1: kinetic energy --- diagonal}

\begin{equation}
\braket{\kk{+}\GG | -\tfrac12\nabla^2 | \kk{+}\GG'}
= \tfrac12 |\kk + \GG|^2\, \delta_{\GG\GG'}
\qquad\Longrightarrow\qquad
(\hat T c)(\GG) = \tfrac12 |\kk+\GG|^2\; c(\GG),
\end{equation}
an elementwise multiply, identical on both spin components.

\section{Term 2: local potential --- FFT convolution}

The total local potential
\begin{equation}
V_{\mathrm{eff,loc}}(\rr) = V_{\mathrm{loc,ion}}(\rr) + \vH(\rr) + \vxc(\rr)
\end{equation}
is multiplicative in real space; in reciprocal space it is a convolution.
Never build the convolution matrix --- apply it as
\begin{equation}
c(\GG) \;\xrightarrow{\;\text{pad + IFFT}\;}\; u(\rr_i)
\;\xrightarrow{\;\times V_{\mathrm{eff,loc}}(\rr_i)\;}\;
\;\xrightarrow{\;\text{FFT + gather}\;}\; (\hat V c)(\GG),
\end{equation}
costing $O(N \log N)$ per band. This FFT sandwich dominates the runtime of a plane-wave code and batches
perfectly over bands, spins, and $\kk$-points
(Algorithm~\ref{alg:applyh}).

\paragraph{Ionic local part via structure factor.} Being a lattice sum of
identical atomic contributions, its Fourier components factorize:
\begin{equation}
V_{\mathrm{loc,ion}}(\GG) = \frac{1}{\Omega}\, S(\GG)\;
\tilde V_{\mathrm{loc}}(|\GG|),
\qquad
S(\GG) = \sum_{a} \ee^{-\ii \GG \cdot \RR_a},
\label{eq:structfact}
\end{equation}
with $\tilde V_{\mathrm{loc}}$ the analytic HGH form
\eqref{eq:hghlocG}--\eqref{eq:hghlocpoly}. For
one atom at the origin, $S(\GG) = 1$. Build it once per geometry on the FFT
box (all $\GG$ of the box, with the $\GG=0$ entry set by rule~2 below) and
inverse-FFT to the real-space grid: $V_{\mathrm{loc,ion}}(\rr_j) =
N_{\mathrm{grid}}\cdot\mathrm{IFFT}[V_{\mathrm{loc,ion}}(\GG)]$, which must
come out real to rounding error.

\section{Term 3: Hartree potential --- Poisson in one line}

The Poisson equation $\nabla^2 \vH = -4\pi n$ diagonalizes in plane waves:
\begin{equation}
\boxed{\;
\vH(\GG) = \frac{4\pi\, n(\GG)}{|\GG|^2}, \qquad \GG \ne 0,
\;}
\qquad
E_{\mathrm H} = \frac{\Omega}{2} \sum_{\GG \ne 0}
\frac{4\pi\, |n(\GG)|^2}{|\GG|^2}.
\label{eq:hartreeG}
\end{equation}
One FFT of the density, one elementwise multiply, one inverse FFT. This
``free'' Poisson solver is the second great gift of periodic boundary
conditions.

\paragraph{The $\GG = 0$ problem.}
\label{sec:g0}
Three quantities diverge individually at $\GG = 0$: the Hartree potential of
the electrons ($+4\pi N_e / \Omega G^2$), the local pseudopotential of the
ions ($-4\pi Z_{\mathrm{ion}} N_{\mathrm{at}} / \Omega G^2$), and the
ion--ion interaction. In a neutral cell ($N_e = Z_{\mathrm{ion}}
N_{\mathrm{at}}$) the divergences cancel exactly. What survives is finite, and the recipe is three rules:
\begin{enumerate}[leftmargin=2em]
\item $\vH(\GG{=}0) = 0$.
\item $V_{\mathrm{loc,ion}}(\GG{=}0) = \frac{1}{\Omega}\sum_a \alpha_a$,
where $\alpha_a$ is what is left of $\tilde V_{\mathrm{loc}}(G\to0)$ after
the $-4\pi Z_{\mathrm{ion}}/G^2$ divergence is removed,
\begin{equation}
\alpha = \int \Big[V_{\mathrm{loc}}(r) + \frac{Z_{\mathrm{ion}}}{r}\Big]
\dd\rr
= 2\pi Z_{\mathrm{ion}}\, r_{\mathrm{loc}}^2
+ (2\pi)^{3/2} r_{\mathrm{loc}}^3\,
\big[C_1 + 3C_2 + 15C_3 + 105C_4\big].
\label{eq:alphaZ}
\end{equation}
\item The ion--ion energy is the Ewald energy of point charges in a uniform
neutralizing background (Section~\ref{sec:ewald}).
\end{enumerate}
With rule 2 the constant $\alpha/\Omega$ is simply part of the potential:
it shifts every eigenvalue by the same amount and contributes
$\frac{N_e}{\Omega}\sum_a\alpha_a$ to the energy through the ordinary
local-potential term --- no separate bookkeeping. (Some texts instead set
$V_{\mathrm{loc,ion}}(0) = 0$ and add this energy by hand; either is fine,
but never do both.) These constants shift total energies, not eigenvalue
differences; because they depend on $\Omega$, a mistake shows up as a wrong
lattice constant.

\section{Term 4: nonlocal pseudopotential --- batched projections}

From Eqs.~\eqref{eq:separable}, \eqref{eq:hghso}, with the reciprocal-space
projectors $\tilde\beta_{i\ell m}(\kk{+}\GG)$ of Eq.~\eqref{eq:betaG}
(built from the Bessel transforms $\tilde p_{i\ell}$ of
Eq.~\eqref{eq:hghprojG} and the real harmonics of
Table~\ref{tab:realylm}), the action on a spinor is
\begin{equation}
(\hat V_{\mathrm{NL}}\, c)^{s}(\GG)
= \sum_{a}\sum_{i\ell m}
\tilde\beta_{i\ell m}(\kk{+}\GG)
\sum_{j m' s'} D^{\ell}_{(i m s),(j m' s')}
\underbrace{\sum_{\GG'} \tilde\beta^{*}_{j\ell m'}(\kk{+}\GG')\,
c^{s'}(\GG')}_{\text{projection } P_{j\ell m' s'}} .
\label{eq:vnl_action}
\end{equation}
Two GEMMs around a tiny block-diagonal matrix. For the gold potentials of
Table~\ref{tab:auhgh} ($s,p,d$ channels, two radial projectors each) there
are $2\,(1 + 3 + 5) = 18$ projector functions, hence a $36\times36$ matrix
$D$ once spin is included --- microscopic compared to $N_{\mathrm{pw}}$.
The nonlocal energy is $E_{\mathrm{NL}} = \sum_{n\kk} w_\kk f_{n\kk}\,
P_{n\kk}^\dagger D\, P_{n\kk}$ with $P$ the vector of projections.

\section{Ion--ion energy: Ewald summation}
\label{sec:ewald}

The electrostatic energy of point charges $Z_{\mathrm{ion}}$ on a lattice
(with neutralizing background) converges conditionally in both real and
reciprocal space. Ewald's trick: add and subtract Gaussian charge clouds of
width $\eta^{-1/2}$ on each ion, splitting the sum into two rapidly
convergent pieces,
\begin{equation}
\begin{aligned}
E_{\mathrm{Ewald}}
={}& \frac{2\pi}{\Omega} \sum_{\GG\neq 0}
\frac{\ee^{-G^2/4\eta}}{G^2}\, \big|S_Z(\GG)\big|^2
+ \frac12 \sum_{a b}{\sum_{\RR}}'
Z_a Z_b\, \frac{\mathrm{erfc}\big(\sqrt{\eta}\,
|\RR + \RR_a - \RR_b|\big)}{|\RR + \RR_a - \RR_b|}\\
&- \sqrt{\frac{\eta}{\pi}} \sum_a Z_a^2
- \frac{\pi\, \big(\sum_a Z_a\big)^2}{2\,\eta\,\Omega},
\end{aligned}
\label{eq:ewald}
\end{equation}
where $S_Z(\GG) = \sum_a Z_a \ee^{-\ii\GG\cdot\RR_a}$ and the primed sum
omits the self term $\RR = 0, a = b$. The result is independent of the
splitting parameter $\eta$ --- vary $\eta$ by a factor of 4 and check
invariance to $10^{-10}$; it is the classic unit test. In practice take
$\eta \approx 0.5$\,Bohr$^{-2}$ for the gold cell and truncate where the
terms fall below rounding error: lattice vectors with
$|\RR| \le 6/\sqrt\eta$ (where $\mathrm{erfc} < 10^{-16}$) and reciprocal
vectors with $|\GG| \le 12\sqrt\eta$ (where
$\ee^{-G^2/4\eta} < 10^{-15}$). For one atom per cell
this is a few thousand terms, computed once per geometry.

\section{The assembled total (free) energy}

The most transparent form adds up the pieces one by one (the ``direct''
form, Eq.~\eqref{eq:etot1}):
\begin{equation}
\boxed{
\begin{aligned}
F ={}& \underbrace{\sum_{n\kk} w_\kk f_{n\kk} \sum_{\GG s}
\tfrac12|\kk{+}\GG|^2\,|c^s_{n\kk}(\GG)|^2}_{T_s}
+ \underbrace{\int V_{\mathrm{loc,ion}}\, n\,\dd\rr}_{E_{\mathrm{loc}}}
+ E_{\mathrm{NL}}\\
&+ E_{\mathrm H}[n] + \Exc[n,\mathbf m]
+ E_{\mathrm{Ewald}}
- \sigma\sum_{n\kk} w_\kk\, s_{n\kk} .
\end{aligned}}
\label{eq:fdirect}
\end{equation}
Equivalently, using the eigenvalues (Eq.~\eqref{eq:etot2}; valid when the
$\alpha/\Omega$ constant is inside the Hamiltonian, rule~2 of
Section~\ref{sec:g0}). The eigenvalue sum already contains $T_s$,
$E_{\mathrm{loc}}$ and $E_{\mathrm{NL}}$ exactly once, so only the Hartree
and XC terms need double-counting corrections:
\begin{equation}
\boxed{
\begin{aligned}
F ={}& \sum_{n\kk} w_{\kk}\, f_{n\kk}\, \varepsilon_{n\kk}
- \sigma \sum_{n\kk} w_{\kk}\, s(f_{n\kk})
\\
&- E_{\mathrm H}[n]
- \int \Big[\vxc\, n + \mathbf B_{\mathrm{xc}}\cdot\mathbf m\Big] \dd\rr
+ \Exc[n, \mathbf m]
+ E_{\mathrm{Ewald}} .
\end{aligned}}
\label{eq:ftot}
\end{equation}
Real-space integrals use the grid ($\int \to \frac{\Omega}{N_{\mathrm{grid}}}
\sum_i$). One subtlety: in Eq.~\eqref{eq:ftot} the eigenvalues were computed
in the potential of the \emph{input} density, so $E_{\mathrm H}$,
$\int\vxc n$ and $\Exc$ there must also use the input density, whereas
Eq.~\eqref{eq:fdirect} is naturally evaluated with the output density
(the one built from the current wavefunctions and occupations).
The two numbers differ during the SCF iterations and coincide at
convergence; their difference shrinks quadratically with the density error
and is a superb convergence diagnostic. (Do not expect $F$ to decrease
monotonically from one iteration to the next: density mixing is not a
minimization, and small oscillations early on are normal.)

\paragraph{Equilibrium lattice constant and bulk modulus.} Compute
$F(\Omega)$ at 7--9 volumes within $\pm 8\%$ of experiment, fit the
Birch--Murnaghan equation of state
\begin{equation}
E(V) = E_0 + \frac{9 V_0 B_0}{16}
\left\{ \Big[\big(\tfrac{V_0}{V}\big)^{2/3} - 1\Big]^3 B_0'
+ \Big[\big(\tfrac{V_0}{V}\big)^{2/3} - 1\Big]^2
\Big[6 - 4\big(\tfrac{V_0}{V}\big)^{2/3}\Big] \right\},
\end{equation}
extracting $V_0$ (hence $a_0 = (4V_0)^{1/3}$ for FCC), $B_0$, $B_0'$ by
nonlinear least squares (a parabola through the lowest three points gives
the starting values). In atomic units $B_0$ comes out in Ha/Bohr$^3$;
$1$\,Ha/Bohr$^3 = 29\,421.0$\,GPa. Targets in
Chapter~\ref{ch:verify}. One caveat: at fixed $E_{\mathrm{cut}}$ the number
of plane waves changes in jumps as the cell is strained, which makes $F(V)$
slightly ragged and can bias the fitted $a_0$ and $B_0$ if the cutoff is too
low. Converge the lattice constant with respect to $E_{\mathrm{cut}}$, not
just the energy at one volume.

\begin{gpubox}{GPU/JAX notes}
The Hamiltonian is never materialized: it exists only as the function
$c \mapsto \hat H c$ composed of (elementwise) $+$ (FFT sandwich) $+$ (two
GEMMs). All three stages batch over the composite index (bands $\times$ spin
$\times$ $\kk$): a single call applies $\hat H$ to \emph{every state in the
calculation at once} --- shape
$[N_k, N_b, 2, N_G^{\max}] \to$ same. Zero-pad/gather between sphere and FFT
box with precomputed index arrays and \texttt{scatter/gather} (\texttt{.at[].add()}, see Section~\ref{sec:jaxprimer});
masks handle the padded region. Ewald and structure factors are geometry-time
constants --- compute in float64 on host or device, once.
\end{gpubox}

\chapter{The Algorithms}
\label{ch:algorithms}

Everything so far defined \emph{what} to compute. This chapter gives the
numerical engines: iterative eigensolvers for the linear problem at fixed
potential, and mixing schemes for the nonlinear self-consistency, with full
pseudocode. Notation: at each $\kk$, $H_{\kk}$ is the (matrix-free)
Hamiltonian on spinor coefficient vectors of length $2N_G$; we seek the
lowest $N_b$ eigenpairs.

\section{Setting up a $\kk$-point: the basis}
\label{sec:setup}

Everything at a given $\kk$ is expressed on its own plane-wave sphere.
Building it is pure integer bookkeeping, done once per geometry.

\begin{algorithm}[H]
\caption{Plane-wave basis and projectors at one $\kk$-point}
\label{alg:basis}
\begin{algorithmic}[1]
\Function{BuildBasis}{$\kk_{\mathrm{frac}}$, $E_{\mathrm{cut}}$}
\State $M_i \gets \big\lceil \sqrt{2E_{\mathrm{cut}}}\,|\mathbf a_i|/2\pi
\big\rceil + 1$ \Comment{no $\GG$ in the sphere has $|m_i| > M_i$}
\State \textbf{for all} integer triples $\mathbf m$ with $|m_i|\le M_i$:
$\qq \gets (\kk_{\mathrm{frac}} + \mathbf m)\,B$;\quad keep $\mathbf m$ if
$\tfrac12|\qq|^2 \le E_{\mathrm{cut}}$
\State store for each kept plane wave: $\mathbf m$;\ $\qq$;\ $T_{\qq} =
\tfrac12|\qq|^2$;\ FFT-box index
$\mathrm{idx} = \mathrm{ravel}(m_1 \bmod N_1,\, m_2 \bmod N_2,\, m_3 \bmod N_3)$
\State projectors $\tilde\beta^{a}_{i\ell m}(\qq)$ from
Eqs.~\eqref{eq:betaG}, \eqref{eq:hghprojG}, Table~\ref{tab:realylm}
\Comment{array $[N_{\mathrm{proj}}, N_G]$}
\State \Return basis record (padded to $N_G^{\max}$ with a mask when
batching over $\kk$)
\EndFunction
\end{algorithmic}
\end{algorithm}

Here $B$ is the matrix whose rows are $\mathbf b_i$. The ``$\bmod N_i$''
places negative frequencies where the FFT library expects them (wrapped
order). Sort the kept plane waves in any fixed order you like; nothing
depends on it. Two checks: the count should be close to
$\Omega(2E_{\mathrm{cut}})^{3/2}/6\pi^2$, and at $\Gamma$ the set must be
symmetric under $\mathbf m\to-\mathbf m$.

\section{Applying the Hamiltonian}
\label{sec:applyh}

The eigensolver below never sees a matrix, only this function. It is
Chapter~\ref{ch:hamiltonian} condensed into one routine; $C$ holds any
number of spinor vectors (bands) at one $\kk$.

\begin{algorithm}[H]
\caption{Matrix-free Hamiltonian application at one $\kk$-point}
\label{alg:applyh}
\begin{algorithmic}[1]
\Function{ApplyH}{$C[n, s, \GG]$;\ basis;\ $V_{\mathrm{eff,loc}}(\rr)$,
$\mathbf B_{\mathrm{xc}}(\rr)$ on the grid}
\State $Y \gets T_{\qq}\, C$ \Comment{kinetic: elementwise, same for both $s$}
\State $u[n,s,\rr] \gets \mathrm{IFFT}\big(\text{scatter } C[n,s,\cdot]
\text{ into a zero box at idx}\big)$ \Comment{to real space}
\State $v[n,s,\rr] \gets V_{\mathrm{eff,loc}}(\rr)\,u[n,s,\rr]
+ \sum_{s'}\big(\mathbf B_{\mathrm{xc}}(\rr)\cdot\bm\sigma\big)_{ss'}\,
u[n,s',\rr]$ \Comment{second term absent for gold}
\State $Y \mathrel{+}= \text{gather } \mathrm{FFT}(v) \text{ at idx}$
\Comment{back to the sphere; discard everything outside it}
\State $P[n,s,p] \gets \sum_{\GG} \tilde\beta_p^{*}(\GG)\, C[n,s,\GG]$
\Comment{projections, $p = (a,i,\ell,m)$}
\State $Q[n,s,p] \gets \sum_{s'p'} D_{(p s),(p' s')}\, P[n,s',p']$
\Comment{coupling matrix, Section~\ref{sec:lsmatrix}}
\State $Y \mathrel{+}= \sum_p \tilde\beta_p(\GG)\, Q[n,s,p]$
\State \Return $Y$
\EndFunction
\end{algorithmic}
\end{algorithm}

The FFT normalizations cancel between lines 3 and 5 (use the library's
unnormalized forward and $1/N$-normalized inverse transform and forget
about them here). Test it before anything else: for random vectors
$\phi,\psi$ the two numbers $\braket{\phi|\hat H\psi}$ and
$\braket{\psi|\hat H\phi}^*$ must agree to rounding error, and for a small
cutoff the matrix assembled by applying \textsc{ApplyH} to unit vectors
must have the same eigenvalues as the iterative solver returns.

\section{Why iterative eigensolvers}

$H_{\kk}$ has dimension $2N_G \sim 5\times10^3$; dense diagonalization
($O(N^3) \sim 10^{11}$ flops \emph{per $\kk$-point per SCF step}) is possible
but wasteful, and becomes impossible for larger cells you will want later. We
need only the lowest $N_b \sim 24$ of $\sim 5\times10^3$ eigenpairs of a Hermitian
operator we can apply fast. That is the textbook setting for
\emph{preconditioned block Krylov / subspace iteration} methods. Two
GPU-relevant choices:

\begin{itemize}[leftmargin=2em]
\item \textbf{Blocked Davidson}: expands a subspace with preconditioned
residuals; robust, standard in plane-wave codes.
\item \textbf{LOBPCG} (Locally Optimal Block Preconditioned Conjugate
Gradient): keeps a fixed-size subspace $[X, W, P]$ (current block,
preconditioned residuals, previous search directions) and Rayleigh--Ritz
minimizes within it. Fixed memory, fixed shapes, everything expressed as
dense GEMMs on tall matrices --- \emph{the} eigensolver for JAX/GPU. We make
it the primary recommendation.
\end{itemize}

\section{The preconditioner}

Iterative methods converge at a rate set by the spread of the spectrum. A
plane-wave Hamiltonian's spectrum is dominated by kinetic energy
$\tfrac12|\kk+\GG|^2$, which ranges up to $E_{\mathrm{cut}}$; but we only
care about the lowest few eV. The cure is a diagonal preconditioner that
approximately inverts the kinetic term at high $\GG$ while leaving low $\GG$
alone. The classic Teter--Payne--Allan form, applied elementwise to a
residual $r(\GG)$ with $x = \tfrac12|\kk+\GG|^2 / E_{\mathrm{kin}}(\psi)$:
\begin{equation}
K(x) = \frac{27 + 18x + 12x^2 + 8x^3}
{27 + 18x + 12x^2 + 8x^3 + 16 x^4},
\label{eq:teter}
\end{equation}
smoothly interpolating between $1$ (small $x$) and $1/(2x)$ (large $x$);
$E_{\mathrm{kin}}(\psi) = \sum_{\GG s}\tfrac12|\kk{+}\GG|^2 |c^s(\GG)|^2$ is
the kinetic energy of the band being corrected.
A simpler alternative, $K = 1/\max\big(1, \tfrac12|\kk+\GG|^2 - \tau\big)$
with shift $\tau \approx$ current Ritz value, also works. The preconditioner
is elementwise and free; it typically cuts iteration counts by
$5$--$10\times$.

\section{LOBPCG}

The idea in one paragraph. Suppose $X$ holds your current guesses for the
lowest $N_b$ eigenvectors. The residual $R = HX - X\Lambda$ tells you in
which direction each guess is wrong. Rather than taking a step of some
guessed length along $R$, build the small subspace spanned by the current
vectors $X$, the (preconditioned) residuals $W$, and the previous step's
update directions $P$, and solve the eigenproblem \emph{exactly inside that
subspace} --- a $3N_b\times3N_b$ dense diagonalization, called the
\textbf{Rayleigh--Ritz} step. The best $N_b$ vectors of the subspace
become the new $X$. Because the subspace contains the old $X$, the Ritz
values can only go down; because it contains $P$, the method has the
memory that makes conjugate gradients fast.

Written naively, with $S = [X, W, P]$ taken as it comes, the Rayleigh--Ritz
step is a \emph{generalized} eigenproblem $S^\dagger H S\, y = \theta\,
S^\dagger S\, y$, and it is the weak point of the method: near convergence
$W$ and $P$ shrink and become almost parallel to each other, the overlap
matrix $S^\dagger S$ becomes numerically singular, and the iteration falls
apart (details below). The cure is to keep the three blocks
\emph{orthonormal and mutually orthogonal at all times}, so that
$S^\dagger S = \mathbf 1$ and the Rayleigh--Ritz step is an ordinary
Hermitian eigenproblem. One small utility does all the orthonormalizing.

\begin{algorithm}[H]
\caption{Orthonormalize the columns of $V$, discarding (zeroing) linearly
dependent directions; optionally carry $HV$ along}
\label{alg:ortho}
\begin{algorithmic}[1]
\Function{Orthonormalize}{$V$, $\tau$ [, $HV$]}
\State $d_j \gets \|V_{:,j}\|$ (set $d_j = 1$ where the column is zero)
\Comment{column scales}
\State $G \gets \mathrm{diag}(d)^{-1}\, V^\dagger V\, \mathrm{diag}(d)^{-1}$
\Comment{small Hermitian matrix with unit diagonal; double precision}
\State $(\sigma, U) \gets \mathrm{eigh}(G)$
\Comment{$\sigma_j \ge 0$ measures how independent direction $j$ is}
\State $T \gets \mathrm{diag}(d)^{-1}\; U\;
\mathrm{diag}\big(\sigma_j > \tau\ ?\ \sigma_j^{-1/2} : 0\big)$
\State \Return $V\,T$ [and $HV\,T$]
\Comment{columns are orthonormal, or exactly zero}
\EndFunction
\end{algorithmic}
\end{algorithm}

Because the same matrix $T$ multiplies $V$ and $HV$, the pair stays
consistent ($HV$ remains the image of $V$ under $\hat H$) without a new
Hamiltonian application. Dependent directions are not removed --- that
would change array shapes --- but replaced by zero columns, which the
eigensolver then ignores. There is no failure branch, unlike a Cholesky
factorization, which simply breaks down on a singular matrix.

\begin{algorithm}[H]
\caption{LOBPCG for the lowest $N_b$ spinor bands at one $\kk$
(all $\kk$ run batched)}
\label{alg:lobpcg}
\begin{algorithmic}[1]
\State $X \gets$ initial block $[2N_G \times N_b]$ (previous SCF
step's eigenvectors; first step: random complex numbers divided by
$1 + \tfrac12|\kk{+}\GG|^2$)
\State $X \gets$ \Call{Orthonormalize}{$X$, $10^{-12}$}, twice;\quad
$HX \gets$ \Call{ApplyH}{$X$}
\State $(\Lambda, Y) \gets \mathrm{eigh}(X^\dagger HX)$;\quad
$X \gets XY$;\ $HX \gets HX\,Y$ \Comment{start from Ritz vectors}
\State $P \gets 0$;\ $HP \gets 0$ \Comment{same shape as $X$}
\For{$t = 1, \dots, t_{\max}$}
\State $R \gets HX - X\,\mathrm{diag}(\Lambda)$;\quad
$r_n \gets \|R_{:,n}\|$ \Comment{residual of each band}
\State \textbf{if} $r_n < \epsilon$ for all bands that are needed
\textbf{then stop}
\State $W \gets K \circ R$ \Comment{preconditioner Eq.~\eqref{eq:teter},
each column with its own $E_{\mathrm{kin}}$}
\State $W_{:,n} \gets 0$ where $r_n < \epsilon/10$
\Comment{``soft locking'': converged bands get no new direction}
\State \textbf{twice:}\ \ $W \gets W - Q\,(Q^\dagger W)$ with $Q = [X, P]$;\quad
$W \gets$ \Call{Orthonormalize}{$W$, $10^{-10}$}
\State $HW \gets$ \Call{ApplyH}{$W$} \Comment{the only Hamiltonian
application per iteration}
\State $S \gets [X, W, P]$;\quad $A \gets S^\dagger\,[HX, HW, HP]$;\quad
$A \gets \tfrac12(A + A^\dagger)$ \Comment{$3N_b\times3N_b$, double precision}
\State $A_{jj} \gets A_{jj} + \Lambda_{\mathrm{big}}$ for every $j$ whose
column of $S$ is zero \Comment{$\Lambda_{\mathrm{big}} = 10^6$\,Ha: park
empty directions out of the way}
\State $(\Theta, Y) \gets \mathrm{eigh}(A)$;\quad
$\Lambda \gets \Theta_{1..N_b}$;\quad keep the first $N_b$ columns of $Y$,
split by rows as $(Y_X; Y_W; Y_P)$
\State $P \gets W Y_W + P\, Y_P$;\quad $HP \gets HW\, Y_W + HP\, Y_P$
\Comment{the step just taken}
\State $X \gets X Y_X + P$;\quad $HX \gets HX\, Y_X + HP$
\Comment{new Ritz vectors; no extra $H$ applications}
\State \textbf{twice:}\ \ $C \gets X^\dagger P$;\ \ $P \gets P - XC$;\ \
$HP \gets HP - HX\,C$;\ \
$(P, HP) \gets$ \Call{Orthonormalize}{$P$, $10^{-6}$, $HP$}
\Comment{$P \perp$ new $X$, orthonormal}
\EndFor
\State \Return $\Lambda$ (eigenvalues, ascending), $X$ (eigenvectors)
\end{algorithmic}
\end{algorithm}

Note the bookkeeping: $HX$ and $HP$ are updated by the same linear
combinations as $X$ and $P$, so the expensive operation is applied only to
$W$ --- and it is applied \emph{after} $W$ has been orthogonalized, so that
$HW$ is exact for the vectors actually used. The first iteration needs no
special case: $P = 0$ contributes zero columns, which line 13 parks. The
new $X$ is orthonormal because $S$ is and the columns of $Y$ are.

Two remarks on the tolerance $\epsilon$. In a fixed-potential calculation
(Section~\ref{sec:nscf}) it is the accuracy you want, typically
$10^{-8}$\,Ha, and the loop stops when it is reached. Inside the SCF loop
only a fixed, small number of iterations is done per step; there set
$\epsilon = 0$, so that nothing is ever locked and there is no stopping
test --- which is also what a jitted program with a static iteration count
wants.

\paragraph{Why ``twice'' appears twice.} Lines 10 and 17 each repeat a
projection and an orthonormalization. This is not belt and braces.
\textsc{Orthonormalize} divides by $\sqrt{\sigma_j}$, which amplifies
rounding errors by up to $\tau^{-1/2}$ ($10^{5}$ for $W$, $10^{3}$ for
$P$, on top of the rescaling by the column norms, which for a converged
band are of order $10^{-10}$). After a single pass a nearly dependent
block is therefore orthonormal only to about $10^{-10}$, not $10^{-16}$.
The Rayleigh--Ritz step \emph{assumes} $S^\dagger S = \mathbf 1$, so the
defect is inherited by the new $X$ and grows from one iteration to the
next. The second pass starts from an almost orthonormal block, amplifies
nothing, and brings the defect down to rounding level. For $P$ the second
pass costs two small matrix products and no Hamiltonian application,
because $HP$ is carried along. Keep both passes: with a single pass on $P$
the iteration is fine while it converges and then, some tens of iterations
\emph{after} reaching rounding-level residuals, loses orthogonality and
diverges. (This was found the hard way: an implementation written from an
earlier version of this text diverged on the cosine-potential test of
Chapter~\ref{ch:verify}, Stage~1, while passing all tests on gold.)

\paragraph{Numerical stability --- read this before simplifying anything.}
LOBPCG as usually presented in three lines (take $S = [X, W, P]$ as is,
Cholesky-factorize $S^\dagger S$, solve) is a trap. It behaves perfectly
for the first ten or twenty iterations, which is all an SCF step uses; then,
once some bands are converged to rounding error, it \emph{diverges} --- not
to a slightly wrong answer but to eigenvalues of $-10^{3}$\,Ha and far
beyond --- usually without any error being raised. For the gold
Hamiltonian of this document, in double precision, this happens within
30--100 iterations at ordinary $\kk$-points, i.e.\ exactly in the tightly
converged band-structure and topology calculations. The mechanism is worth
understanding, because it defeats the obvious repairs too (for instance,
replacing Cholesky by an eigendecomposition of $S^\dagger S$ with a small
cutoff): when $W$ and $P$ are nearly parallel, the Ritz coefficients $Y$
become huge ($10^5$) with opposite signs, the new $P$ is a small difference
of large vectors, and the stored $HP$ --- updated by recurrence, never
recomputed --- stops being the image of $P$ under $\hat H$. That
inconsistency is amplified at every iteration, leaks into $HX$, and from
then on the ``Rayleigh quotients'' are no longer bounded below by the true
lowest eigenvalue. The design of Algorithm~\ref{alg:lobpcg} removes the
cause: all coefficient matrices are $O(1)$ because $S$ is orthonormal; the
only block whose image is obtained by recurrence, $P$, is transformed
exclusively by well-conditioned operations (line 17 discards directions
below $10^{-6}$ instead of amplifying them); and nearly dependent
directions in $W$ are discarded \emph{before} $\hat H$ is applied. In
tests on the gold Hamiltonian --- $\Gamma$, $X$, $L$, $W$, $K$ and generic
$\kk$, with and without SOC, several random starts each --- and on the
highly degenerate nearly-free-electron problems of Stage~1, with
$\epsilon = 10^{-10}$, $10^{-6}$ and $0$, this version converged in every
case and stayed converged (eigenvalues within $2\times10^{-9}$\,Ha of dense
diagonalization, $X^\dagger X = \mathbf 1$ to $10^{-13}$) for 400
iterations. One thing it does not promise: once the residuals have
reached the rounding floor ($10^{-9}$--$10^{-10}$\,Ha for a level with
many degenerate partners) they wander there instead of decreasing, so ask
for $\epsilon \ge 10^{-9}$\,Ha and stop when you get it.

\paragraph{What to expect.} For gold with $N_b = 20$ spinor bands and a
random start, the residuals of the lowest 14 bands fall below
$10^{-9}$\,Ha in 25--45 iterations; the eigenvalues are converged to
$10^{-10}$\,Ha well before that (the eigenvalue error is quadratic in the
residual). The top few bands of the block always converge slowest ---
which is why $N_b$ should exceed the number of bands you need by a buffer
of four or more.

\paragraph{How many iterations?} Inside an SCF loop you do \emph{not} need
tight convergence: 2--4 LOBPCG iterations per SCF step suffice (after
$\sim8$ in the very first step, which starts from random vectors), because
the potential will move anyway and the residuals keep shrinking from one
SCF step to the next. Warm-starting from the previous step's $X$ is what
makes the whole calculation fast.

\paragraph{Blocked Davidson (alternative).} Expand $V \gets [V, K\circ R]$ up
to a max subspace ($\sim 3 N_b$), Rayleigh--Ritz (Algorithms~\ref{alg:lobpcg}
and \ref{alg:ortho} serve as building blocks), restart with Ritz vectors.
Same kernels, but the growing/shrinking subspace means dynamic shapes ---
awkward under JIT; if you implement it, fix the schedule (e.g.\ always
exactly two expansions per restart) so shapes are static.

\section{Occupations, density, and the SCF loop}

After diagonalization at all $\kk$: find $\mu$, set occupations, and
accumulate the density matrix on the grid.

\begin{algorithm}[H]
\caption{Fermi level and occupations}
\label{alg:fermi}
\begin{algorithmic}[1]
\Function{Occupy}{$\varepsilon_{n\kk}$, $w_\kk$, $N_e$, $\sigma$}
\State $\mu_{\mathrm{lo}} \gets \min\varepsilon - 10\sigma$;\quad
$\mu_{\mathrm{hi}} \gets \max\varepsilon + 10\sigma$
\For{$80$ iterations} \Comment{fixed count: interval shrinks by $2^{80}$}
\State $\mu \gets \tfrac12(\mu_{\mathrm{lo}}+\mu_{\mathrm{hi}})$;\quad
$N(\mu) \gets \sum_{n\kk} w_\kk\, f\big((\varepsilon_{n\kk}-\mu)/\sigma\big)$
\State \textbf{if} $N(\mu) > N_e$ \textbf{then} $\mu_{\mathrm{hi}}\gets\mu$
\textbf{else} $\mu_{\mathrm{lo}}\gets\mu$
\EndFor
\State $x_{n\kk} \gets (\varepsilon_{n\kk}-\mu)/\sigma$;\quad
$f_{n\kk} \gets f(x_{n\kk})$;\quad
$S \gets \sum_{n\kk} w_\kk\, s(x_{n\kk})$
\State \Return $\mu$, $f$, $S$ \Comment{entropy term of the free energy is
$-\sigma S$}
\EndFunction
\end{algorithmic}
\end{algorithm}

\begin{algorithm}[H]
\caption{Density matrix from the spinor bands}
\label{alg:density}
\begin{algorithmic}[1]
\State $n_{\alpha\beta}(\rr) \gets 0$
\For{each $\kk$ (in chunks), each band $n$}
\State $u^{s}(\rr) \gets \dfrac{N_{\mathrm{grid}}}{\sqrt\Omega}\,
\mathrm{IFFT}\big(\text{scatter } c^{s}_{n\kk}\big)$ for $s =
\uparrow,\downarrow$ \Comment{normalization of Eq.~\eqref{eq:fourierconv}}
\State $n_{\alpha\beta}(\rr) \mathrel{+}= w_\kk f_{n\kk}\,
u^{\alpha}(\rr)\,\big[u^{\beta}(\rr)\big]^{*}$
\EndFor
\State $n \gets n_{\uparrow\uparrow}+n_{\downarrow\downarrow}$;\quad
$m_x \gets 2\,\mathrm{Re}\,n_{\uparrow\downarrow}$;\
$m_y \gets -2\,\mathrm{Im}\,n_{\uparrow\downarrow}$;\
$m_z \gets n_{\uparrow\uparrow}-n_{\downarrow\downarrow}$
\end{algorithmic}
\end{algorithm}

The check: $\frac{\Omega}{N_{\mathrm{grid}}}\sum_{\rr} n(\rr) = N_e$ to
rounding error, every iteration. For gold $\mathbf m$ comes out zero to
rounding error \emph{only} if both members of every Kramers pair are
equally occupied \emph{and} both $\kk$ and $-\kk$ are in the mesh. On a
time-reversal-reduced mesh (one point of each $\pm\kk$ pair with doubled
weight, as in the reference cases of Chapter~\ref{ch:verify}) the doubled
weights give the correct $n$ but a spurious, nonzero $\mathbf m$
(Eq.~\eqref{eq:mpair}): there, accumulate only
$n = \sum |u^\uparrow|^2 + |u^\downarrow|^2$ and set $\mathbf m = 0$
--- which is also the right choice for a first version of the code. (The
reference calculations of Chapter~\ref{ch:verify} do exactly this.)

\subsection{Why simple iteration diverges for metals: charge sloshing}

The naive update $n_{\mathrm{in}}^{(m+1)} = n_{\mathrm{out}}^{(m)}$ diverges
whenever the dielectric response is strong. To see why, perturb the input
density by $\delta n_{\mathrm{in}}$. The potential changes by
$K_{\mathrm{Hxc}}\,\delta n_{\mathrm{in}}$ (dominated by the Hartree
kernel $4\pi/G^2$), and the electrons respond to that with
$\delta n_{\mathrm{out}} = \chi_0 K_{\mathrm{Hxc}}\,\delta n_{\mathrm{in}}$,
where $\chi_0$ is the density response of the non-interacting system. So an
error in the input is multiplied by $\chi_0 K_{\mathrm{Hxc}}$ each
iteration. In a metal $\chi_0(G\to0) \to -g(\varepsilon_F)$, a finite
negative number, so for a Fourier component of wavevector $G$ the
multiplier is $\approx -k_{TF}^2/G^2$ with $k_{TF}^2 = 4\pi
g(\varepsilon_F)$: negative (the error changes sign --- charge
\emph{sloshes} back and forth) and larger than one in magnitude whenever
$G < k_{TF}$ (the oscillation grows). Simple linear mixing
$n_{\mathrm{in}} \gets n_{\mathrm{in}} + A\,R$ multiplies the error by
$1 - A\,(1 + k_{TF}^2/G^2)$ instead, which is stable only if
$A < 2/(1 + k_{TF}^2/G_{\min}^2)$ --- tiny for a long cell.

How bad is it for us? In the gold primitive cell the smallest nonzero
reciprocal vector is $|\mathbf b_i| = 1.41$\,Bohr$^{-1}$, larger than
$k_{TF} \approx 0.9$\,Bohr$^{-1}$, so even plain mixing with
$A \approx 0.5$ converges. Measured on a small gold test case (scalar
relativistic, $E_{\mathrm{cut}} = 20$\,Ha, $4^3$ mesh, iterations to
$\|R\| < 10^{-7}$): plain mixing with $A = 0.5$ needs 24 iterations, with
$A = 0.8$ it has not converged after 60, and with $A = 1$ it does not
converge at all (the $d$ electrons make gold's true screening stronger
than the free-electron estimate of $k_{TF}$); Kerker mixing with $A = 0.8$
needs 16; Kerker--Pulay needs 10. Sloshing becomes the dominant problem the
moment you build a supercell or a slab (a ten-layer slab with vacuum has
$G_{\min} \approx 0.1$\,Bohr$^{-1}$). Implement the two standard fixes now, used
together, so the code is ready:

\paragraph{(a) Kerker preconditioning.} Damp the long-wavelength part of the
density update in Fourier space:
\begin{equation}
\delta n(\GG) \;\gets\; A\, \frac{G^2}{G^2 + k_{TF}^2}\; \delta n(\GG),
\qquad k_{TF}^2 = 4\pi\, g(\varepsilon_F) \ \ (\text{use } k_{TF} \approx
1\ \mathrm{Bohr}^{-1},\ A \approx 0.8).
\label{eq:kerker}
\end{equation}
This is precisely the inverse of the factor $1 + k_{TF}^2/G^2$ found above:
after preconditioning, every wavelength converges at the same rate.

\paragraph{(b) Pulay / Anderson mixing (DIIS).} Keep the last $M \sim 8$
pairs of (input density, residual $R^{(m)} = n_{\mathrm{out}}^{(m)} -
n_{\mathrm{in}}^{(m)}$); find the linear combination of past inputs whose
extrapolated residual is smallest:
\begin{equation}
\min_{\{\gamma\}} \Big\| \sum_m \gamma_m R^{(m)} \Big\|^2
\quad \text{s.t.} \sum_m \gamma_m = 1
\;\;\Longrightarrow\;\;
\begin{pmatrix} \mathsf{B} & \mathbf 1\\ \mathbf 1^{\mathsf T} & 0
\end{pmatrix}
\begin{pmatrix} \bm\gamma \\ \lambda\end{pmatrix}
= \begin{pmatrix}\mathbf 0\\ 1\end{pmatrix},
\qquad
\mathsf B_{mm'} = \braket{R^{(m)} | R^{(m')}},
\end{equation}
then update
\begin{equation}
n_{\mathrm{in}}^{(\mathrm{new})}
= \sum_m \gamma_m\Big[ n_{\mathrm{in}}^{(m)}
+ \mathcal K\, R^{(m)} \Big],
\label{eq:pulayupdate}
\end{equation}
with $\mathcal K$ the Kerker operator \eqref{eq:kerker}. The logic: if the
residual depended linearly on the input, the combination
$\sum_m\gamma_m n_{\mathrm{in}}^{(m)}$ would have residual
$\sum_m\gamma_m R^{(m)}$, the smallest reachable from the history; one
preconditioned step is then taken from there.

\begin{algorithm}[H]
\caption{Kerker-preconditioned Pulay mixing (all arrays in Fourier space)}
\label{alg:pulay}
\begin{algorithmic}[1]
\Function{Mix}{$n_{\mathrm{in}}(\GG)$, $R(\GG)$, history}
\State append $(n_{\mathrm{in}}, R)$ to the history; drop the oldest pair if
more than $M$ are stored
\State $\mathsf B_{mm'} \gets \mathrm{Re}\sum_{\GG\ne0}
\dfrac{4\pi}{G^2}\,\big[R^{(m)}(\GG)\big]^{*} R^{(m')}(\GG)$;\quad
$\mathsf B \gets \mathsf B/\max_{mm'}|\mathsf B_{mm'}|$
\State solve the bordered system for $\bm\gamma$ by least squares
\Comment{it becomes ill-conditioned near convergence}
\State \Return $\sum_m \gamma_m\big[n^{(m)}_{\mathrm{in}}(\GG)
+ A\,\tfrac{G^2}{G^2+k_{TF}^2}\,R^{(m)}(\GG)\big]$
\EndFunction
\end{algorithmic}
\end{algorithm}

The metric $4\pi/G^2$ weights long wavelengths, matching the physics (it
is the Hartree energy of the residual). The $\GG=0$ component needs no
special treatment: every input and output density integrates to $N_e$, so
$R(\GG{=}0)=0$ identically, and because $\sum_m\gamma_m = 1$ the mixed
density keeps the correct electron count. The mixed density can dip
slightly negative at isolated grid points; the clamp inside the XC routine
(Chapter~\ref{ch:xc}) is all that is needed. With the first history entry
alone, $\gamma_1 = 1$ and the step is plain Kerker mixing. Mix all four
components $(n, \mathbf m)$ with the same $\bm\gamma$ (the $\mathbf m$
components without the Kerker factor); for gold, $\mathbf m$ stays zero but
the code path should exist.

\paragraph{Measuring convergence.} Throughout this document the residual
norm is
\begin{equation}
\|R\| = \Big[\int_\Omega \big|n_{\mathrm{out}}(\rr)-n_{\mathrm{in}}(\rr)\big|^2
\dd\rr\Big]^{1/2}
= \Big[\frac{\Omega}{N_{\mathrm{grid}}}\sum_{\rr}|R(\rr)|^2\Big]^{1/2}.
\label{eq:resnorm}
\end{equation}
Because the energy is stationary at self-consistency, its error is
\emph{quadratic} in $\|R\|$: $\|R\|\sim10^{-4}$ already gives the energy to
$\sim10^{-8}$\,Ha.

\subsection{The complete SCF driver}

\begin{algorithm}[H]
\caption{Self-consistent field for spinor plane-wave DFT}
\begin{algorithmic}[1]
\State \textbf{Geometry-time setup:} lattice, reciprocal vectors, $\GG$
spheres and masks per $\kk$, FFT grid and pad/gather indices, structure
factor, $V_{\mathrm{loc,ion}}(\GG)$, projector arrays
$\tilde\beta(\kk{+}\GG)$, $D^{\ell}$ blocks (with SOC), Ewald energy
\State $n_{\mathrm{in}}(\GG) \gets \frac{1}{\Omega}\sum_a Z_{\mathrm{ion}}\,
\ee^{-G^2 w^2/2}\,\ee^{-\ii\GG\cdot\RR_a}$ with $w = 1$\,Bohr
(a Gaussian blob per atom, integrating to $N_e$); $\mathbf m \gets 0$
\State $X_{\kk} \gets$ random initial blocks (Algorithm~\ref{alg:lobpcg},
line 1)
\For{$m = 1, \dots, m_{\max}$}
\State $\vH \gets 4\pi n_{\mathrm{in}}(\GG)/G^2$;\quad
$(\vxc, \mathbf B_{\mathrm{xc}}) \gets$ XC$(n_{\mathrm{in}}, \mathbf m, \nabla n_{\mathrm{in}})$
\State $V_{\mathrm{eff,loc}} \gets V_{\mathrm{loc,ion}} + \vH + \vxc$
(real-space grid)
\State \textbf{for all $\kk$ in parallel:} Algorithm~\ref{alg:lobpcg}
started from the current $X_\kk$, with $\epsilon = 0$ and $t_{\max} = 8$ in
the first pass, $3$ afterwards, on
$H_{\kk}[V_{\mathrm{eff,loc}}, \mathbf B_{\mathrm{xc}}]$
$\Rightarrow \varepsilon_{n\kk}, X_{\kk}$
\State $(\mu, f_{n\kk}, S) \gets$ \Call{Occupy}{$\varepsilon, w, N_e, \sigma$}
\Comment{Algorithm~\ref{alg:fermi}}
\State $(n_{\mathrm{out}}, \mathbf m_{\mathrm{out}}) \gets$ density from
$\{X_{\kk}, f\}$ \Comment{Algorithm~\ref{alg:density}}
\State $R \gets n_{\mathrm{out}} - n_{\mathrm{in}}$;\quad record
$\|R\|$ and both free energies \Comment{Eqs.~\eqref{eq:resnorm},
\eqref{eq:fdirect}, \eqref{eq:ftot}}
\State \textbf{if} $\|R\| < 10^{-7}$ \textbf{stop}
\State $n_{\mathrm{in}} \gets$ \Call{Mix}{$n_{\mathrm{in}}, R$, history}
\Comment{Algorithm~\ref{alg:pulay}, $M = 8$}
\EndFor
\State Keep $n_{\mathrm{in}}$, $V_{\mathrm{eff,loc}}$ and $\mu$; all post-processing
(band structure, DOS, parities, Wilson loops) uses
Section~\ref{sec:nscf}
\end{algorithmic}
\end{algorithm}

\paragraph{Convergence pathologies and cures.}
\begin{itemize}[leftmargin=2em]
\item \emph{Oscillating $F$, residual stuck}: smearing too small or
$\kk$-mesh too coarse; raise $\sigma$ to $0.02$\,Ha to converge, then
re-refine.
\item \emph{Slow creep}: mixing history too short, or Kerker $k_{TF}$
mismatched; try $A \in [0.3, 0.8]$.
\item \emph{Occupied/empty band flip-flops}: not enough empty bands; ensure
the highest computed band is $\gtrsim 10\sigma$ above $\mu$.
\item \emph{LOBPCG stagnation}: block too small (use $N_b \ge N_{\mathrm{occ}}
+ 8$).
\item \emph{Eigenvalues suddenly huge and negative}: the stored $HX$ or
$HP$ is no longer consistent with $X$ or $P$; you simplified
Algorithm~\ref{alg:lobpcg} (see ``Numerical stability'').
\item \emph{Electron count drifts}: the density normalization of
Algorithm~\ref{alg:density} or the $\kk$ weights are off; fix this before
looking at anything else.
\end{itemize}

\section{After self-consistency: calculations at fixed potential}
\label{sec:nscf}

Once the SCF has converged, the density and hence $V_{\mathrm{eff,loc}}$
are known and \emph{frozen}. The Kohn--Sham Hamiltonian is then a fixed
linear operator, and its eigenstates can be computed at any $\kk$
whatsoever --- not only on the SCF mesh --- at the cost of one eigensolve
each. Every remaining result of this document is obtained this way.

\begin{algorithm}[H]
\caption{Bands at arbitrary $\kk$-points in a fixed potential}
\label{alg:nscf}
\begin{algorithmic}[1]
\Function{SolveBands}{list of $\kk$;\ $V_{\mathrm{eff,loc}}$;\
$N_{\mathrm{want}}$;\ $\epsilon$}
\For{each $\kk$ (batched)}
\State basis $\gets$ \Call{BuildBasis}{$\kk$, $E_{\mathrm{cut}}$}
\Comment{same cutoff and FFT grid as the SCF}
\State run Algorithm~\ref{alg:lobpcg} from a random start with
$N_b = N_{\mathrm{want}}+6$ until the lowest $N_{\mathrm{want}}$ residuals
are below $\epsilon$
\EndFor
\State \Return $\varepsilon_{n\kk}$ and coefficients $c^s_{n\kk}(\GG)$ for
$n \le N_{\mathrm{want}}$
\EndFunction
\end{algorithmic}
\end{algorithm}

Use double precision and $\epsilon \sim 10^{-8}$\,Ha here. Its uses:
\begin{itemize}[leftmargin=2em]
\item \textbf{Band structure.} A path
$\Gamma$--$X$--$W$--$L$--$\Gamma$--$K$ through the FCC high-symmetry points
of Section~\ref{sec:bzintegration}, 20--30 points per segment (interpolate
linearly in fractional coordinates; plot against the cumulative Cartesian
path length), energies measured from $\mu$.
\item \textbf{Density of states.} On a mesh at least as fine as the SCF
mesh,
\begin{equation}
g(\varepsilon) = \sum_{n\kk} w_{\kk}\,
\delta_\sigma(\varepsilon - \varepsilon_{n\kk}),
\end{equation}
with $\delta_\sigma$ a normalized Gaussian of width $\sim 0.005$\,Ha. The
$d$-band complex appears as a massive peak $2$--$7$\,eV below $E_F$, split
by SOC.
\item \textbf{Wavefunction diagnostics.} Parities at the eight TRIM
(Section~\ref{sec:symmetry}) and overlaps between neighbouring
$\kk$-points (Chapter~\ref{ch:topology}) need the coefficients, not just
the eigenvalues --- which is why \textsc{SolveBands} returns them.
\end{itemize}
One caution: eigenvectors within a degenerate level --- and in gold every
level is at least two-fold degenerate --- are returned in an arbitrary
orthonormal combination, with arbitrary overall phases, different at every
$\kk$. Any quantity you compute from wavefunctions must be insensitive to
this freedom. Eigenvalues, densities, and traces or eigenvalues of
matrices built from \emph{complete} degenerate groups are; individual
matrix elements are not.

\begin{gpubox}{Parallel structure of the whole SCF step}
One SCF iteration is: [pointwise XC kernel] $\to$ [a few FFTs for
$\vH$, gradients] $\to$ [batched LOBPCG: per iteration one batched
$\hat H$-apply $=$ $2 N_k N_b$ FFT pairs $+$ a handful of batched GEMMs,
then one $3N_b$-dim and four $N_b$-dim \texttt{eigh} per $\kk$ (batched)] $\to$ [device-wide reductions
for $\mu$] $\to$ [$2 N_k N_b$ IFFTs $+$ fused reduction for the density]
$\to$ [Fourier-space mixing]. Every stage is either elementwise, an FFT
batch, a GEMM batch, or a small batched \texttt{eigh} --- there is no
sequential inner structure anywhere except the SCF iteration count itself
($\sim 20$--$40$ iterations) and the handful of LOBPCG steps inside each.
\end{gpubox}

\chapter{Mapping the Code onto JAX and the GPU}
\label{ch:gpu}

This chapter is the implementation blueprint: the JAX vocabulary you need,
data layouts, batching strategy, precision policy, and JIT discipline. No
program listings --- but every decision you will face is settled here.
Nothing in it requires a GPU: the same program runs unchanged on a laptop
CPU, where the small cases of Chapter~\ref{ch:verify} take a minute or two.
Develop there; move to an accelerator when the physics is right.

\section{The JAX you need}
\label{sec:jaxprimer}

JAX is NumPy with three additions: it runs on accelerators, it can
differentiate functions, and it can compile them. If you can write NumPy
you can write JAX (\texttt{import jax.numpy as jnp} and use \texttt{jnp}
where you would use \texttt{np}), provided you respect a few rules that
differ from NumPy. The complete list of what this project uses:
{\small
\begin{longtable}{@{}>{\raggedright\arraybackslash}p{0.27\textwidth}p{0.68\textwidth}@{}}
\toprule
JAX feature & What it does, and where you use it\\
\midrule
\endfirsthead
\toprule
JAX feature & What it does, and where you use it (continued)\\
\midrule
\endhead
\bottomrule
\endlastfoot
\texttt{jax.config.update(} \texttt{"jax\_enable\_x64", True)} & Must be the first
statement of the program. Without it JAX silently computes everything in
single precision and none of the reference numbers can be reproduced.\\
\texttt{jnp.fft.fftn}, \texttt{ifftn} with \texttt{axes=} & 3D
transforms over the last three axes of a batch; same normalization as
NumPy (Fourier conventions of Chapter~\ref{ch:solids}).\\
\texttt{jnp.einsum}, \texttt{@} & All projections, Gram matrices and
subspace rotations.\\
\texttt{jnp.linalg.eigh} & Hermitian eigenproblems (complex input is
fine); eigenvalues come back ascending.\\
\texttt{x.at[idx].set(v)}, \texttt{.add(v)} & JAX arrays are immutable:
\texttt{x[idx] = v} is an error. This form returns a modified copy; it
is how coefficients are scattered into the FFT box. Reading
(\texttt{x[idx]}) works as in NumPy. \emph{Pitfall:} the padded entries
of a $\kk$-point's index array must point somewhere; if they point to box
index 0 and you use \texttt{.set}, a padded zero can overwrite the true
$\GG = 0$ coefficient. Use \texttt{.add} (padded coefficients are zero, so
they add nothing).\\
\texttt{jnp.where(cond, a, b)} & Replaces \texttt{if} on array values
(masks, thresholds, soft locking). Both branches are evaluated, so make
both finite: for ``$1/s$ where $s > \tau$, else $0$'' write
\texttt{jnp.where(s > tau, 1/jnp.where(s > tau, s, 1.0), 0.0)} --- the inner
\texttt{where} only makes the discarded branch finite, the outer one still
supplies the zeros (dropping it gives $1$, not $0$, below the threshold:
in Algorithm~\ref{alg:ortho} that would keep dependent directions). Never
write \texttt{jnp.where(s > tau, 1/s, 0.0)} with a zero in \texttt{s}.\\
\texttt{jax.jit(f)} & Compiles \texttt{f} for fixed array shapes and
dtypes. Inside, Python \texttt{if}/\texttt{while} may not depend on array
\emph{values}; loops with a fixed count are simply unrolled.\\
\texttt{jax.vmap(f)} & Turns a function written for one $\kk$-point (or
one grid point) into one for a whole batch. Write Algorithms
\ref{alg:applyh}--\ref{alg:lobpcg} for a single $\kk$ and
\texttt{vmap} them.\\
\texttt{jax.lax.map(f, xs)} & Like \texttt{vmap} but sequential: use it
where a full batch would not fit in memory (density build, large
meshes).\\
\texttt{jax.grad(f)} & Derivative of a scalar function. Gives $\vxc$
from the XC energy density (Chapter~\ref{ch:xc}): \texttt{grad} of
$f(n,\varsigma)$ with respect to each argument, then \texttt{vmap} over
grid points.\\
\end{longtable}}
\noindent Three habits prevent most JAX-specific trouble. First, do all
geometry-time setup (the integer bookkeeping of Algorithm~\ref{alg:basis},
projector tables, Ewald) in ordinary NumPy and convert the finished arrays
with \texttt{jnp.array}; only the numerical kernels need to be JAX.
Second, keep the outer SCF loop, the Fermi-level bisection and the Pulay
solve in plain Python at first, calling jitted kernels from it --- you can
print and inspect everything. Third, the first call of a jitted function
is slow (it compiles) and it recompiles whenever a shape changes; give
every $\kk$-point the same padded shape so that this happens once.

\section{The parallelism inventory}

Plane-wave DFT exposes four nested, independent axes of parallelism:
\begin{center}
\begin{tabular}{@{}llll@{}}
\toprule
Axis & Size (Au, $q{=}11$) & Independence & GPU mapping\\
\midrule
$\kk$-points & $N_k \sim 10^3$ (full $12^3$ mesh)
& fully independent linear algebra & leading batch dim\\
Bands & $N_b \sim 24$ & coupled only via Rayleigh--Ritz & block/GEMM dim\\
Spinor & 2 & coupled by SOC and $\mathbf B_{\mathrm{xc}}$
& folded into band dim\\
$\GG$ / grid & $N_G \sim 2{\times}10^3$, grid $\sim 40^3$
& elementwise / FFT & thread-level (XLA)\\
\bottomrule
\end{tabular}
\end{center}
The design rule: \textbf{never write a loop over any of these axes ---
express every operation as one array program over the full
$[N_k, N_b, 2, N_G^{\max}]$ tensor}, and let vectorized primitives
(\texttt{vmap}, batched FFT, batched GEMM, batched \texttt{eigh}) fill the
device. A $12^3$ mesh $\times$ 24 bands $\times$ 2 spin components $\approx
8\times10^4$ FFTs of size $\sim 36^3$ per Hamiltonian application: enough
to saturate any current accelerator (and too many to hold in memory at
once --- see the memory budget below).

\section{Data layout and static shapes}

\begin{itemize}[leftmargin=2em]
\item \textbf{Wavefunctions} $C$: complex array $[N_k, N_b, 2, N_G^{\max}]$.
Each $\kk$-point's cutoff sphere has a different size $N_G(\kk)$; pad to
$N_G^{\max} = \max_\kk N_G(\kk)$ and carry a boolean mask
$[N_k, N_G^{\max}]$. Masked entries are kept exactly zero (multiply by the
mask after every operation that could pollute them). Padding overhead is a
few percent; the reward is one static shape for everything, which JAX/XLA
requires for compilation and which makes batched FFT/GEMM trivial.
\item \textbf{Sphere $\leftrightarrow$ box maps}: for each $\kk$, a
precomputed integer array of linear FFT-box indices for its $\GG$ sphere.
Pad-and-scatter ($\texttt{zeros(box).at[idx].add(coeffs)}$-style) and
gather are the only ``indexing'' operations in the code.
\item \textbf{Fields} ($n$, $\mathbf m$, potentials): real/complex arrays on
the dense grid $[N_1, N_2, N_3]$ (and their Fourier images). Mixing history:
$[M_{\mathrm{hist}}, 4, N_1, N_2, N_3]$ in Fourier space.
\item \textbf{Projectors}: complex $[N_k, N_{\mathrm{proj}}, N_G^{\max}]$;
coupling blocks $D$: complex
$[N_{\mathrm{proj}} \times 2, N_{\mathrm{proj}} \times 2]$ block-diagonal
(store dense; it is tiny).
\end{itemize}

\section{The seven kernels}

The entire program reduces to seven device kernels; everything else is glue.
\begin{enumerate}[leftmargin=2em]
\item \textbf{Batched zero-pad / gather} between sphere and FFT box.
\item \textbf{Batched 3D FFT / IFFT} over the $[N_k \times N_b \times 2]$
batch. In JAX, \texttt{jnp.fft.fftn} with \texttt{axes=(-3,-2,-1)} transforms
only the last three axes of a batched array such as
$[N_k, N_b, 2, N_1, N_2, N_3]$ and is therefore already batched (without
\texttt{axes=} it would transform every axis); sizes with small prime factors keep cuFFT on fast paths.
\item \textbf{Elementwise field kernels}: kinetic multiply, local-potential
multiply, XC evaluation, Kerker filter, preconditioner, smearing functions.
XLA fuses chains of these into single kernels automatically.
\item \textbf{Batched tall GEMMs}: projections $B^{\dagger}C$ and expansions
$B(DP)$; subspace products $S^{\dagger}(HS)$; block updates $SY$. These are
\texttt{einsum}s over the $[N_k]$ batch.
\item \textbf{Batched small \texttt{eigh}}: the $3N_b \times 3N_b$
Rayleigh--Ritz problem and the $N_b\times N_b$ Gram matrices of
\textsc{Orthonormalize} (Algorithm~\ref{alg:ortho}; four calls per
iteration), batched over $\kk$
(\texttt{vmap(eigh)}). No Cholesky and no QR anywhere: \texttt{eigh} never
fails, which matters because a jitted program cannot branch on a failed
factorization.
\item \textbf{Global reductions}: density accumulation
($\sum_{n\kk} w f |\psi|^2$ --- an einsum-reduction after the IFFT batch),
electron count for Fermi bisection, residual norms, energies.
\item \textbf{Scatter of constants}: structure factors, projector tables,
Ewald --- geometry-time, negligible.
\end{enumerate}

\section{Precision policy}

\begin{itemize}[leftmargin=2em]
\item Run wavefunctions and FFTs in \textbf{complex64} (single precision)
during the bulk of the SCF: half the memory and typically $2$--$4\times$
the speed of double precision (more on consumer GPUs, whose double-precision
throughput is throttled).
\item Single precision has a floor: residuals cannot be pushed much below
$\sim10^{-5}$ and energies carry $\sim10^{-6}$ relative noise. Finish with
a few SCF iterations --- and do all band-structure, $E(V)$ and topology
post-processing --- in \textbf{complex128}. The tight tolerances quoted
elsewhere in this document ($10^{-8}$ and below) refer to double
precision.
\item Accumulate \textbf{energies, the density, and the small Gram and
Rayleigh--Ritz matrices in float64/complex128}. The subspace matrices are
tiny, so promoting them is free and removes the main weakness of
single-precision LOBPCG (loss of orthogonality). In single precision you
will also need to raise the discard thresholds of
Algorithm~\ref{alg:lobpcg} --- as a starting point, to about
$10^{-5}$ (for $W$) and $10^{-3}$ (for $P$): a direction that is linearly
dependent to seven digits carries no information in a seven-digit
arithmetic.
\item \textbf{Write and validate the whole program in double precision
first} (every reference number of Chapter~\ref{ch:verify} is a
double-precision number). Single precision is an optimization for later:
be \emph{explicit} about dtypes at every array creation, keep full
float64 as a debug switch --- energy differences between the two runs
localize precision bugs instantly.
\item Target after the double-precision finish: total energy reproducible
to $\lesssim 10^{-6}$\,Ha/atom; eigenvalues to $10^{-7}$\,Ha.
\end{itemize}

\section{JIT discipline}

\begin{itemize}[leftmargin=2em]
\item \textbf{One jitted SCF step}: a single function
(fields, $X$, history) $\mapsto$ (fields', $X'$, history', scalars),
containing the eigensolver steps, Fermi search, density build, and mixing.
The Python-level SCF loop calls it repeatedly; convergence is checked on the
returned scalars. (A fully on-device \texttt{lax.while\_loop} SCF is
possible but makes debugging miserable --- do it last, if ever.)
\item \textbf{Static shapes everywhere}: fixed $N_b$, fixed LOBPCG
iterations per SCF step, fixed mixing history length (a ring buffer),
fixed bisection iteration count. Convergence-dependent behavior is expressed
through \texttt{where}-masks, not shape changes.
\item \textbf{No host--device sync inside the step}: never branch Python
control flow on device values inside the jitted region; return diagnostics
as arrays.
\item \textbf{Multi-GPU} (when you outgrow one device): shard the leading
$\kk$ axis (\texttt{shard\_map}); the only cross-device
communication per SCF step is (i) an all-reduce of the density and (ii) an
all-reduce of the eigenvalue histogram for the Fermi search --- a few MB.
This near-embarrassing parallelism over $\kk$ is why the algorithmic choices
of Chapter~\ref{ch:algorithms} scale.
\item \textbf{Autodiff opportunities}: XC potentials via \texttt{grad} of the
energy density (Ch.~\ref{ch:xc}); later, forces and stress via
differentiating the total energy with respect to atomic positions and
strain --- the Hellmann--Feynman theorem emerges automatically, provided
you differentiate the energy expression at \emph{fixed}, converged
wavefunctions and occupations (differentiating through a truncated SCF
loop gives something else). Design the
energy as a pure function of (parameters, density, wavefunctions) from day
one.
\end{itemize}

\section{Memory budget (sanity check)}

Full $\Gamma$-centered $12^3$ mesh, $N_b = 24$, $N_G^{\max} \approx 1500$ ($E_{\mathrm{cut}}
= 40$\,Ha; $\approx 2500$ at $60$\,Ha), complex64:
wavefunctions $1728 \times 24 \times 2 \times 1500 \times 8\,\mathrm{B}
\approx 1.0$\,GB. LOBPCG holds six such arrays ($X, W, P$ and their images
under $H$), so budget $\sim 6$\,GB; time-reversal symmetry halves $N_k$
for free (Section~\ref{sec:symmetry}). The FFT boxes are the real constraint: all states at once on a
$36^3$ grid would be $1728\times24\times2\times36^3\times8\,\mathrm{B}
\approx 31$\,GB, so the $\hat H$-apply and the density build \emph{must} be
processed in $\kk$-chunks (a static chunk size, scanned with
\texttt{lax.map}; 64 $\kk$-points per chunk is $\sim 1$\,GB). Fields and
mixing history are megabytes.

\begin{keybox}{Recommended parameters for the first converged gold run}
HGH PBE $q{=}11$; $E_{\mathrm{cut}} = 40$\,Ha (FFT grid $36^3$);
$\Gamma$-centered $12{\times}12{\times}12$ mesh; Methfessel--Paxton
order 1, $\sigma = 0.01$\,Ha; $N_b = 24$ spinor bands; LOBPCG: 8
iterations in the first SCF step, 3 per step afterwards, Teter
preconditioner; Pulay history 8, Kerker $A = 0.8$, $k_{TF} = 1.0$;
convergence $\|R\| < 10^{-7}$ (Eq.~\eqref{eq:resnorm}). Expect about ten
SCF iterations. Then converge upward: $E_{\mathrm{cut}} \to 60$, mesh
$\to 16^3$, and verify observables move by less than your target
accuracy. For debugging, the small case of Chapter~\ref{ch:verify}
(Stage~3, Table~\ref{tab:smallcases}) runs in about a minute on a CPU.
\end{keybox}

\chapter{The Topology of Gold's Bands}
\label{ch:topology}

Everything computed so far --- energies, bands, densities of states --- is a
\emph{number that varies continuously} when you change the lattice constant
or the functional a little. This chapter extracts a quantity of a different
kind: an \emph{integer}, $0$ or $1$, that is a property of the
wavefunctions rather than of the energies and that cannot change under
small perturbations. It is the $\mathbb{Z}_2$ topological invariant.
Gold's value is $1$ (``nontrivial''), and the well-known surface state of
the Au(111) surface is its experimental fingerprint.

Nothing new has to be built. The three ingredients are already in your
hands: the even/odd character of zone-boundary states
(Section~\ref{sec:nfe}), parity and time reversal as permutations of the
coefficient array (Section~\ref{sec:symmetry}), and the fixed-potential
band solver \textsc{SolveBands} (Section~\ref{sec:nscf}). The chapter
has two halves. Sections~\ref{sec:berry}--\ref{sec:z2def} explain what the
invariant \emph{is}, in terms of where the electronic charge sits in the
unit cell and how it moves; this leads to a general algorithm, the Wilson
loop. Sections~\ref{sec:fukane}--\ref{sec:goldparity} then show that
inversion symmetry collapses the whole computation to multiplying eight
signs. You will do both, and they must agree.

\section{What ``topological'' means}

A doughnut and a coffee mug both have exactly one hole. You can deform one
into the other smoothly; you cannot get rid of the hole without tearing.
The number of holes is a \emph{topological invariant}: an integer that
survives every smooth deformation.

Band structures have invariants of the same kind. Take a group of bands
separated from all higher bands by an energy gap. Deform the Hamiltonian
smoothly (strain the crystal, change the potential) without ever closing
that gap and without breaking time-reversal symmetry. Some properties of
the group of wavefunctions cannot change under such a deformation, and for
systems with spin--orbit coupling one of them is a single bit
$\nu_0 \in \{0, 1\}$.
\begin{itemize}[leftmargin=2em]
\item $\nu_0 = 0$: the bands can be deformed into those of isolated atoms
pulled far apart. Nothing special.
\item $\nu_0 = 1$: they cannot. Somewhere in the Brillouin zone the bands
are ordered ``the wrong way round'' compared with separated atoms --- a
\textbf{band inversion} --- and no smooth change can undo it.
\end{itemize}
Why care about a bit? Because of what happens at a surface. Outside the
crystal is vacuum, which is trivially $\nu_0 = 0$. If the crystal has
$\nu_0 = 1$, the bit has to change across the surface, and it can only
change where the gap closes. So the surface is \emph{forced} to host
electronic states inside the bulk gap, and SOC makes them spin-polarized
with the spin locked to the direction of motion. Those surface states are
what experiments see.

To say precisely what the bit measures we need one concept that has not
appeared yet.

\section{Where is the charge? Berry phases and Wannier centers}
\label{sec:berry}

\paragraph{The question.} A filled band contributes a lattice-periodic
charge density. Ask a seemingly innocent question: \emph{where, within the
unit cell, is the center of that band's charge?} The density alone cannot
answer it --- any periodic density can be assigned to cells in many ways.
The answer is hidden in how the cell-periodic functions $u_{n\kk}$
\emph{change with} $\kk$.

\paragraph{One band in one dimension.} Take a single isolated band on a 1D
lattice of period $a$, with $u_k(x)$ on a mesh $k_j = 2\pi j/(N a)$,
$j = 0,\dots,N{-}1$. Each $u_{k_j}$ comes out of the eigensolver with an
arbitrary phase, so an individual overlap $\braket{u_{k_j}|u_{k_{j+1}}}$
means nothing by itself. But multiply the overlaps around the closed loop
of the Brillouin zone,
\begin{equation}
\mathcal W = \braket{u_{k_0}|u_{k_1}}\braket{u_{k_1}|u_{k_2}}\cdots
\braket{u_{k_{N-1}}|u_{k_N}},
\qquad u_{k_N}(x) \equiv \ee^{-\ii 2\pi x/a}\,u_{k_0}(x),
\label{eq:berry1d}
\end{equation}
and every state appears once as a bra and once as a ket: all the arbitrary
phases cancel. (The last factor closes the loop: $k_N = k_0 + 2\pi/a$ is
the same Bloch state as $k_0$, and $\psi = \ee^{\ii kx}u_k$ then fixes
$u_{k_N}$ as stated.) The phase of this product,
\begin{equation}
\phi = -\,\mathrm{Im}\ln \mathcal W \pmod{2\pi},
\end{equation}
is the \textbf{Berry phase} of the band (in one dimension also called the
Zak phase). It is a genuine, computable property of the band, and its
meaning is the answer to our question:
\begin{equation}
\boxed{\;\bar x = \frac{\phi}{2\pi}\,a \pmod a\;}
\label{eq:wcc}
\end{equation}
is the center of charge of the band, measured from the origin of the cell
--- the \textbf{Wannier charge center}. (``Wannier'' because it is the
mean position of the localized Wannier function built from the band;
``mod $a$'' because the charge center is only defined up to a lattice
translation.)

\emph{Why this is plausible.} In $\kk$-space the position operator acts as
$\ii\,\partial/\partial k$ on $u_k$ (just as momentum acts as
$-\ii\,\partial/\partial x$ in real space). For neighbouring mesh points,
$\braket{u_k|u_{k+\delta}} \approx 1 + \delta\braket{u_k|\partial_k u_k}
\approx \exp\!\big(-\ii\delta\, \braket{u_k|\ii\partial_k|u_k}\big)$, so
$-\mathrm{Im}\ln$ of the product adds up
$\delta \times(\text{local expectation of position})$ around the zone;
dividing by the length $2\pi/a$ of the zone gives the average
position. A check you can do by hand: for a band made of one atomic
orbital sitting at position $x_0$ in each cell (tight-binding limit),
$u_k(x) \propto \sum_R \ee^{-\ii k(x - x_0 - R)}\varphi(x-x_0-R)$, and
Eq.~\eqref{eq:berry1d} gives $\phi = 2\pi x_0/a$ (all interior overlaps
are real and positive on a fine mesh; the whole phase comes from the
closing factor): $\bar x = x_0$. The charge is where the orbital is.

\paragraph{Several bands: the Wilson loop.} For a group of $N_{\mathrm{occ}}$
bands (which may cross or be degenerate among themselves, as long as the
group is separated from the bands above), individual bands have no
identity, so the overlaps become matrices
\begin{equation}
M^{(\kk,\kk')}_{mn} = \braket{u_{m\kk}|u_{n\kk'}}
= \sum_{\GG,s} c^{s\,*}_{m\kk}(\GG)\, c^{s}_{n\kk'}(\GG),
\qquad m,n = 1,\dots,N_{\mathrm{occ}},
\label{eq:overlapM}
\end{equation}
(sum over the plane waves $\GG$ common to both spheres), and the loop
product becomes a matrix, the \textbf{Wilson loop}
\begin{equation}
\mathcal W = M^{(\kk_0,\kk_1)}\, M^{(\kk_1,\kk_2)} \cdots
M^{(\kk_{N-1},\,\kk_N)} .
\label{eq:wilson}
\end{equation}
Under a change of basis within the group at any intermediate $\kk_j$ the
product is unchanged; a change at $\kk_0$ conjugates $\mathcal W$ by a
unitary matrix. Its \emph{eigenvalues} are therefore well defined. On a
fine mesh they have modulus $\approx 1$, and their phases
$\theta_j \in (-\pi,\pi]$, $j = 1,\dots,N_{\mathrm{occ}}$, give
$N_{\mathrm{occ}}$ charge centers $\bar x_j = a\,\theta_j/2\pi$: the
positions of the $N_{\mathrm{occ}}$ electrons per cell that the group of
bands holds. (Sign conventions for $\theta$ differ between texts; none of
what follows depends on the choice.)

\paragraph{In three dimensions: centers that depend on $\kk_\perp$.} In
the crystal, run the loop along one reciprocal direction only: fix the
fractional coordinates $k_2, k_3$ and let $k_1$ go once around,
$\kk_j = (j/N_1,\,k_2,\,k_3)$. The last overlap again closes the loop with
the \emph{same} states as at $\kk_0$, relabeled: since
$u_{\kk+\mathbf b_1}(\rr) = \ee^{-\ii\mathbf b_1\cdot\rr}u_{\kk}(\rr)$,
\begin{equation}
c_{\kk_0+\mathbf b_1}(\GG) = c_{\kk_0}(\GG + \mathbf b_1)
\label{eq:closure}
\end{equation}
--- a shift of the Miller index $m_1$ by one; do \emph{not} solve again at
$\kk_0 + \mathbf b_1$, which would return the same states with new random
phases and destroy the cancellation. The result is a set of centers
$\bar x_j(k_2,k_3)$ along $\mathbf a_1$ that depend on the transverse
momentum: for each $(k_2,k_3)$ the crystal behaves like a 1D chain along
$\mathbf a_1$, and $\theta_j(k_2,k_3)$ tells you where its electrons sit.

\section{The \texorpdfstring{$\mathbb{Z}_2$}{Z2} invariant: how the centers flow}
\label{sec:z2def}

Fix $k_3$ and watch the $N_{\mathrm{occ}}$ centers $\theta_j(k_2)$ as $k_2$
goes across the zone. They trace continuous curves on a cylinder ($\theta$
is an angle). Time-reversal symmetry constrains these curves in two ways.
\begin{itemize}[leftmargin=2em]
\item $\theta_j(-k_2) = \theta_j(k_2)$ as a set: the picture for
$k_2\in[-\tfrac12,0]$ is the mirror image of that for $[0,\tfrac12]$, so
half the zone suffices.
\item At $k_2 = 0$ and $k_2 = \tfrac12$ the 1D chain is itself
time-reversal invariant, and by Kramers' theorem its centers are
\emph{doubly degenerate}: the curves meet in pairs at both ends of the
half zone.
\end{itemize}
Between the two ends the pairs split. The only question topology asks is
how they recombine (compare the two panels of Fig.~\ref{fig:wcc}): each curve can return to the
partner it started with, or the curves can \emph{switch partners}, each
ending up paired with a curve that started in a different pair. Switching
partners cannot be undone by any smooth deformation --- to do so two curves
would have to pass through each other, and Kramers degeneracy protects the
crossings at the ends. The parity of the number of partner switches is the
invariant of the plane:
\begin{equation}
z(k_3) = \Big(\text{number of crossings of the line }
\theta = \theta_{\mathrm{ref}} \text{ by the curves } \theta_j(k_2),\
0\le k_2\le\tfrac12\Big) \bmod 2 .
\label{eq:zcount}
\end{equation}
If every curve returns to its partner, any horizontal line is crossed an
even number of times; if partners are switched, some curve must travel
from one pair to the next and crosses the line an odd number of times.
The result does not depend on $\theta_{\mathrm{ref}}$. Physically: as $k_2$
advances through half the zone, a trivial band group shuffles its charge
centers and puts them back; a nontrivial one hands a charge center over
from one Kramers pair to the next, a net flow that cannot be stopped
without closing the gap.

A three-dimensional crystal has two independent time-reversal-invariant
planes of this kind for each direction, $k_3 = 0$ and $k_3 = \tfrac12$.
The \textbf{strong index} and one of the three \textbf{weak indices} are
\begin{equation}
\nu_0 = \big[z(k_3{=}0) + z(k_3{=}\tfrac12)\big] \bmod 2,
\qquad
\nu_3 = z(k_3{=}\tfrac12),
\label{eq:nu0fromz}
\end{equation}
and $\nu_1$, $\nu_2$ follow from the planes $k_1 = \tfrac12$,
$k_2 = \tfrac12$. The strong index is the robust one: $\nu_0 = 1$ means
the two parallel planes disagree, which is what guarantees a protected
state on \emph{every} surface (in a metal such as gold: inside the
projected direct gap, not in an absolute gap;
Section~\ref{sec:directgap}). The full label is written
$(\nu_0;\nu_1\nu_2\nu_3)$.

\section{But gold is a metal}
\label{sec:directgap}

The construction needs a group of bands separated from the rest by a gap,
and a metal has no gap at the Fermi level. Look closer, though: the
definition never mentions the Fermi level or the occupations. It needs
only that the highest band of the group and the next band do not touch
\emph{at any single $\kk$} --- a ``direct gap'' everywhere, even if the
top of the group at one $\kk$ lies above the bottom of the next band at
another. Gold has exactly this between Kramers pair 6 and pair 7 (spinor
bands 12 and 13). Pairs 1--6 are the five $d$ bands plus the $sp$ band
that crosses $E_F$; pair 7 lies above them at every $\kk$:
by $5.1$\,eV at $X$, $4.3$\,eV at $L$ and $19$\,eV at $\Gamma$ in the
reference calculation of Chapter~\ref{ch:verify}. Away from the TRIM the
gap is much smaller: it falls to $0.07$\,eV on the line $W$--$L$, and a
search over the zone found a minimum of about $0.02$\,eV at a point on a
mirror plane, $\kk \approx \tfrac{2\pi}{a}(0.68, 0.68, 0.18)$ (equivalent,
by a reciprocal lattice vector and cubic symmetry, to
$\tfrac{2\pi}{a}(0.82, 0.32, 0.32)$ inside the first zone). Small, but
not zero.
So the group is $N_{\mathrm{occ}} = 12$ spinor bands, six Kramers pairs,
regardless of the fact that part of pair 6 is empty. (That gold's
invariant defined this way is nontrivial was pointed out by Yan and
co-workers in 2015.) The first step of any topology calculation on a metal
is therefore to \emph{verify the direct gap} on a fine mesh.

Two practical warnings follow from these numbers. First, a coarse mesh
does not see the small gaps: the $8^3$ SCF mesh reports a minimum of
$0.54$\,eV and a $20^3$ mesh $0.27$\,eV, so finish the check with a local
minimization of $\varepsilon_{13}-\varepsilon_{12}$ started from the
smallest values found. Second, it is spin--orbit coupling that keeps the
two groups apart at such points: with $k^\ell = 0$ the same search finds
the sixth and seventh bands \emph{touching} (gap below $10^{-4}$\,eV) at
several points of the zone --- without SOC, bands of different mirror
symmetry are free to cross. The invariant of gold is a spin--orbit effect not because
SOC changes any parity, but because SOC makes the group of bands well
defined.

\section{The shortcut for inversion-symmetric crystals: parities}
\label{sec:fukane}

Following charge centers across the zone works for any crystal. For a
crystal with inversion symmetry Fu and Kane (2007) proved that the answer
is already fixed by the parities $\xi_m(\Gamma_i) = \pm1$ of the Kramers
pairs at the eight TRIM $\Gamma_i$ (Section~\ref{sec:symmetry}). Multiply
the parities of the $N$ pairs of the group at each TRIM, then multiply the
eight results:
\begin{equation}
\delta_i = \prod_{m=1}^{N} \xi_m(\Gamma_i) ,
\qquad\qquad
\boxed{\;(-1)^{\nu_0} = \prod_{i=1}^{8} \delta_i\;}
\label{eq:fukane}
\end{equation}
The weak index $\nu_j$ is the same product restricted to the four TRIM
with $n_j = 1$.

\paragraph{Why parities know about charge centers.} The connection to the
previous section is short enough to sketch. Consider a Wilson loop along
$k_1$ at one of the special values of $(k_2,k_3)$ where the 1D chain is
time-reversal invariant; the loop passes through exactly two TRIM,
$\Gamma_a$ (at $k_1 = 0$) and $\Gamma_b$ (at $k_1=\tfrac12$). Inversion
maps a charge center $\bar x$ to $-\bar x$, so a center that is alone must
sit at a point that is its own mirror image: $\bar x = 0$ (on the atom,
$\theta = 0$) or $\bar x = a/2$ (between atoms, $\theta = \pi$). Which one?
Exactly as in the two-plane-wave model of Section~\ref{sec:nfe}: a band
whose states at $\Gamma_a$ and $\Gamma_b$ have the \emph{same} parity is
centered on the atom; if the parities are \emph{opposite}, on the midpoint.
For the whole group (centers may also come in mirror-image sets
$\pm\theta$, which cancel),
\begin{equation}
\prod_{\text{pairs } m} \ee^{\ii\theta_m}
= \prod_{m=1}^{N} \xi_m(\Gamma_a)\,\xi_m(\Gamma_b) = \delta_a\delta_b ,
\label{eq:zakparity}
\end{equation}
one factor per Kramers pair: the number of pairs centered at
$\theta = \pi$ is odd precisely when $\delta_a\delta_b = -1$. Now compare
the two ends of the half zone in Eq.~\eqref{eq:zcount}. If the number of
pairs sitting at $\theta=\pi$ changes from even at $k_2 = 0$ to odd at
$k_2=\tfrac12$, one center must have travelled from $0$ to $\pi$ and
thereby crossed a reference line once: $z = 1$. Hence
$(-1)^{z} = \delta_a\delta_b\delta_c\delta_d$ over the four TRIM of the
plane, and Eq.~\eqref{eq:nu0fromz} multiplies the two planes together into
Eq.~\eqref{eq:fukane}. You will verify Eq.~\eqref{eq:zakparity}
numerically below.

\paragraph{The intuition.} For atoms pulled far apart, every band is built
from one kind of atomic orbital and has the same parity $(-1)^\ell$ at
every TRIM, so each band contributes $(\pm1)^8 = +1$ and $\nu_0 = 0$. To
get $\nu_0 = 1$, an even state and an odd state must have swapped places,
relative to the atomic ordering, at an odd number of TRIM. The formula
counts band inversions.

\section{Algorithm 1: the invariant from parities}
\label{sec:goldparity}

Your eigensolver returns \emph{some} orthonormal basis of each degenerate
level, not parity eigenstates. That does not matter. With the permutation
of Eq.~\eqref{eq:parityperm}, form the matrix of $\hat P$ among the lowest
$2N$ spinor bands at a TRIM,
\begin{equation}
\Pi_{mn} = \braket{\psi_m|\hat P|\psi_n}
= \sum_{\GG, s} c^{s\,*}_m(\GG)\; c^{s}_n(-\GG-\GG_0),
\qquad m,n = 1,\dots,2N,
\end{equation}
and diagonalize it. Because $\hat P$ commutes with $\hat H$ and the set of
$2N$ bands is not cut through the middle of a degenerate level, the
eigenvalues come out as exactly $\pm1$, each in Kramers-degenerate
pairs. If $N_-$ of the $2N$ eigenvalues are $-1$, then
\begin{equation}
\delta_i = (-1)^{N_-/2}.
\label{eq:deltafromcount}
\end{equation}

\begin{algorithm}[H]
\caption{$\mathbb{Z}_2$ invariant of gold from parities}
\label{alg:z2}
\begin{algorithmic}[1]
\State Converge the SCF (SOC on); keep $V_{\mathrm{eff,loc}}$ fixed from
here on (Section~\ref{sec:nscf})
\State $N \gets 6$ \Comment{Kramers pairs below the direct gap}
\State \textbf{Direct-gap check:} \Call{SolveBands}{} on a fine mesh and on
the high-symmetry path; verify
$\min_\kk\,(\varepsilon_{2N+1,\kk} - \varepsilon_{2N,\kk}) > 0$
\For{each of the 8 TRIM $\kk_i = \tfrac12(n_1,n_2,n_3)$}
\State $(\varepsilon, C) \gets$ \Call{SolveBands}{$\kk_i$;\
$N_{\mathrm{want}} = 2N{+}2$;\ $\epsilon = 10^{-9}$}
\State \textbf{assert} $\varepsilon_{2N+1} - \varepsilon_{2N} > 10^{-3}$\,Ha
and bands pair up ($|\varepsilon_{2j} - \varepsilon_{2j-1}| < 10^{-8}$\,Ha)
\State $\mathrm{perm}[g] \gets$ index of plane wave
$(-m_1{-}n_1, -m_2{-}n_2, -m_3{-}n_3)$ for each plane wave $g = (m_1,m_2,m_3)$
\State $\Pi_{mn} \gets \sum_{s,g}\, \overline{C[m,s,g]}\; C[n,s,\mathrm{perm}[g]]$
for $m,n \le 2N$
\Comment{$2N\times2N$, Hermitian}
\State $\lambda \gets \mathrm{eigvalsh}(\Pi)$;\quad \textbf{assert} every
$\big||\lambda_j| - 1\big| < 10^{-6}$
\State $N_- \gets \#\{\lambda_j < 0\}$;\quad \textbf{assert} $N_-$ even;\quad
$\delta_i \gets (-1)^{N_-/2}$
\EndFor
\State $\nu_0 \gets 0$ if $\prod_{i=1}^{8}\delta_i = +1$, else $1$
\State $\nu_j \gets 0$ if $\prod_{i:\,n_j = 1}\delta_i = +1$, else $1$,
for $j = 1,2,3$ \Comment{weak indices: four TRIM each}
\State \Return $(\nu_0;\,\nu_1\nu_2\nu_3)$
\end{algorithmic}
\end{algorithm}

The assertions are not decoration. $|\lambda| \ne 1$ means the band set
cuts a degenerate level, or the eigenvectors are not converged, or the
permutation is wrong; an odd $N_-$ means Kramers pairing is broken.

\section{What to expect for gold, and why}

\begin{center}
\begin{tabular}{@{}lllcc@{}}
\toprule
TRIM & Even pairs (eV) & Odd pairs (eV) & $N_-/2$ & $\delta_i$\\
\midrule
$\Gamma$ ($\times 1$) & $-10.06$, $-5.45\,(\times2)$, $-4.24$, $-3.02\,(\times2)$ & none below $+15.97$ & 0 & $+1$\\
$X$ ($\times 3$) & $-7.32$, $-7.03$, $-2.31$, $-2.09$, $-1.05$ & $+0.71$ & 1 & $-1$\\
$L$ ($\times 4$) & $-7.56$, $-5.57$, $-4.38$, $-2.36$, $-1.63$ & $-1.22$ & 1 & $-1$\\
\bottomrule
\end{tabular}
\end{center}
The table lists the six lowest Kramers pairs at each TRIM with their
energies relative to $E_F$ and their computed parities (reference
calculation: LDA, $E_{\mathrm{cut}}=40$\,Ha). The seventh pair is even at
$X$ ($+5.78$\,eV) and $L$ ($+3.13$\,eV) and odd at $\Gamma$.

\begin{equation}
(-1)^{\nu_0} = \underbrace{(+1)}_{\Gamma}\cdot\underbrace{(-1)^3}_{3X}
\cdot\underbrace{(-1)^4}_{4L} = -1
\qquad\Longrightarrow\qquad \nu_0 = 1 .
\end{equation}
Each weak index involves two $X$ and two $L$ points, giving $+1$: the
result is $(\nu_0;\nu_1\nu_2\nu_3) = (1;000)$. Treat the table as a
prediction to be confirmed by your own numbers, not as an input. The
energies depend on the functional and the cutoff at the level of
$0.1$\,eV; the parities do not.

The pattern can be understood with nothing more than
Sections~\ref{sec:nfe} and \ref{sec:symmetry}:
\begin{itemize}[leftmargin=2em]
\item \textbf{$d$ states are even everywhere.} A Bloch sum of an
atom-centered orbital, evaluated at a TRIM, inherits the parity
$(-1)^\ell$ of the orbital. Five $d$ orbitals give five even pairs at
every TRIM.
\item \textbf{At $\Gamma$} the lowest free-electron state is the constant
plane wave $\GG = 0$, which is even ($s$-like). Odd states are far above
($16$\,eV above $E_F$ in the table).
\item \textbf{At a zone-boundary point} ($X$ or $L$) the two lowest free-
electron states are degenerate, and the crystal potential splits them
into the even, $s$-like $\cos$ state and the odd, $p$-like $\sin$ state
of Eq.~\eqref{eq:nfeparity}. Which one is lower? In the two-plane-wave
model that was decided by the sign of one Fourier component of the
potential. In gold something else dominates: \emph{states of equal parity
repel each other}. The even combination is pushed \emph{up} by the five
even $d$ bands lying just below it; the odd one cannot mix with them and
stays put. In gold (as in copper and silver) this repulsion is strong
enough that the odd state ends up \emph{below} the even one at both $X$
and $L$ --- the inverted order of Section~\ref{sec:nfe}. That is the band
inversion. At $L$ the odd level sits $1.2$\,eV below $E_F$ and the
even one $3.1$\,eV above; the gap between them is the ``$L$ gap''
through which the necks of gold's Fermi surface pass.
\end{itemize}
Because there are four $L$ points, their signs cancel in $\nu_0$: the
strong index is decided by $\Gamma$ versus $X$. The ordering at $L$
matters for \emph{where} on a given surface the protected state appears.

\paragraph{The measurable consequence.} Cut the crystal along a (111)
plane. Photoemission on Au(111) shows a free-electron-like \emph{surface
state} centered at the surface-zone center, about $0.5$\,eV below $E_F$,
sitting inside the bulk $L$ gap. It is split into two parabolas of
opposite spin, shifted in momentum (the Rashba splitting, first resolved
by LaShell, McDougall and Jensen in 1996), with the spin locked
perpendicular to the momentum. Yan \emph{et al.}\ showed that this
long-known ``Shockley state'' is the topological surface state demanded by
$\nu_0 = 1$: the bulk invariant computed here from eight parities is why
that surface band must exist and must be spin--momentum locked. (To see
the surface state itself you would need a slab calculation --- a natural
sequel once the bulk code works.)

\section{Algorithm 2: the invariant from Wilson loops}
\label{sec:wilson}

The parity shortcut exists only for inversion-symmetric crystals. The
general method of Sections~\ref{sec:berry}--\ref{sec:z2def} is worth
implementing: it is what you will need for any crystal without inversion,
it gives an independent check on gold, and it lets you \emph{see} the
topology.

\begin{algorithm}[H]
\caption{Wilson-loop $\mathbb{Z}_2$ on one plane (fixed $k_3 \in
\{0,\tfrac12\}$)}
\label{alg:wilson}
\begin{algorithmic}[1]
\For{$i = 0,\dots,N_2$} \Comment{$k_2 = i/(2N_2)$: from $0$ to $\tfrac12$
inclusive}
\State $(\varepsilon, C_j) \gets$ \Call{SolveBands}{} for the lowest $2N$
bands at $\kk_j = (j/N_1,\, k_2,\, k_3)$, $j = 0,\dots,N_1{-}1$
\State $C_{N_1} \gets C_0$ with every Miller index $m_1$ lowered by one
\Comment{Eq.~\eqref{eq:closure}; do not re-solve}
\State $\mathcal W \gets \mathbf 1_{2N}$
\For{$j = 0,\dots,N_1{-}1$}
\State $M \gets$ overlap matrix \eqref{eq:overlapM} between $C_j$ and
$C_{j+1}$
\State $(U, \Sigma, V^\dagger) \gets \mathrm{svd}(M)$;\quad
record $\min\Sigma$;\quad
$\mathcal W \gets \mathcal W\, U V^\dagger$
\EndFor
\State $\theta[i,\cdot] \gets$ phases of the eigenvalues of $\mathcal W$,
sorted
\EndFor
\State $z \gets$ \Call{CountCrossings}{$\theta$}
\Comment{Algorithm~\ref{alg:count}}
\end{algorithmic}
\end{algorithm}

Line 7 replaces each overlap matrix by its nearest unitary matrix
$UV^\dagger$. On a finite mesh $M$ is not quite unitary (its singular
values are slightly below one, by an amount of order $1/N_1^2$); keeping
only the unitary part makes $\mathcal W$ exactly unitary, so its
eigenvalues lie exactly on the unit circle. A singular value near zero
signals that the group of $2N$ bands is not the same group at the two
neighbouring points --- the direct gap has closed, or the mesh is far too
coarse. To compute the overlaps, scatter both coefficient arrays into the
FFT box (or any common integer index) so that equal $\GG$ line up.

Counting crossings ``by eye'' on a plot is how one checks the result;
a program needs a rule. The standard one (Soluyanov and Vanderbilt)
uses, as the reference line, the middle of the \emph{largest gap} between
adjacent centers, and lets it move with $k_2$: whenever the gap center
jumps from one step to the next, count the centers it jumped over.

\begin{algorithm}[H]
\caption{Counting partner switches by following the largest gap}
\label{alg:count}
\begin{algorithmic}[1]
\Function{CountCrossings}{$\theta[i,j]$,\ $i = 0..N_2$,\ $j = 1..2N$}
\Function{GapCenter}{$\vartheta$} \Comment{$\vartheta$: the $2N$ sorted angles of one row}
\State find adjacent angles (cyclically, the last and the first are
adjacent across $\pm\pi$) with the largest separation;\ \Return their
midpoint on the circle
\EndFunction
\State $z \gets 0$;\quad $g \gets$ \Call{GapCenter}{$\theta[0,\cdot]$}
\For{$i = 1,\dots,N_2$}
\State $g' \gets$ \Call{GapCenter}{$\theta[i,\cdot]$}
\State $z \gets z + \#\{\,j : \theta[i,j] \text{ lies on the shorter arc
between } g \text{ and } g'\,\}$
\State $g \gets g'$
\EndFor
\State \Return $z \bmod 2$
\EndFunction
\end{algorithmic}
\end{algorithm}

Use $N_1 \approx 20$--$40$ points along the loop and enough $k_2$ values
(start with 80) that the curves are smooth; if the count changes when
$N_2$ is doubled, the mesh was too coarse. Run it on the planes $k_3 = 0$
and $k_3 = \tfrac12$ and combine with Eq.~\eqref{eq:nu0fromz}.

\paragraph{What the reference calculation finds.} Figure~\ref{fig:wcc}
shows the twelve charge centers of gold on both planes ($N_1 = 24$,
$N_2 = 80$). Because gold has inversion symmetry the picture is symmetric
under $\theta\to-\theta$ at every $k_2$; most centers stay near
$\theta = 0$ (the $d$ electrons, sitting on the atom) and the action is in
the outermost curves.
\begin{itemize}[leftmargin=2em]
\item \textbf{Plane $k_3 = 0$.} At $k_2 = 0$ the loop passes through
$\Gamma$ and an $L$ point ($\delta_\Gamma\delta_L = -1$): five Kramers
pairs of centers sit at $\theta = 0$ and one at $\theta = \pi$, as
Eq.~\eqref{eq:zakparity} demands. At $k_2 = \tfrac12$ it passes through an
$L$ and an $X$ point ($\delta_L\delta_X = +1$): four pairs at $\theta = 0$
and a mirror-image set at $\theta = \pm0.51\pi$ --- no pair at $\pi$. The
pair that started at $\pi$ has come apart and its members have joined
other partners: $z = 1$.
\item \textbf{Plane $k_3 = \tfrac12$.} Both ends pass through one $L$ and
one $X$ point and look identical (four pairs at $0$, a set at
$\pm0.51\pi$). The curves make large excursions in between but return to
the partners they started with: $z = 0$.
\end{itemize}
Hence $\nu_0 = (1+0) \bmod 2 = 1$ and $\nu_3 = 0$, in agreement with the
parities.

Now look at Fig.~\ref{fig:wcc} critically, because gold is a \emph{hard}
case for this method and the figure shows it. The outer curves do not
move smoothly: near $k_2 \approx 0.26$ and $0.43$ (left) and
$k_2\approx 0.14$ and $0.37$ (right) they appear to jump. True charge
centers are continuous functions of $k_2$ as long as the direct gap is
open; the apparent jumps occur where the loop passes close to the points
at which the gap nearly closes (Section~\ref{sec:directgap}: $0.02$\,eV).
There the wavefunctions of bands 12 and 13 exchange character within a
tiny region of $\kk$-space, the centers sweep across the circle within a
correspondingly tiny interval of $k_2$, and the smallest singular value
of the overlap matrices drops to $0.03$ --- far below the healthy range.
Three consequences:
\begin{itemize}[leftmargin=2em]
\item The \emph{end points} $k_2 = 0, \tfrac12$ are unaffected (the gap
there is large) and, through Eq.~\eqref{eq:zakparity}, already determine
$z$ for an inversion-symmetric crystal. This is the solid part of the
result.
\item A crossing count that runs across an unresolved jump is not
trustworthy. With $N_2 = 80$ Algorithm~\ref{alg:count} returned the
values above; with $N_2 = 20$ it returned $z = 1$ on \emph{both} planes,
i.e.\ the wrong $\nu_0$. An invariant that changes when the mesh is
refined was not converged.
\item The cure is local refinement: wherever the gap center of
Algorithm~\ref{alg:count} moves by more than a fraction of the gap width
between two $k_2$ values, or a singular value falls below $\sim0.3$,
insert intermediate $k_2$ points (and increase $N_1$) until the curves
are continuous. For a material with a comfortable direct gap none of this
is needed; for gold it is the difference between a picture and a proof,
and it is the reason the parity method is the primary tool here.
\end{itemize}
(In an inversion-symmetric crystal the largest gap between centers is
degenerate with its mirror image, so the reference line of
Algorithm~\ref{alg:count} may hop between the two; each hop passes an even
number of centers and does not affect $z$.)

\begin{figure}[H]
\centering
\pgfplotsset{wccplot/.style={width=0.5\textwidth, height=0.42\textwidth,
  xmin=0, xmax=0.5, ymin=-1, ymax=1, xlabel={$k_2$},
  ytick={-1,-0.5,0,0.5,1},
  tick label style={font=\small}, label style={font=\small},
  title style={font=\small}}}
\begin{tikzpicture}
\begin{axis}[wccplot, ylabel={$\theta/\pi$}, title={$k_3 = 0$:\ \ $z = 1$}]
\addplot[only marks, mark=*, mark size=0.6pt, red!70!black] table[x index=0,y index=1] {figdata/wcc_k3_0.dat};
\end{axis}
\end{tikzpicture}\hfill
\begin{tikzpicture}
\begin{axis}[wccplot, yticklabels={}, title={$k_3 = \tfrac12$:\ \ $z = 0$}]
\addplot[only marks, mark=*, mark size=0.6pt, red!70!black] table[x index=0,y index=1] {figdata/wcc_k3_half.dat};
\end{axis}
\end{tikzpicture}
\caption{Wannier charge centers (Wilson-loop phases) of the twelve lowest
spinor bands of gold across half the zone, from the reference calculation
($E_{\mathrm{cut}} = 20$\,Ha, fixed self-consistent potential). $\theta$
is periodic: $+\pi$ and $-\pi$ are the same point. Left: the Kramers pair
at $\theta = \pm\pi$ for $k_2 = 0$ has no counterpart at $k_2 = \tfrac12$
--- partners are switched. Right: the two ends match. The apparent jumps
are unresolved rapid motion near the points of smallest direct gap (see
text).}
\label{fig:wcc}
\end{figure}

\section{Sanity checks}

\begin{enumerate}[leftmargin=2em]
\item \textbf{Empty lattice.} With $V = 0$ at $\Gamma$ the lowest state is
even. At $X$ the two lowest spatial states are degenerate ($\times2$ for
spin): $\Pi$ on those four spinor states must have eigenvalues
$\{+1,+1,-1,-1\}$. This tests the permutation \eqref{eq:parityperm}
without any physics.
\item \textbf{The two-plane-wave model, live.} With the single-cosine
potential of Stage~1 (Chapter~\ref{ch:verify}),
$V(\rr) = 2V_0\cos(\GG_1\cdot\rr)$ and $\kk = \tfrac12\GG_1$, the two
lowest levels must have parities $(+,-)$ in ascending order for $V_0<0$
and $(-,+)$ for $V_0>0$, Eq.~\eqref{eq:nfeparity}. A band inversion made
by hand.
\item \textbf{$\Pi^2 = \mathbf 1$ and $\Pi = \Pi^\dagger$} to rounding,
when the block of bands is not cut through a degenerate level.
\item \textbf{$[\hat P, \hat H] = 0$ and $[\Theta,\hat H]=0$} for random
spinors at each TRIM (Section~\ref{sec:symmetry}).
\item \textbf{Count of even states.} At $\Gamma$ the lowest twelve spinor
bands of gold must all be even.
\item \textbf{SOC off.} The parity counts at the TRIM cannot depend on
SOC: there the groups are separated by $4$--$19$\,eV, far more than any
spin--orbit shift, so setting $k^\ell = 0$ must leave every $\delta_i$
unchanged. What changes is the direct gap away from the TRIM
(Section~\ref{sec:directgap}).
\item \textbf{Parities versus Wilson loops.} On each of the lines
$k_2 \in\{0,\tfrac12\}$ of each plane, the Wilson-loop phases must be
pairwise degenerate, symmetric under $\theta\to-\theta$, and satisfy
Eq.~\eqref{eq:zakparity} with the $\delta_i$ from Algorithm~\ref{alg:z2}.
\item \textbf{Robustness.} Change $E_{\mathrm{cut}}$, the functional, or
the lattice constant by a few percent: every $\delta_i$ and every $z$ must
stay the same. A topological invariant that changes with the cutoff is a
bug.
\end{enumerate}

\begin{gpubox}{GPU/JAX notes}
The parity method is eight small eigenproblems plus one gather
(\texttt{C[..., perm]}) and one $12\times12$ \texttt{eigvalsh} each ---
negligible cost, and it batches over the eight TRIM like any other
$\kk$ axis. The Wilson loop is a fixed-potential band calculation on
$N_1 \times (N_2{+}1)$ points per plane, followed by a \texttt{scan} of
small matrix products; the overlaps are one \texttt{einsum} over the
coefficient arrays once both spheres are scattered onto a common index
set (the FFT box works). The $12\times12$ SVDs and the final
non-Hermitian \texttt{eig} are cheapest on the host. Both methods must run
in double precision with tightly converged eigenvectors.
\end{gpubox}

\begin{keybox}{Summary}
The centers of charge of a group of bands are computable from overlaps of
neighbouring wavefunctions (Wilson loops). With time-reversal symmetry
those centers come in Kramers pairs at the edges of the half zone, and the
$\mathbb{Z}_2$ invariant records whether they switch partners in between.
With inversion symmetry the same bit is the product of the parities of the
Kramers pairs at the eight TRIM, Eq.~\eqref{eq:fukane}, obtained by
applying the permutation $c(\GG) \to c(-\GG-2\kk)$ to your coefficient
arrays. For gold both routes give $\nu_0 = 1$: the odd, $p$-like band lies
below the even, $s$-like band at the zone boundary, and the spin-split
surface state of Au(111) is the visible consequence.
\end{keybox}

\chapter{Electrical Conductivity: What DFT Can and Cannot Tell You}
\label{ch:conductivity}
% generated by verification/optics_report.py -- do not edit
\newcommand{\GF}{0.264}
\newcommand{\GFb}{0.276}
\newcommand{\VF}{1.33}
\newcommand{\VFb}{1.31}
\newcommand{\WP}{8.47}
\newcommand{\WPb}{8.55}
\newcommand{\ST}{1.47}
\newcommand{\STb}{1.49}
\newcommand{\NEFF}{0.90}
\newcommand{\TAU}{30}
\newcommand{\MFP}{40}
\newcommand{\LAMTR}{0.14}
\newcommand{\NKFINE}{1661}
\newcommand{\NKCOARSE}{413}
\newcommand{\NTRANS}{199}
\newcommand{\EDGE}{1.9}
\newcommand{\ONSET}{1.5}
\newcommand{\EDGES}{2.1}
\newcommand{\EPSZERO}{9.9}
\newcommand{\NEFFINTER}{1.4}
\newcommand{\RGBLDA}{195,182,177}
\newcommand{\RLDAb}{0.44}
\newcommand{\RLDAg}{0.46}
\newcommand{\RLDAr}{0.65}
\newcommand{\RGBSHa}{243,162,157}
\newcommand{\RSHab}{0.36}
\newcommand{\RSHag}{0.38}
\newcommand{\RSHar}{0.84}
\newcommand{\RGBSHb}{255,180,129}
\newcommand{\RSHbb}{0.26}
\newcommand{\RSHbg}{0.52}
\newcommand{\RSHbr}{0.92}
\newcommand{\RGBDR}{253,253,253}
\newcommand{\RDRb}{0.98}
\newcommand{\RDRg}{0.98}
\newcommand{\RDRr}{0.98}
\newcommand{\RGBCO}{179,184,168}
\newcommand{\RCOb}{0.36}
\newcommand{\RCOg}{0.44}
\newcommand{\RCOr}{0.63}
\newcommand{\RGBETA}{205,183,175}
\newcommand{\RETAb}{0.44}
\newcommand{\RETAg}{0.48}
\newcommand{\RETAr}{0.69}
\definecolor{goldLDA}{RGB}{195,182,177}
\definecolor{goldSHa}{RGB}{243,162,157}
\definecolor{goldSHb}{RGB}{255,180,129}
\definecolor{goldDR}{RGB}{253,253,253}
\newcommand{\RGBEXP}{255,222,163}
\definecolor{goldEXP}{RGB}{255,222,163}
\newcommand{\REXPb}{0.41}
\newcommand{\REXPg}{0.79}
\newcommand{\REXPr}{0.96}
\newcommand{\EDGEEXP}{2.4}
\newcommand{\EPSEXP}{5.7}
\newcommand{\EPSCALC}{8.3}
\newcommand{\EDGENOSOC}{2.04}
\newcommand{\EDGESOC}{1.85}


Gold is one of the best conductors there is: at room temperature its
resistivity is $\rho = 2.2\,\mu\Omega\,$cm ($22$\,n$\Omega$\,m at
$293$\,K), i.e.\ a conductivity
$\sigma = 1/\rho = 4.5\times10^{7}$\,S/m. Can the program you have written
predict this number? The honest answer has two halves, and separating them
is the point of this chapter.
\begin{itemize}[leftmargin=2em]
\item \textbf{Yes:} everything about conduction that is a property of the
\emph{perfect crystal} --- which electrons carry the current, how fast they
move, how strongly they respond to a field --- follows from the bands and
wavefunctions you already have. These combine into one number, the
\emph{plasma frequency} $\omega_p$, and you will compute it.
\item \textbf{Not by itself:} in a perfect, rigid crystal the conductivity
is \emph{infinite}. The measured value is finite only because electrons are
scattered --- at room temperature by the thermal vibrations of the lattice
--- and the scattering rate is not a ground-state property of the frozen
crystal. It \emph{can} be computed from first principles, but that needs
the lattice vibrations and their coupling to the electrons, a substantial
extension of the code that is outlined, not built, in
Section~\ref{sec:tau}.
\end{itemize}
The one new tool introduced here, the velocity operator, is also the key to
the optical properties of the next chapter.

\section{Why a perfect crystal does not resist}

A Bloch state $\psi_{n\kk}$ is a stationary state of the periodic
potential: an electron in it moves through the lattice forever with the
mean velocity (derived in the next section)
\begin{equation}
\mathbf v_{n\kk} = \pdv{\varepsilon_{n\kk}}{\kk}
\label{eq:groupvel}
\end{equation}
--- the group velocity of a wave packet built from states near $\kk$. The
periodic potential of the ions, however strong, does not scatter it; that
is the content of Bloch's theorem. In equilibrium the currents of all
occupied states cancel, because $\varepsilon_{n,-\kk} = \varepsilon_{n\kk}$
implies $\mathbf v_{n,-\kk} = -\mathbf v_{n\kk}$.

Now apply a uniform electric field $\mathbf E$. The semiclassical equation
of motion of a Bloch electron (charge $-1$ in atomic units) is
\begin{equation}
\frac{\dd\kk}{\dd t} = -\mathbf E :
\end{equation}
the force changes the \emph{crystal momentum}, so every electron drifts
through $\kk$-space at the same rate. The occupied region of the Brillouin
zone shifts as a whole, the cancellation between $\kk$ and $-\kk$ is lost,
and a current flows that keeps growing as long as the field is on. Nothing
in the perfect crystal stops it. A finite, steady current requires
processes that knock electrons back toward equilibrium: collisions with
lattice vibrations (phonons), with impurities and defects, and with each
other. Summarize them, crudely but usefully, by a single \emph{relaxation
time} $\tau$: the average time an electron travels between collisions.

\section{The velocity operator}
\label{sec:velocity}

Both this chapter and the next need matrix elements of the electron
velocity between Bloch states. At fixed $\kk$ your program represents the
Hamiltonian as a matrix $H_{\kk}$ acting on the coefficients $c(\GG)$ ---
the Hamiltonian for the periodic functions $u_{n\kk}$. The velocity
operator in the same representation is simply its derivative,
\begin{equation}
\boxed{\;\hat{\mathbf v} = \pdv{H_{\kk}}{\kk}\;}
\qquad\qquad
\mathbf v_{mn}(\kk) = \sum_{\GG\GG' ss'} c^{s\,*}_{m\kk}(\GG)\,
\pdv{H^{ss'}_{\GG\GG'}(\kk)}{\kk}\, c^{s'}_{n\kk}(\GG') .
\label{eq:velop}
\end{equation}
\emph{Why.} Differentiate the eigenvalue $\varepsilon_{n\kk} =
\braket{u_{n\kk}|H_\kk|u_{n\kk}}$ with respect to $\kk$. Three terms
appear; the two containing $\partial u/\partial\kk$ add up to
$\varepsilon_{n\kk}\,\partial_\kk\braket{u_{n\kk}|u_{n\kk}} = 0$ because
the state stays normalized. What remains is the
\textbf{Hellmann--Feynman theorem},
\begin{equation}
\pdv{\varepsilon_{n\kk}}{\kk} = \Braket{u_{n\kk}|\pdv{H_\kk}{\kk}|u_{n\kk}}
= \mathbf v_{nn}(\kk),
\label{eq:hf}
\end{equation}
which is Eq.~\eqref{eq:groupvel}: the diagonal matrix element of
$\partial H_\kk/\partial\kk$ is the group velocity. The off-diagonal
elements $\mathbf v_{mn}$, $m\ne n$, govern transitions between bands
(Chapter~\ref{ch:colour}).

\paragraph{Term by term.} Go through the Hamiltonian of
Chapter~\ref{ch:hamiltonian} and ask where $\kk$ appears.
\begin{itemize}[leftmargin=2em]
\item \emph{Kinetic energy} $\tfrac12|\kk+\GG|^2$: derivative $\kk+\GG$.
This is the momentum, as for a free particle.
\item \emph{Local potential} $V_{\mathrm{eff,loc}}(\GG-\GG')$: no $\kk$
dependence, no contribution.
\item \emph{Nonlocal pseudopotential}: the projectors
$\tilde\beta_p(\kk+\GG)$ of Eq.~\eqref{eq:betaG} depend on $\kk$, so this
term \emph{does} contribute. (Physically: a nonlocal potential does not
commute with the position operator, so velocity is no longer momentum
alone. Leaving this term out is the classic error in pseudopotential
optics codes.)
\end{itemize}
Writing $\mathbf q = \kk+\GG$ and $\partial_\alpha\tilde\beta_p \equiv
\partial\tilde\beta_p(\mathbf q)/\partial q_\alpha$ for the Cartesian
directions $\alpha = x,y,z$, and with the projections
$P_{n} = \sum_\GG\tilde\beta^*(\mathbf q)\,c_n(\GG)$ of
Algorithm~\ref{alg:applyh} and their analogues
$\partial_\alpha P_{n} = \sum_\GG[\partial_\alpha\tilde\beta(\mathbf q)]^*
c_n(\GG)$ (both are vectors over the composite index: projector $\times$
spin),
\begin{equation}
v^{\alpha}_{mn} = \sum_{\GG s} c^{s\,*}_{m}(\GG)\,(\kk+\GG)_\alpha\,
c^{s}_{n}(\GG)
\;+\; (\partial_\alpha P_m)^\dagger D\, P_n
\;+\; P_m^\dagger D\, (\partial_\alpha P_n).
\label{eq:vmn}
\end{equation}
The derivative of a projector is the derivative of an explicit function,
$\tilde\beta_p(\mathbf q) \propto \tilde p_{i\ell}(|\mathbf q|)\,
Y_{\ell m}(\hat{\mathbf q})$: obtain it with automatic differentiation
(\texttt{jax.jacfwd} of the function that returns the projector table for
given $\mathbf q$), or with a central difference
$[\tilde\beta(\mathbf q + \delta\,\hat{\mathbf e}_\alpha)
- \tilde\beta(\mathbf q - \delta\,\hat{\mathbf e}_\alpha)]/2\delta$,
$\delta = 10^{-4}$, which is accurate to $10^{-8}$ for these smooth
functions.

\begin{algorithm}[H]
\caption{Velocity matrix elements at one $\kk$-point}
\label{alg:velocity}
\begin{algorithmic}[1]
\Function{Velocity}{$C[n,s,\GG]$;\ basis}
\State $P[n,(s,p)] \gets \sum_\GG \tilde\beta_p^*(\mathbf q)\,C[n,s,\GG]$;\quad
$DP \gets P\,D^{\mathsf T}$ \Comment{as in \textsc{ApplyH}}
\For{$\alpha = x, y, z$}
\State $\partial_\alpha\tilde\beta_p(\mathbf q) \gets$ derivative of the
projector table with respect to $q_\alpha$
\State $\partial P[n,(s,p)] \gets \sum_\GG
[\partial_\alpha\tilde\beta_p(\mathbf q)]^*\,C[n,s,\GG]$
\State $v^\alpha_{mn} \gets \sum_{\GG s}\overline{C[m,s,\GG]}\; q_\alpha\, C[n,s,\GG]$
\Comment{kinetic part}
\State $v^\alpha_{mn} \mathrel{+}= \sum_{a}\overline{\partial P[m,a]}\,DP[n,a]
+ \sum_{a}\overline{P[m,a]}\,(\partial P\,D^{\mathsf T})[n,a]$
\Comment{nonlocal part}
\EndFor
\State \Return $v^\alpha_{mn}$ \Comment{three Hermitian $N_b\times N_b$ matrices}
\EndFunction
\end{algorithmic}
\end{algorithm}

\paragraph{Degenerate bands.} Every band of gold is a Kramers pair, and
within a degenerate pair the eigensolver returns an arbitrary basis. The
group velocity of the pair is therefore not an individual diagonal element
but the $2\times2$ block of $v^\alpha$ within the pair. For gold the
combined symmetry $\hat P\Theta$ (Section~\ref{sec:symmetry}) forces that
block to be a multiple of the unit matrix: both partners move with the
same velocity, and the block is the same in any basis. Use
$\tfrac12\mathrm{tr}$ of the block for the velocity and
$\tfrac12\sum_{m,n\in\text{pair}}|v^\alpha_{mn}|^2$ for its square.

\paragraph{Test.} Equation~\eqref{eq:hf} is a complete check of
Algorithm~\ref{alg:velocity}, nonlocal term included: solve at
$\kk \pm h\,\hat{\mathbf e}_\alpha$ ($h = 10^{-4}$\,Bohr$^{-1}$) and compare
$[\varepsilon_n(\kk+h\hat{\mathbf e}_\alpha) -
\varepsilon_n(\kk-h\hat{\mathbf e}_\alpha)]/2h$ with the pair-averaged
$v^\alpha_{nn}$. In the reference calculation they agree to better
than $10^{-6}$\,a.u.\ for all bands at a generic $\kk$, the matrices are
Hermitian to rounding error, and the off-diagonal element inside each
Kramers pair is below $10^{-8}$. Without the nonlocal term the
discrepancy is $0.1$--$0.3$\,a.u.\ for every band --- as large as the
velocities themselves.

\section{Conductivity in the relaxation-time approximation}

In a steady field, an electron accumulates the shift $\delta\kk = -\mathbf
E\,\tau$ before it is scattered. The occupation of state $n\kk$ is then
that of the equilibrium distribution $f^0$ at the point the electron
started from,
\begin{equation}
f_{n\kk} = f^0(\varepsilon_{n,\kk+\mathbf E\tau})
\approx f^0(\varepsilon_{n\kk})
+ \pdv{f^0}{\varepsilon}\;\mathbf v_{n\kk}\cdot\mathbf E\,\tau .
\end{equation}
The current density is the sum of charge ($-1$) times velocity over all
occupied states, per unit volume; with our conventions (one electron per
spinor state, weights $\sum_\kk w_\kk = 1$, cell volume $\Omega$) that is
$\mathbf j = -\frac{1}{\Omega}\sum_{n\kk}w_\kk f_{n\kk}\mathbf v_{n\kk}$.
The equilibrium part cancels, and what is left is linear in the field,
$j_\alpha = \sum_\beta\sigma_{\alpha\beta}E_\beta$, with
\begin{equation}
\sigma_{\alpha\beta} = \frac{\tau}{\Omega}\sum_{n\kk} w_\kk
\Big(-\pdv{f^0}{\varepsilon}\Big)_{\varepsilon_{n\kk}}
v^\alpha_{n\kk}\,v^\beta_{n\kk} .
\label{eq:sigmatensor}
\end{equation}
The factor $-\partial f^0/\partial\varepsilon$ is a narrow peak at the
Fermi level (a $\delta$-function at zero temperature): \emph{only
electrons at the Fermi surface conduct}. In a cubic crystal the tensor is
isotropic, $\sigma_{\alpha\beta} = \sigma\,\delta_{\alpha\beta}$, and it is
customary to write
\begin{equation}
\boxed{\;
\sigma = \frac{\omega_p^2\,\tau}{4\pi},
\qquad
\omega_p^2 = \frac{4\pi}{\Omega}\sum_{n\kk} w_\kk\,
\delta(\varepsilon_{n\kk}-\mu)\;\frac{|\mathbf v_{n\kk}|^2}{3}\;}
\label{eq:plasma}
\end{equation}
(atomic units). $\omega_p$ is the \textbf{Drude plasma frequency}. It
packages everything the band structure contributes to conduction, and it
is a ground-state quantity your code can deliver.

\paragraph{Sanity check: free electrons.} For a free-electron gas of density
$n$, with $g(\varepsilon_F) = 3n/2\varepsilon_F$ states per unit energy
and volume and $v_F^2 = 2\varepsilon_F$, Eq.~\eqref{eq:plasma} gives
$\omega_p^2 = \tfrac{4\pi}{3}\,g(\varepsilon_F)\,v_F^2 = 4\pi n$ and
$\sigma = n\tau$ --- in SI units the Drude formula $\sigma = ne^2\tau/m$
of elementary solid-state physics. For gold with one $6s$ electron per
atom, $n = 1/\Omega$: $\hbar\omega_p = 9.02$\,eV and
$v_F = (3\pi^2 n)^{1/3} = 1.39\times10^6$\,m/s. The $d$ electrons, whose
bands lie entirely below $E_F$, do not conduct.

\paragraph{Evaluating the Fermi-surface sum.} The $\delta$-function is
replaced by a normalized Gaussian of width $s$,
$\delta_s(x) = \ee^{-x^2/s^2}/(s\sqrt\pi)$, and the sum is taken over a
mesh much denser than the SCF mesh, using \textsc{SolveBands} in the fixed
self-consistent potential (Section~\ref{sec:nscf}). The integrand is
invariant under the 48 operations of the cubic point group (because it
contains $|\mathbf v|^2$, not individual components), so it suffices to
visit the irreducible wedge: group the mesh points into sets related by
the 48 signed permutations of the Cartesian components, keep one point of
each set, and give it the weight (size of the set)$/N_\kk$. For a
$24^3$ mesh that is 413 points instead of 13\,824.

\begin{algorithm}[H]
\caption{Drude weight, density of states and Fermi velocity}
\label{alg:drude}
\begin{algorithmic}[1]
\State irreducible $\kk$-points and weights $w_\kk$ of an $N^3$
$\Gamma$-centered mesh
\For{each $\kk$}
\State $(\varepsilon, C) \gets$ \Call{SolveBands}{$\kk$};\quad
$v^\alpha_{mn} \gets$ \Call{Velocity}{$C$}
\State $V_n \gets \tfrac13\cdot\tfrac12\sum_\alpha\sum_{m,m'\in
\text{pair}(n)}|v^\alpha_{mm'}|^2$ \Comment{$|\mathbf v_{n\kk}|^2/3$}
\State store $\varepsilon_{n\kk}$, $V_{n\kk}$, and
$\tfrac13\sum_\alpha|v^\alpha_{mn}|^2$ for all pairs $m,n$
\Comment{the last for Chapter~\ref{ch:colour}}
\EndFor
\State $g(\mu) \gets \sum_{n\kk}w_\kk\,\delta_s(\varepsilon_{n\kk}-\mu)$
\Comment{states per Ha per cell}
\State $\omega_p^2 \gets \dfrac{4\pi}{\Omega}\sum_{n\kk}w_\kk\,
\delta_s(\varepsilon_{n\kk}-\mu)\,V_{n\kk}$;\quad
$\langle v_F^2\rangle \gets 3\sum w\,\delta_s V\big/\sum w\,\delta_s$
\end{algorithmic}
\end{algorithm}

\section{Results for gold}

Reference calculation: fixed potential of Case~B (Chapter~\ref{ch:verify};
LDA, spin--orbit coupling on), $\Gamma$-centered meshes reduced to the
irreducible wedge.
\begin{center}
\begin{tabular}{@{}lcccc@{}}
\toprule
 & \multicolumn{2}{c}{this calculation} & free electrons & experiment\\
 & $24^3$ & $40^3$ & ($6s^1$) & \\
\midrule
$g(E_F)$ (states/eV/atom) & \GF & \GFb & $0.27$ & $\approx 0.3$\\
r.m.s.\ Fermi velocity ($10^6$\,m/s) & \VF & \VFb & $1.39$ & $\approx 1.4$\\
$\hbar\omega_p$ (eV) & \WP & \WPb & $9.02$ & $8.45$--$9.0$\\
$\sigma/\tau$ ($10^{6}$\,S\,m$^{-1}$\,fs$^{-1}$) & \ST & \STb & $1.66$ & ---\\
\bottomrule
\end{tabular}
\end{center}
(Gaussian width $s = 0.005$\,Ha. On the $40^3$ mesh, $s = 0.01$\,Ha
raises $\hbar\omega_p$ by $0.2$\,eV and $s = 0.003$\,Ha lowers it by
$0.5$\,eV: so narrow a Gaussian samples the Fermi surface unevenly on a
mesh of this density. Treat the numbers as accurate to a few tenths of an
eV. Experimental values:
$g(E_F)$ from the electronic specific heat coefficient
$\gamma = 0.69$\,mJ\,mol$^{-1}$K$^{-2}$, which includes an enhancement
by the electron--phonon interaction of the order of 10\%; $\omega_p$
from Drude fits to optical data, $8.45$\,eV (Olmon \emph{et al.}, 2012)
to $9.0$\,eV (Johnson and Christy, 1972).) Three
things to take from the table. First, gold conducts like a free-electron
metal with close to one electron per atom: the effective number of
carriers, $\omega_p^2\Omega/4\pi$, comes out as \NEFF{}. The single band
crossing $E_F$ in Fig.~\ref{fig:bands} is the $6s$ band. Second,
Kohn--Sham theory is on firm ground here: unlike the position of the $d$
bands, the Fermi surface of a simple metal is described well. Third, the
band structure alone fixes the ratio $\sigma/\tau$ --- but not $\sigma$.

\paragraph{Unit conversions.} One atomic unit of conductivity is
$e^2/(a_0\hbar) = 4.5998\times10^{6}$\,S/m; of time,
$\hbar/\mathrm{Ha} = 0.024189$\,fs; of velocity,
$2.1877\times10^{6}$\,m/s. Example: $\hbar\omega_p = 8.5$\,eV is
$\omega_p = 0.3124$\,a.u.; with $\tau = 30$\,fs $= 1240$\,a.u.,
$\sigma = \omega_p^2\tau/4\pi = 9.63$\,a.u.\ $= 4.4\times10^7$\,S/m.

\section{The missing ingredient: the relaxation time}
\label{sec:tau}

\paragraph{From experiment.} Turn the argument around: the measured
room-temperature resistivity and the computed plasma frequency give
\begin{equation}
\tau = \frac{4\pi\sigma}{\omega_p^2} \approx \TAU\,\mathrm{fs},
\qquad
\ell = v_F\,\tau \approx \MFP\,\mathrm{nm} .
\end{equation}
An electron at the Fermi surface of gold travels about a hundred lattice
constants between collisions: the picture of nearly free Bloch waves,
occasionally scattered, is self-consistent. (This is also why thin gold
films and nanowires conduct worse than bulk gold: once the sample is not
much thicker than $\ell$, surface scattering matters.)

\paragraph{From first principles (outline only).} At room temperature
$\tau$ in a pure metal is set by electron--phonon scattering. Above the
Debye temperature (about 165\,K for gold) the result of the full theory is
remarkably simple,
\begin{equation}
\frac{1}{\tau} = 2\pi\,\lambda_{\mathrm{tr}}\,k_BT
\qquad\Longrightarrow\qquad
\rho(T) = \frac{4\pi}{\omega_p^2}\cdot 2\pi\,\lambda_{\mathrm{tr}}\,k_BT ,
\label{eq:rhoT}
\end{equation}
(atomic units; $k_BT = 9.34\times10^{-4}$\,Ha at 295\,K): the resistivity
is linear in temperature, as observed, with a slope fixed by $\omega_p$
and one dimensionless number, the \emph{transport electron--phonon
coupling constant} $\lambda_{\mathrm{tr}}$. Inserting the measured
resistivity and the computed $\omega_p$ gives
$\lambda_{\mathrm{tr}}\approx\LAMTR$ --- a small value: gold is a good
conductor because its electrons couple weakly to its lattice vibrations
(the same reason it is not a superconductor at any accessible
temperature). This is exactly how ``empirical'' coupling constants are
extracted in the literature: Savrasov and Savrasov (1996, Eq.~23 there)
use $\lambda_{\mathrm{tr}} = c_1\omega_p^2/8\pi^2k_B$ with $c_1$ the
measured slope of $\rho(T)$, which is Eq.~\eqref{eq:rhoT} solved for
$\lambda_{\mathrm{tr}}$, and find $0.12$ for copper (and values above 1
for good superconductors such as lead and niobium).

Computing $\lambda_{\mathrm{tr}}$ is a research-grade but well-defined DFT
calculation. It requires, beyond this document: (i) forces on atoms
(derivatives of the total energy with respect to atomic positions --- in
JAX, \texttt{grad} of your energy function); (ii) the phonon frequencies
and displacement patterns, from the forces in supercells with one atom
displaced, or from density-functional perturbation theory; (iii) the
change of the self-consistent potential caused by each phonon, and its
matrix elements between Bloch states at the Fermi surface; (iv) a double
Fermi-surface integral of those matrix elements weighted by the change of
velocity in each scattering event. Every step uses machinery you now
understand --- supercells, fixed-potential band calculations, velocity
matrix elements, Fermi-surface sums --- but none of it has been carried
out for this document, and Eq.~\eqref{eq:rhoT} is quoted from the
literature, not derived. That the programme works is established: a
published first-principles Boltzmann calculation for gold (Gunst, Blom
and Stokbro, 2018) obtains $27.0$\,n$\Omega$\,m at 300\,K, some $20$--$30\%$ above the
measured room-temperature value ($\approx 22$\,n$\Omega$\,m), and reproduces the ordering Au $>$ Cu $>$ Ag. The residual resistivity at low temperature is different again:
it is caused by impurities and defects and is a property of the sample,
not of gold.

\begin{keybox}{Summary}
The DC conductivity factorizes, $\sigma = \omega_p^2\tau/4\pi$. The plasma
frequency $\omega_p$ is a Fermi-surface average of squared group
velocities, computed from $\hat{\mathbf v} = \partial H_\kk/\partial\kk$
(kinetic \emph{and} nonlocal-pseudopotential terms) in a fixed-potential
calculation on a dense mesh; for gold it is that of one nearly free
electron per atom. The relaxation time $\tau$ is not a property of the
perfect frozen crystal: take it from experiment, or compute the
electron--phonon coupling --- a separate project.
\end{keybox}

\begin{gpubox}{GPU/JAX notes}
The $\kk$-points of the dense mesh are independent: this is
\textsc{SolveBands} batched over $\kk$ (\texttt{vmap}, in chunks), followed
by three \texttt{einsum}s per $\kk$ for the velocity matrices. The
projector derivative is one \texttt{jax.jacfwd} of a function you already
have; because it is evaluated once per $\kk$-point, its cost is
irrelevant. Store $\varepsilon_{n\kk}$ and the real arrays
$\tfrac13\sum_\alpha|v^\alpha_{mn}|^2$ --- a few megabytes --- and do all
of the analysis of this chapter and the next on the host.
\end{gpubox}

\chapter{The Colour of Gold}
\label{ch:colour}

Chapter~\ref{ch:intro} promised that the program would explain why gold is
yellow. This chapter keeps the promise as far as an honest calculation
can: it derives, implements and tests the complete chain
\[
\text{bands and wavefunctions}
\;\to\; \varepsilon(\omega)
\;\to\; R(\omega)
\;\to\; \text{colour seen by a human eye},
\]
and then looks critically at what comes out. Every step is a
post-processing of the fixed-potential calculation of
Chapter~\ref{ch:conductivity}: no new self-consistent calculation is
needed.

\section{What ``colour'' means physically}

A thick piece of metal transmits nothing; what you see is the light it
reflects. Its colour is therefore determined by the
\textbf{reflectivity} $R(\omega)$, the fraction of the incident intensity
that is reflected at each photon energy $\hbar\omega$. Visible light spans
$\hbar\omega = 1.6$\,eV (red, $\lambda = 780$\,nm) to $3.3$\,eV (violet,
$380$\,nm); the conversion is $\lambda[\mathrm{nm}] =
1239.84/\hbar\omega[\mathrm{eV}]$. A metal that reflects all visible
frequencies equally is white-grey (silver, aluminium). Gold reflects red,
orange and yellow light almost completely and blue light poorly; the eye
reads that mixture as ``golden yellow''. Copper cuts off at a lower
energy still and looks reddish.

The reflectivity in turn is fixed by a single complex function, the
\textbf{dielectric function} $\varepsilon(\omega) = \varepsilon_1(\omega)
+ \ii\varepsilon_2(\omega)$, which describes how the electrons respond to
an oscillating electric field. So the task is to compute
$\varepsilon(\omega)$.

\section{The dielectric function from the band structure}

An electromagnetic wave of frequency $\omega$ can do two things to the
electrons of a metal.

\paragraph{Intraband response: the Drude term.} It can accelerate the
electrons at the Fermi surface, exactly as a static field does in
Chapter~\ref{ch:conductivity}. Repeating that derivation with
$\mathbf E\propto\ee^{-\ii\omega t}$ gives the conductivity
$\sigma(\omega) = \sigma(0)/(1-\ii\omega\tau)$, and since a conductivity
contributes $4\pi\ii\sigma/\omega$ to the dielectric function,
\begin{equation}
\varepsilon_{\mathrm{intra}}(\omega) = -\frac{\omega_p^2}{\omega\,(\omega +
\ii\gamma)},
\qquad \gamma = 1/\tau ,
\label{eq:drude}
\end{equation}
with the plasma frequency of Eq.~\eqref{eq:plasma}. For
$\gamma\ll\omega\ll\omega_p$ this is large and negative: the free
electrons move so as to cancel the field inside the metal, and the wave is
reflected. This is why metals are shiny.

\paragraph{Interband response: absorption.} It can lift an electron from
an occupied state $n\kk$ to an empty state $m\kk$ of another band. Light
carries negligible momentum on the scale of the Brillouin zone, so $\kk$
does not change: transitions are \emph{vertical} in a band-structure plot.
The photon is absorbed if its energy matches,
$\hbar\omega = \varepsilon_{m\kk}-\varepsilon_{n\kk}\equiv\omega_{mn}$,
with a rate given by Fermi's golden rule. The perturbation is the coupling
of the field to the electron's velocity, so the rate is proportional to
$|\mathbf v_{mn}|^2$, the off-diagonal velocity matrix elements of
Eq.~\eqref{eq:vmn}. Summing the absorbed power over all transitions and
dividing by the energy flux of the wave gives the absorptive part of the
dielectric function (cubic crystal, atomic units, one electron per spinor
state):
\begin{equation}
\boxed{\;
\varepsilon_2^{\mathrm{inter}}(\omega) = \frac{4\pi^2}{\Omega}
\sum_{\kk} w_\kk \sum_{n,m}\big(f_{n\kk}-f_{m\kk}\big)\,
\frac{|\mathbf v_{mn}(\kk)|^2}{3\,\omega_{mn}^2}\;
\delta\big(\omega - \omega_{mn}\big),
\qquad \omega_{mn} > 0 .\;}
\label{eq:eps2}
\end{equation}
Read it as a recipe: for every $\kk$ and every pair (occupied $n$, empty
$m$), put a spike at the transition energy $\omega_{mn}$ with strength
$|\mathbf v_{mn}|^2/3\omega_{mn}^2$. The factor $f_n - f_m$ enforces the
Pauli principle (from a filled to an empty state); $\tfrac13|\mathbf
v|^2$ is the average over the three polarization directions, which are
equivalent in a cubic crystal.

\paragraph{The real part.} Causality ties the real part of a response
function to its imaginary part (the Kramers--Kronig relation,
$\varepsilon_1(\omega) = 1 + \tfrac{2}{\pi}\,\mathcal P\!\!\int_0^\infty
\frac{\omega'\varepsilon_2(\omega')}{\omega'^2-\omega^2}\,\dd\omega'$).
Inserting Eq.~\eqref{eq:eps2} does the integral analytically and gives
both parts in one complex formula, which is what you should program:
\begin{equation}
\varepsilon_{\mathrm{inter}}(\omega) = 1 + \frac{8\pi}{\Omega}
\sum_\kk w_\kk \sum_{\substack{n,m\\ \omega_{mn}>0}}
\big(f_{n\kk}-f_{m\kk}\big)\,
\frac{|\mathbf v_{mn}(\kk)|^2/3}
{\omega_{mn}\,\big[\omega_{mn}^2 - (\omega + \ii\eta)^2\big]} .
\label{eq:epsinter}
\end{equation}
The small positive $\eta$ turns each $\delta$-spike of
Eq.~\eqref{eq:eps2} into a Lorentzian of width $\eta$ (check: for
$\eta\to0$, $\mathrm{Im}\,[\omega_{mn}^2-(\omega+\ii\eta)^2]^{-1} \to
\pi\,\delta(\omega-\omega_{mn})/2\omega_{mn}$, which reproduces
Eq.~\eqref{eq:eps2}). It represents the finite lifetime of the excited
states and, in practice, smooths the result of a finite $\kk$-mesh. The
total dielectric function is
\begin{equation}
\varepsilon(\omega) = \varepsilon_{\mathrm{inter}}(\omega)
+ \varepsilon_{\mathrm{intra}}(\omega) .
\label{eq:epstot}
\end{equation}

\paragraph{What has been approximated.} Equations
\eqref{eq:drude}--\eqref{eq:epstot} are the \emph{independent-particle
approximation}: (i) Kohn--Sham eigenvalue differences are used as
excitation energies --- an assumption, not a theorem (see the
interpretation warnings of Chapter~\ref{ch:hkks}); (ii) the electron and
the hole it leaves behind are assumed not to interact, which is a good
approximation in a metal, where they are screened; (iii) the microscopic
fields that vary on the scale of the unit cell (``local fields'') are
neglected; (iv) only as many empty bands as were computed are included,
which slightly underestimates $\varepsilon_1$. For the colour of a noble
metal (i) is by far the most important, as we shall see.

\section{From the dielectric function to the reflectivity}

Maxwell's equations for a wave hitting a flat surface at normal incidence
give the reflected fraction in terms of the complex refractive index
$\tilde n = \sqrt{\varepsilon}$:
\begin{equation}
R(\omega) = \left|\frac{\sqrt{\varepsilon(\omega)} - 1}
{\sqrt{\varepsilon(\omega)} + 1}\right|^2 .
\label{eq:reflectivity}
\end{equation}
Two limits explain everything that follows. If $\varepsilon$ is real and
negative, $\sqrt\varepsilon$ is purely imaginary and $R = 1$: total
reflection, the Drude regime. If absorption sets in
($\varepsilon_2$ no longer small compared with $|\varepsilon_1|$), $R$
drops. \emph{The colour of a noble metal is decided by the photon energy
at which interband absorption begins}: below that threshold the metal is
a near-perfect mirror, above it a mediocre one. In gold the threshold is
the energy needed to lift an electron from the top of the $5d$ bands to
the Fermi level (Fig.~\ref{fig:bands}). Experimentally it lies at about
$2.4$\,eV ($\approx 520$\,nm, blue-green), so blue light is absorbed (the measured reflectivity falls through $0.7$ at
$2.33$\,eV for gold, $2.16$\,eV for copper and $3.71$\,eV for silver,
interpolating between the tabulated points of Johnson and Christy). In silver, whose $4d$ bands lie deeper, it
is at about $3.8$\,eV, in the ultraviolet --- so silver reflects the whole
visible spectrum and is white. The difference between the two is
relativity (Chapter~\ref{ch:soc}): the relativistic contraction of the
$6s$ shell and expansion of the $5d$ shell bring gold's $d$ bands close
to the Fermi level.

\section{From the reflectivity to a colour}

Human colour vision has three kinds of receptors. Any spectrum of light
$S(\lambda)$ is perceived through just three numbers, the
\emph{tristimulus values}
\begin{equation}
X = \int S(\lambda)\,\bar x(\lambda)\,\dd\lambda,\qquad
Y = \int S(\lambda)\,\bar y(\lambda)\,\dd\lambda,\qquad
Z = \int S(\lambda)\,\bar z(\lambda)\,\dd\lambda ,
\end{equation}
where $\bar x,\bar y,\bar z$ are the standardized \emph{colour-matching
functions} of the CIE 1931 standard observer ($Y$ measures brightness).
For a reflecting object under an illuminant of spectrum $I(\lambda)$, the
light reaching the eye is $S = I\,R$. Four ingredients make this
computable in a few lines.

\paragraph{1. Colour-matching functions.} They are tabulated, but an
accurate analytic fit exists (Wyman, Sloan and Shirley, 2013). With
$\lambda$ in nm and the asymmetric Gaussian
$g(\lambda;\mu,\sigma_1,\sigma_2) =
\exp\!\big[-\tfrac12(\lambda-\mu)^2/\sigma^2\big]$, $\sigma = \sigma_1$
for $\lambda<\mu$ and $\sigma_2$ otherwise:
\begin{equation}
\begin{aligned}
\bar x(\lambda) &= 1.056\,g(\lambda;599.8,37.9,31.0)
+ 0.362\,g(\lambda;442.0,16.0,26.7) - 0.065\,g(\lambda;501.1,20.4,26.2),\\
\bar y(\lambda) &= 0.821\,g(\lambda;568.8,46.9,40.5)
+ 0.286\,g(\lambda;530.9,16.3,31.1),\\
\bar z(\lambda) &= 1.217\,g(\lambda;437.0,11.8,36.0)
+ 0.681\,g(\lambda;459.0,26.0,13.8).
\end{aligned}
\label{eq:cmf}
\end{equation}
(Checked against the CIE table: the largest deviations are $0.015$,
$0.007$ and $0.024$ for functions that peak at $1.06$, $1.00$ and
$1.78$.)

\paragraph{2. Illuminant.} Standard daylight (CIE D65) is approximated here
by a black body at $T = 6500$\,K (the true D65 spectrum differs, mostly
in the violet; tabulated D65 may be substituted, but all
reference colours in this document assume the black body), $I(\lambda)\propto\lambda^{-5}/
[\exp(c_2/\lambda T)-1]$ with $c_2 = 1.438777\times10^{7}$\,nm\,K.

\paragraph{3. Normalization.} Divide by the values a perfect mirror
($R = 1$) would give, so that white stays white whatever the illuminant,
and express the result relative to the white point of the display
standard:
\begin{equation}
X = 0.95047\,\frac{\int I R\,\bar x\,\dd\lambda}{\int I\,\bar x\,\dd\lambda},
\qquad
Y = \frac{\int I R\,\bar y\,\dd\lambda}{\int I\,\bar y\,\dd\lambda},
\qquad
Z = 1.08883\,\frac{\int I R\,\bar z\,\dd\lambda}{\int I\,\bar z\,\dd\lambda}.
\label{eq:xyz}
\end{equation}
The integrals run from 380 to 780\,nm; a sum over 1\,nm steps suffices.
Scaling each channel separately in this way is the simplest (``XYZ
scaling'') form of \emph{chromatic adaptation}; colour scientists use
more refined transforms (e.g.\ Bradford), which shift the final sRGB
values by a few units but change nothing in the physics.

\paragraph{4. Display colour (sRGB).} Screens use the sRGB standard. The
linear red, green and blue intensities are
\begin{equation}
\begin{pmatrix} r\\ g\\ b\end{pmatrix} =
\begin{pmatrix}
\phantom{-}3.2406 & -1.5372 & -0.4986\\
-0.9689 & \phantom{-}1.8758 & \phantom{-}0.0415\\
\phantom{-}0.0557 & -0.2040 & \phantom{-}1.0570
\end{pmatrix}
\begin{pmatrix} X\\ Y\\ Z\end{pmatrix},
\label{eq:srgbmatrix}
\end{equation}
each clipped to $[0,1]$ and then encoded with the sRGB transfer function
$c \mapsto 12.92\,c$ for $c\le0.0031308$ and $1.055\,c^{1/2.4}-0.055$
otherwise; multiply by 255 and round to obtain the familiar 8-bit triple.
Two unit tests: $R\equiv1$ must give $(255,255,255)$, and $R\equiv0.5$
must give $(188,188,188)$.

\begin{algorithm}[H]
\caption{Dielectric function, reflectivity and colour}
\label{alg:colour}
\begin{algorithmic}[1]
\State \textbf{input:} $\varepsilon_{n\kk}$, $w_\kk$,
$V_{mn}(\kk) = \tfrac13\sum_\alpha|v^\alpha_{mn}|^2$ and $\omega_p$ from
Algorithm~\ref{alg:drude};\ $\mu$;\ $\eta$, $\gamma$
\State $f_{n\kk}\gets 1$ if $\varepsilon_{n\kk}<\mu$ else $0$
\Comment{sharp occupations: no smearing here}
\State collect all transitions with $\omega_{mn} =
\varepsilon_{m\kk}-\varepsilon_{n\kk} > 10^{-5}$\,Ha and $f_n - f_m > 0$:
energies $\omega_t$, strengths $S_t = w_\kk\,(f_n-f_m)\,V_{mn}(\kk)$
\For{each $\omega$ on a grid, $0 < \hbar\omega \le 8$\,eV, step $0.01$\,eV}
\State $\varepsilon(\omega) \gets 1 + \dfrac{8\pi}{\Omega}\sum_t
\dfrac{S_t}{\omega_t\,[\omega_t^2-(\omega+\ii\eta)^2]}
\;-\;\dfrac{\omega_p^2}{\omega(\omega+\ii\gamma)}$
\State $R(\omega)\gets\big|(\sqrt{\varepsilon}-1)/(\sqrt{\varepsilon}+1)\big|^2$
\EndFor
\State interpolate $R$ to $\lambda = 380,381,\dots,780$\,nm
($\hbar\omega = 1239.84/\lambda$ eV)
\State $X,Y,Z$ from Eqs.~\eqref{eq:cmf}--\eqref{eq:xyz};\quad sRGB from
Eq.~\eqref{eq:srgbmatrix} and the transfer function
\end{algorithmic}
\end{algorithm}

The transitions within a Kramers pair have $\omega_{mn} = 0$ and are
excluded by the threshold in line 3; they are the intraband term. The
sum over $\kk$ uses the irreducible wedge with its weights, as in
Chapter~\ref{ch:conductivity}.

\section{Results: the computed spectrum of gold}
\label{sec:colourresults}

The reference calculation uses the fixed self-consistent potential of
Case~B (LDA, spin--orbit coupling on), the irreducible wedge of a
$\Gamma$-centered $40^3$ mesh (\NKFINE{} points, 22 spinor bands, about
\NTRANS\,000 transitions), $\eta = 0.1$\,eV and $\hbar\gamma = 0.07$\,eV
(a typical value from optical measurements on gold; the colour is
insensitive to it). Figure~\ref{fig:optics} shows the result.

\begin{figure}[tbp]
\centering
\pgfplotsset{optplot/.style={width=0.5\textwidth, height=0.42\textwidth,
  xmin=0.5, xmax=6, xlabel={$\hbar\omega$ (eV)},
  tick label style={font=\small}, label style={font=\small},
  title style={font=\small}, legend style={font=\scriptsize, draw=none, fill=none}, legend cell align={left},
  no markers}}
\begin{tikzpicture}
\begin{axis}[optplot, ymin=0, ymax=13, title={Interband absorption $\varepsilon_2^{\mathrm{inter}}$},
  legend pos=north east]
\fill[yellow!25] (axis cs:1.6,0) rectangle (axis cs:3.3,13);
\addplot[red!70!black, thick] table[x index=0, y index=3] {figdata/optics_spec.dat};
\addplot[black!55, thin] table[x index=0, y index=9] {figdata/optics_spec.dat};
\addplot[black, only marks, mark=o, mark size=1.2pt] table[x index=0, y index=2] {figdata/optics_exp.dat};
\legend{$40^3$ mesh, $24^3$ mesh, measured $\varepsilon_2$}
\end{axis}
\end{tikzpicture}\hfill
\begin{tikzpicture}
\begin{axis}[optplot, ymin=0, ymax=1.05, title={Reflectivity $R$},
  legend style={at={(0.97,0.86)}, anchor=north east}]
\fill[yellow!25] (axis cs:1.6,0) rectangle (axis cs:3.3,1.05);
\addplot[red!70!black, thick] table[x index=0, y index=4] {figdata/optics_spec.dat};
\addplot[blue!60!black, thick, dashed] table[x index=0, y index=6] {figdata/optics_spec.dat};
\addplot[black!55, thin] table[x index=0, y index=7] {figdata/optics_spec.dat};
\addplot[black, only marks, mark=o, mark size=1.2pt] table[x index=0, y index=3] {figdata/optics_exp.dat};
\legend{computed (LDA), transitions $+0.75$\,eV, Drude term only, measured}
\end{axis}
\end{tikzpicture}
\caption{Optical response of gold from the reference calculation. The
shaded band is the visible range ($1.6$--$3.3$\,eV). Left: the interband
part of $\varepsilon_2$, Eq.~\eqref{eq:eps2}, for two $\kk$-meshes; the
remaining differences are the fine structure a finite mesh cannot
resolve. Right: reflectivity, Eq.~\eqref{eq:reflectivity}. Without
interband transitions gold would be a featureless mirror (grey line);
the $5d\to6s$ transitions cut the reflectivity down inside the visible
range (red). Dashed: the same calculation with all interband transition
energies raised by $0.75$\,eV, to show what a correctly placed absorption
edge does (see text). Circles: measured optical constants of Johnson and
Christy (1972); the measured $\varepsilon_2$ in the left panel is the
total, including the Drude part, which dominates below about $1.3$\,eV
(the measured point at $0.64$\,eV, $\varepsilon_2 = 25$, lies above the
frame).}
\label{fig:optics}
\end{figure}

\paragraph{What the calculation gets right.} Below about $1.3$\,eV
interband absorption is weak and gold is an excellent mirror
($R > 0.95$): this is the Drude regime, governed by $\omega_p$ alone.
Then the $5d\to6s$ transitions set in. $\varepsilon_2^{\mathrm{inter}}$
rises steeply, and the reflectivity collapses from above $0.9$ to below
$0.5$ within half an electron volt. Above the edge it stays between
$0.4$ and $0.5$ throughout the visible, close to the measured reflectivity
of gold for blue light (\REXPb{} at 450\,nm). The mechanism of gold's colour ---
a reflectivity edge in the visible, created by the $d$ bands lying only a
couple of eV below the Fermi level --- comes out of the band structure
with no adjustable parameter. So does the integrated strength: the
interband transitions below 8\,eV account for \NEFFINTER{} electrons per
atom in the sum rule, on top of the \NEFF{} conduction electrons.

\paragraph{What it gets wrong, and why.} The \emph{position} of the edge.
The computed reflectivity falls through $0.7$ at \EDGE\,eV; in real gold
(circles in Fig.~\ref{fig:optics}) the corresponding drop is at
$2.33$\,eV. The cause is not the
optics but the band structure it is fed: the $5d$ bands of this
calculation lie too close to the Fermi level. Their upper edge is at
$-1.05$\,eV at $X$ and $-1.63$\,eV at $L$ (Table~\ref{tab:goldlevels}),
whereas experiment places them $0.5$--$0.9$\,eV deeper
(Table~\ref{tab:dbandexp}). This is the ``known shortcoming of LDA
and GGA for the $d$ bands of noble metals'' met in
Chapter~\ref{ch:soc}, made somewhat worse here by the 11-electron
pseudopotential, and it is approximation (i) of the list above at work:
Kohn--Sham eigenvalues are not excitation energies.

The consequence for the colour is drastic, because the edge sits where the
eye is most sensitive. With the edge at \EDGE\,eV instead of $2.33$\,eV,
green and yellow light is absorbed as strongly as blue, and only red is
reflected well (Table~\ref{tab:colours}).

\begin{table}[tbp]
\centering
\begin{tabular}{@{}lcccc@{}}
\toprule
 & \multicolumn{3}{c}{reflectivity at} & colour\\
 & 450\,nm & 550\,nm & 650\,nm & sRGB\\
\midrule
Drude term only (no $d$ bands) & \RDRb & \RDRg & \RDRr &
\fcolorbox{black}{goldDR}{\rule{0pt}{1.6ex}\hspace{2.2em}}\ \ $(\RGBDR)$\\
computed, LDA & \RLDAb & \RLDAg & \RLDAr &
\fcolorbox{black}{goldLDA}{\rule{0pt}{1.6ex}\hspace{2.2em}}\ \ $(\RGBLDA)$\\
interband transitions $+0.5$\,eV & \RSHab & \RSHag & \RSHar &
\fcolorbox{black}{goldSHa}{\rule{0pt}{1.6ex}\hspace{2.2em}}\ \ $(\RGBSHa)$\\
interband transitions $+0.75$\,eV & \RSHbb & \RSHbg & \RSHbr &
\fcolorbox{black}{goldSHb}{\rule{0pt}{1.6ex}\hspace{2.2em}}\ \ $(\RGBSHb)$\\
measured (Johnson and Christy) & \REXPb & \REXPg & \REXPr &
\fcolorbox{black}{goldEXP}{\rule{0pt}{1.6ex}\hspace{2.2em}}\ \ $(\RGBEXP)$\\
\bottomrule
\end{tabular}
\caption{Reflectivity at three wavelengths and the resulting display colour
(swatch and 8-bit sRGB triple) for the reference calculation, two
illustrative shifts of the interband transition energies, the Drude term
alone, and the measured optical constants.}
\label{tab:colours}
\end{table}

The last row of Table~\ref{tab:colours} is itself a test of the colour chain of
Algorithm~\ref{alg:colour}: feeding \emph{measured} reflectivities
through it gives
\definecolor{agEXP}{RGB}{254,253,252}\definecolor{cuEXP}{RGB}{246,208,191}%
\fcolorbox{black}{goldEXP}{\rule{0pt}{1.6ex}\hspace{1.4em}} $(\RGBEXP)$
for gold,
\fcolorbox{black}{agEXP}{\rule{0pt}{1.6ex}\hspace{1.4em}} $(254,253,252)$
for silver and
\fcolorbox{black}{cuEXP}{\rule{0pt}{1.6ex}\hspace{1.4em}} $(246,208,191)$
for copper (optical constants of Johnson and Christy for all three) ---
the colours of the three metals. The parameter-free LDA result, by
contrast, is a pale, pinkish metal --- warmer than
silver, far from yellow: the calculation predicts \emph{that} gold is
coloured and \emph{why}, but not the right hue. The two ``shifted'' rows
are an illustration, not a prediction: every interband transition energy
$\omega_{mn}$ in the denominators of Eq.~\eqref{eq:epsinter} is raised by
a constant $\Delta$ while the strengths $|\mathbf v_{mn}|^2/\omega_{mn}^2$
are kept (in Algorithm~\ref{alg:colour}, replace
$S_t/(\omega_t[\omega_t^2-\dots])$ by
$(S_t/\omega_t^2)\,(\omega_t+\Delta)/[(\omega_t+\Delta)^2-\dots]$). As
the edge moves toward its experimental position, green and yellow light
is reflected again and the colour turns from pink through salmon to an
orange gold. A rigid shift is a caricature of the true correction (it
moves the edge by less than $\Delta$, because $\varepsilon_1$ changes
with it, and it leaves the shape of the spectrum unchanged), but it makes
the essential point: \emph{the colour of gold is a sensitive measurement
of the energy of its $d$ bands}, more sensitive than the accuracy of
LDA.

\begin{table}[H]
\centering\small
\begin{tabular}{@{}lccc@{}}
\toprule
Level (Kramers pair) & this calculation & PBE, semicore & experiment\\
\midrule
$\Gamma$: $d$ levels & $-5.45$, $-4.24$, $-3.02$ & $-5.67$, $-4.46$, $-3.27$ & $-6.0$, $-4.6$, $-3.7$\\
$L$: upper $d$ levels (pairs 4, 5) & $-2.36$, $-1.63$ & $-2.61$, $-1.90$ & $-3.2$, $-2.3$\\
$L$: $p$-like level (pair 6) & $-1.22$ & $-1.37$ & $-1.0$\\
$L$: $s$-like level (pair 7) & $+3.13$ & $+2.93$ & $+3.6$\\
\bottomrule
\end{tabular}
\caption{Band energies of gold in eV relative to $E_F$: the reference
calculation of this document (LDA, 11-electron pseudopotential), a
published all-electron-quality PBE calculation with semicore states, and
angle-resolved photoemission and optical experiments (the last two columns
from Rangel \emph{et al.}, Phys.\ Rev.\ B \textbf{86}, 125125 (2012),
Table~II; experimental values from different measurements scatter by
about $0.1$\,eV). The $d$ levels of both density-functional calculations
are too shallow; those of this calculation are a further
$0.2$--$0.3$\,eV shallower than the published PBE ones.}
\label{tab:dbandexp}
\end{table}

The \emph{strength} of the absorption is also overestimated: at 3\,eV the
computed $\varepsilon_2$ is \EPSCALC{} against a measured \EPSEXP. This
too is typical of the independent-particle approximation with LDA bands
for the noble metals.

\paragraph{How one would do better.} In order of increasing effort:
(a) the 19-electron pseudopotential of Appendix~\ref{app:q19} and the PBE
functional, which change the $d$-band position by tenths of an eV;
(b) a functional or method that describes localized $d$ states better
than LDA and GGA --- hybrid functionals or DFT$+U$; (c) many-body
perturbation theory ($GW$), in which quasiparticle energies replace
Kohn--Sham eigenvalues and which, for gold, widens the $5d$--$6sp$
separation by $0.35$--$0.75$\,eV (Rangel \emph{et al.}, cited above). All three change the
\emph{input} to this chapter; Algorithm~\ref{alg:colour} stays as it is.
None of them has been carried out for this document.

\paragraph{The role of relativity.} The argument of
Chapter~\ref{ch:soc} can now be made quantitative in the opposite
direction. Everything that distinguishes the red curve of
Fig.~\ref{fig:optics} from the grey one is the $5d$ band lying within
2\,eV of the Fermi level, and that proximity is the relativistic
contraction of the $6s$ shell and expansion of the $5d$ shell, built into
the pseudopotential. Silver, with the same crystal structure and nearly
the same lattice constant but weaker relativistic effects, has its $4d$
bands about 2\,eV deeper: its edge lies in the ultraviolet and its
visible reflectivity is the grey curve. Spin--orbit coupling itself plays
a secondary role for the colour: it moves the upper $d$-band edge by a few
tenths of an eV (compare the two halves of Table~\ref{tab:goldlevels}),
and the reflectivity edge with it (\EDGENOSOC\,eV without SOC against
\EDGESOC\,eV with SOC on the same $16^3$ mesh).

\paragraph{Numerical convergence.} The sRGB values change by up to about
ten units out of 255 when $\eta$ is varied between $0.05$ and $0.2$\,eV
($(\RGBETA)$ for $\eta = 0.2$\,eV) and by up to about fifteen units
between the $24^3$ and $40^3$ meshes ($(\RGBCO)$ on the coarser one); the Drude
damping $\gamma$ has no visible effect. These numerical uncertainties are
real but small compared with the effect of the misplaced $d$ bands.

\section{Sanity checks}

\begin{enumerate}[leftmargin=2em]
\item \textbf{Velocity matrix elements}: the Hellmann--Feynman test of
Section~\ref{sec:velocity}.
\item \textbf{Drude limit}: with the interband sum switched off,
$R(\omega)$ must be within a few percent of 1 for
$\gamma\ll\omega<\omega_p$ and drop steeply once $\omega$ exceeds
$\omega_p$; the colour must come out as a neutral near-white
(the reference calculation gives sRGB $(\RGBDR)$).
\item \textbf{Sum rule}: $N_{\mathrm{eff}}(\omega_{\max}) =
\frac{\Omega}{2\pi^2}\int_0^{\omega_{\max}}\omega\,\varepsilon_2(\omega)\,
\dd\omega$ counts the electrons per cell that have absorbed below
$\omega_{\max}$. The intraband term alone must give
$\omega_p^2\Omega/4\pi$ (about one electron), and the total must stay
below the number of valence electrons, 11.
\item \textbf{Colour chain}: white and 50\% grey as above; a reflectivity
that is 1 above 600\,nm and 0 below must come out red, and so on.
\item \textbf{Kramers--Kronig}: $\varepsilon_1$ from
Eq.~\eqref{eq:epsinter} must agree with a numerical Kramers--Kronig
transform of $\varepsilon_2$.
\item \textbf{Convergence}: repeat with a denser $\kk$-mesh and with
$\eta$ halved. The positions of the absorption edge and of the
reflectivity drop must not move; fine structure will.
\end{enumerate}

\begin{gpubox}{GPU/JAX notes}
The expensive part --- bands and velocity matrices on the dense mesh ---
was done in Chapter~\ref{ch:conductivity}. Everything in this chapter is a
sum over a list of a few hundred thousand transitions for each of a few
hundred frequencies: an outer product that fits one \texttt{vmap} over $\omega$,
or a few seconds of NumPy on the host. A faster and equally accurate
route is to histogram the strengths $S_t/\omega_t^2$ on a fine energy
grid once and convolve.
\end{gpubox}

\chapter{Verification: From Free Electrons to Relativistic Gold}
\label{ch:verify}

Build and test in stages; each stage has an exact or well-established
target. Do not proceed to a stage until the previous one passes.

\paragraph{About the reference numbers in this chapter.} Every formula,
table and algorithm in this document has been exercised by an independent
reference implementation that follows the text literally (plain NumPy,
double precision, no symmetry beyond time reversal). Numbers quoted
``from the reference calculation'' come from it. They are not
experimental or literature values; they are what \emph{these equations}
give with \emph{these parameters}, and a correct implementation must
reproduce them to the stated digits. Where a number depends on
discretization choices (FFT grid size, mesh), those are stated.

\section{Stage 0: components in isolation}
\begin{enumerate}[leftmargin=2em]
\item \textbf{Basis}: plane-wave counts per $\kk$ close to
$\Omega(2E_{\mathrm{cut}})^{3/2}/6\pi^2$ (gold cell:
$a = 7.70668$\,Bohr, $\Omega = 114.4306$\,Bohr$^3$, so about $1380$ at
$E_{\mathrm{cut}} = 40$\,Ha).
\item \textbf{FFT round trip}: sphere $\to$ box $\to$ IFFT $\to$ FFT $\to$
gather returns the input to machine precision; Parseval's theorem holds.
\item \textbf{XC values}: the anchors of Chapter~\ref{ch:xc}
($\varepsilon_x, \varepsilon_c$ at $r_s = 2, 3$ to $10^{-5}$; the PBE
point to $10^{-6}$); AD-generated $\vxc$ matches the analytic derivative.
\item \textbf{Ewald}: independence of $\eta$ to $10^{-10}$. Quantitative
target for the gold cell ($Z_{\mathrm{ion}} = 11$,
$a = 4.0782$\,\AA, one atom): $E_{\mathrm{Ewald}} = -35.992686$\,Ha,
equivalent to the FCC Madelung constant,
$E_{\mathrm{Ewald}} = -0.8958736\, Z_{\mathrm{ion}}^2 / r_{\mathrm{ws}}$
with $r_{\mathrm{ws}} = (3\Omega/4\pi)^{1/3}$.
\item \textbf{Hermiticity}: $\braket{\phi|\hat H \psi} =
\overline{\braket{\psi | \hat H \phi}}$ for random spinors, with SOC on ---
catches conjugation bugs in projectors and $\mathbf L\cdot\mathbf S$.
\item \textbf{$\mathbf L\cdot\mathbf S$ blocks}: the three checks of
Section~\ref{sec:lsmatrix}, in the real-harmonic basis you actually use;
then the $10\times10$ crystal-field model of Chapter~\ref{ch:soc} must give
$\{-1, +1, +2.5\}\,\xi$ with degeneracies $4, 2, 4$ for $\Delta = 2\xi$.
\item \textbf{Symmetry}: $[\hat P,\hat H]=0$, $[\Theta,\hat H]=0$,
$\Theta^2=-1$ on random spinors at the eight TRIM
(Section~\ref{sec:symmetry}).
\item \textbf{Pseudopotential pieces}: the closed forms
\eqref{eq:hghlocG}--\eqref{eq:hghlocpoly} and \eqref{eq:hghprojG} against
brute-force radial quadrature of Eqs.~\eqref{eq:hghloc} and
\eqref{eq:besseltransform}; projector normalization
$\int p_{i\ell}^2 r^2 \dd r = 1$; and $\alpha$ of Eq.~\eqref{eq:alphaZ}:
$\alpha = 61.595082$ for the LDA parameters of Table~\ref{tab:auhgh}.
\item \textbf{Smearing}: $\sum_{n\kk} w_\kk f_{n\kk} = N_e$ after the
Fermi search; $\dd s/\dd x = x\,\dd f/\dd x$ numerically for each scheme
(this is the identity that makes $F$ variational).
\item \textbf{Eigensolver}: for a small cutoff, build the dense matrix by
applying \textsc{ApplyH} to unit vectors and compare \texttt{eigh} with
Algorithm~\ref{alg:lobpcg} --- after 30 iterations \emph{and} after 300.
\end{enumerate}

\section{Stage 1: exactly solvable systems}
\begin{enumerate}[leftmargin=2em]
\item \textbf{Empty lattice} ($V = 0$): eigenvalues must equal
$\tfrac12|\kk+\GG|^2$ exactly; band structure is the folded free-electron
parabola. Tests kinetic term, eigensolver, $\kk$ bookkeeping.
\item \textbf{Single cosine potential}: $V(\rr) = 2V_0\cos(\GG_1\cdot\rr)$
with $\GG_1 = \mathbf b_1$. At $\kk = \tfrac12\mathbf b_1$ (an $L$ point)
the two lowest levels must be $\tfrac18|\mathbf b_1|^2 \pm V_0$ up to
corrections of order $V_0^2/|\mathbf b_1|^2$, i.e.\ a gap
$\approx 2|V_0|$, in agreement with Section~\ref{sec:nfe} and with a dense
diagonalization; and their parities (Eq.~\eqref{eq:parityperm}) must be
even-below-odd for $V_0<0$ and the reverse for $V_0>0$.
\item \textbf{A screened ion in a big box.} This is the decisive test of
the pseudopotential code, SOC included, against an oracle that shares no
code with it. Put one gold ion in a large FCC cell ($a = 24$\,Bohr),
switch XC off, and use as the only other potential the Hartree potential
of a \emph{fixed} Gaussian charge of width $w = 2$\,Bohr,
$n(\GG) = (Z_{\mathrm{ion}}/\Omega)\,\ee^{-G^2w^2/2}$ --- the SCF starting
density with a different width. The total potential is then short-ranged
and spherically symmetric about the atom,
\[
V(r) = V_{\mathrm{loc}}(r) + \frac{Z_{\mathrm{ion}}}{r}\,
\mathrm{erf}\!\Big(\frac{r}{\sqrt2\,w}\Big) \;+\; \hat V_{\mathrm{NL}},
\]
so its bound states are those of a one-dimensional radial problem for each
$(\ell, j)$. Solve that problem independently: on a uniform radial grid
$r_i = ih$, the function $u(r) = r\,R(r)$ obeys a matrix eigenproblem with
\begin{multline*}
H_{ii'} = -\frac{\delta_{i,i'+1} - 2\delta_{ii'} + \delta_{i,i'-1}}{2h^2}
+ \Big[\frac{\ell(\ell+1)}{2r_i^2} + V(r_i)\Big]\delta_{ii'}\\
+ h\sum_{a,b=1}^{2} r_i\,p_{a\ell}(r_i)\;\big(h^\ell_{ab} + \lambda_{\ell j}\,k^\ell_{ab}\big)\;
r_{i'}\,p_{b\ell}(r_{i'}),
\end{multline*}
where $\lambda_{\ell j} = \ell/2$ for $j = \ell+\tfrac12$ and
$-(\ell+1)/2$ for $j = \ell-\tfrac12$ (Eq.~\eqref{eq:ls_eigen}) --- thirty
lines of code, one dense \texttt{eigh} per channel ($h = 0.01$,
$r_{\max} = 30$). Reference values (extrapolated to $h\to0$):
\begin{center}
\begin{tabular}{@{}lccccc@{}}
\toprule
 & $s_{1/2}$ & $p_{1/2}$ & $p_{3/2}$ & $d_{3/2}$ & $d_{5/2}$\\
\midrule
$\varepsilon$ (Ha) & $-1.01843$ & $-0.60281$ & $-0.59335$ & $-1.99901$ & $-1.89459$\\
degeneracy in the box & 2 & 2 & 4 & 4 & 6\\
\bottomrule
\end{tabular}
\end{center}
The plane-wave eigenvalues at $\Gamma$ must equal these \emph{plus}
$2\pi Z_{\mathrm{ion}}w^2/\Omega$, with the degeneracies shown. (The shift
is the price of setting $\vH(\GG{=}0)=0$: the plane-wave potential differs
from the true one by exactly this constant --- a useful exercise in
Section~\ref{sec:g0}.) The reference plane-wave calculation ($E_{\mathrm{cut}} = 40$\,Ha,
$41\,655$ plane waves, FFT grid $98^3$, shift $0.079994$\,Ha) gives
$-1.01843$, $-0.60281$, $-0.59335$, $-1.99900$, $-1.89458$\,Ha: agreement
to $10^{-5}$\,Ha in every channel. The $d_{5/2}$--$d_{3/2}$
splitting, $0.1044$\,Ha with $j = \tfrac52$ \emph{above} $j=\tfrac32$,
fixes the overall sign of the $k^\ell_{ij}$ term.
\item \textbf{The self-consistent atom (oracle only).} Adding a Hartree and
LDA potential to the radial solver and iterating gives the neutral
pseudo-atom. For the configuration $5d_{3/2}^{4}\,5d_{5/2}^{6}\,6s^1$
(spherical, not spin-polarized, LDA parameters):
$\varepsilon_{6s} = -0.2241$, $\varepsilon_{5d_{3/2}} = -0.2971$,
$\varepsilon_{5d_{5/2}} = -0.2405$\,Ha, hence a spin--orbit splitting of
the $5d$ level of $1.54$\,eV --- the number quoted in
Chapter~\ref{ch:soc}, here derived from the twenty numbers of
Table~\ref{tab:auhgh}. (The unoccupied $6p$ levels come out at $-0.0486$
and $-0.0278$\,Ha, split by $0.57$\,eV with $j=\tfrac12$ lower.)
\end{enumerate}

\section{Stage 2: self-consistency on easy targets}
\begin{enumerate}[leftmargin=2em]
\item \textbf{Jellium in a cell}: no ions, uniform positive background
(i.e.\ just $\vH(\GG{=}0) = 0$ and no ionic potential), $N_e$ electrons.
The SCF must return the uniform density $N_e/\Omega$, eigenvalues
$\tfrac12|\kk+\GG|^2 + \vxc(n)$, and the analytic HEG energy per electron
$\tfrac{3}{10}k_F^2 + \varepsilon_{\mathrm{xc}}(n)$ in the limit of a dense
mesh and small smearing.
\item \textbf{Silicon (bonus)}: if you want a soft interim target, HGH
silicon converges with $E_{\mathrm{cut}} \approx 12$\,Ha; indirect
gap $\approx 0.5$--$0.6$\,eV (the famous LDA/PBE underestimate), lattice
constant PBE $\approx 5.47$\,\AA. (Two atoms per cell: your structure
factors and projector phases get their first real test.)
\end{enumerate}

\section{Stage 3: gold, self-consistent --- regression numbers}
\label{sec:stage3}

Start small, and start with spin--orbit coupling switched off
($k^\ell = 0$): the three cases of Table~\ref{tab:smallcases} run in a
minute on a laptop CPU and each isolates one ingredient. Work through them
in order --- scalar-relativistic LDA first, then SOC, then PBE.

\begin{table}[H]
\centering\small
\begin{tabular}{@{}lccl@{}}
\toprule
Case & $F$ (Ha) & $\mu$ (Ha) & levels at $\Gamma$ in Ha (number of spinor states)\\
\midrule
LDA, SOC off & $-33.291161$ & $0.613238$ & $0.213490\,(2)$, $0.406739\,(6)$, $0.463560\,(4)$\\
LDA, SOC on & $-33.311413$ & $0.609888$ & $0.213402\,(2)$, $0.381623\,(4)$, $0.425863\,(2)$, $0.470840\,(4)$\\
PBE, SOC on & $-33.254853$ & $0.624314$ & $0.229208\,(2)$, $0.394565\,(4)$, $0.438491\,(2)$, $0.482719\,(4)$\\
\bottomrule
\end{tabular}
\caption{Small regression cases. Common settings: $q = 11$ parameters of
Table~\ref{tab:auhgh} matching the functional, $a = 4.0782$\,\AA{}
$= 7.706681$\,Bohr, $E_{\mathrm{cut}} = 20$\,Ha, FFT grid $24^3$ ($531$
plane waves at $\Gamma$), $\Gamma$-centered $4^3$ mesh ($36$ points after
time-reversal reduction; density accumulated as $n$ only, with
$\mathbf m = 0$, as Section~\ref{sec:symmetry} requires on such a mesh), Methfessel--Paxton order 1 with
$\sigma = 0.01$\,Ha, 20 spinor bands, SCF as in the driver of
Chapter~\ref{ch:algorithms}, converged to $\|R\| < 10^{-8}$ (10--12
iterations). The levels at $\Gamma$ are the eigenvalues of the converged
Hamiltonian; with SOC off the spinor code returns every level of a scalar
code twice (the $t_{2g}$ triplet as six states, the $e_g$ doublet as
four). The cutoff is far too low for converged physics --- these are
regression numbers.}
\label{tab:smallcases}
\end{table}

A correct implementation reproduces $F$, $\mu$ and the levels to the last
digit shown: both the NumPy reference and an independent JAX program
written from this text do. If you are off, the size of the discrepancy
tells you where to look:
\begin{itemize}[leftmargin=2em]
\item \emph{$10^{-6}$--$10^{-5}$\,Ha in $F$}: rounded constants (PBE
$\beta,\mu$; pseudopotential digits; PW92 $A$) or a different FFT grid.
\item \emph{$10^{-3}$\,Ha or more in $F$ with correct eigenvalue
differences}: a constant --- $\alpha$ of Eq.~\eqref{eq:alphaZ}, the Ewald
energy, the entropy term, or double counting in Eq.~\eqref{eq:ftot}.
\item \emph{Wrong degeneracies at $\Gamma$} (anything other than $2,6,4$
without and $2,4,2,4$ with SOC): projectors, real harmonics or the
$\mathbf L\cdot\mathbf S$ blocks; go back to Stage~0 and Stage~1.3.
\item \emph{SOC-off case right, SOC-on wrong}: the $k^\ell_{ij}$ values or
the assembly of $D$.
\item \emph{LDA right, PBE wrong}: the gradient or divergence step of
Eq.~\eqref{eq:vxcgga2}; test the PBE point of Chapter~\ref{ch:xc} first.
\end{itemize}

Two larger cases with spin--orbit coupling on, for the production
settings:
\begin{center}\small
\begin{tabular}{@{}lll@{}}
\toprule
Setting & Case A & Case B\\
\midrule
functional, parameters & LDA, $q=11$ & LDA, $q=11$\\
$a$ & $7.706681$\,Bohr & $7.706681$\,Bohr\\
$E_{\mathrm{cut}}$; FFT grid & $25$\,Ha; $27^3$ & $40$\,Ha; $35^3$\\
mesh (time-reversal reduced) & $4^3$, Eq.~\eqref{eq:mpmesh}: 32 points & $8^3$, $\Gamma$-centered: 260 points\\
smearing & MP1, $\sigma=0.01$\,Ha & MP1, $\sigma=0.01$\,Ha\\
spinor bands & 20 & 20\\
\midrule
$F$ (Ha) & $-33.324309$ & $-33.325815$\\
$\mu$ (Ha) & $0.58409$ & $0.58322$\\
SCF iterations to $\|R\|<10^{-7}$ & 10 & 10\\
\bottomrule
\end{tabular}
\end{center}
(Throughout: $a_0 = 0.529177211$\,\AA; PW92 constants of
Table~\ref{tab:pw92}. Case~A deliberately uses the shifted mesh of
Eq.~\eqref{eq:mpmesh} to exercise that code path; it is a regression
number, not a recommendation --- see Section~\ref{sec:bzintegration}.)
Case~B is the calculation behind
Table~\ref{tab:goldlevels}, Fig.~\ref{fig:bands} and
Chapter~\ref{ch:topology}.

\paragraph{Then the physics.} Lattice constant and bulk modulus from
equation-of-state fits (end of Chapter~\ref{ch:hamiltonian}; seven
lattice constants spanning $\pm4\%$, scalar-relativistic,
$E_{\mathrm{cut}} = 45$\,Ha, $\Gamma$-centered $8^3$ mesh, MP1
$\sigma = 0.01$\,Ha), and the position of the $d$ bands:
\begin{center}\small\setlength{\tabcolsep}{4pt}
\begin{tabular}{@{}lccccc@{}}
\toprule
 & \multicolumn{2}{c}{LDA} & \multicolumn{2}{c}{PBE} & Experiment\\
 & this potential & literature & this potential & literature & (room temp.)\\
\midrule
Lattice constant $a_0$ (\AA) & $4.093$ & $4.053$ & $4.200$ & $4.173$ & $4.078$\\
Bulk modulus $B_0$ (GPa) & $184$ & $195$ & $133$ & $143$ & $173$\\
$B_0'$ & $5.7$ & & $6.1$ & & \\
$d$-band top at $L$ (eV below $E_F$) & $1.63$ & & & $1.90$ & $2.3$\\
\bottomrule
\end{tabular}
\end{center}
``This potential'' is the reference calculation with the 11-electron
parameters of Table~\ref{tab:auhgh}; ``literature'' are
all-electron-quality values (Ambrosetti and Silvestrelli 2016 for $a_0$
and $B_0$; Rangel \emph{et al.}\ 2012 for the $d$ band, which is quoted
with spin--orbit coupling). Three lessons.
(i) The familiar statement ``LDA overbinds, PBE underbinds'' is visible in
both pairs of columns: PBE gives a lattice constant $0.11$--$0.12$\,\AA{}
larger and a bulk modulus about $50$\,GPa smaller than LDA. (ii) The
11-electron pseudopotential adds $0.04$\,\AA{} (LDA) and $0.03$\,\AA{} (PBE)
to $a_0$ and takes
about $10$\,GPa off $B_0$ compared with all-electron-quality
calculations --- enough to move the LDA lattice constant from slightly
below to slightly above experiment. This is a property of the
pseudopotential (no semicore states, no nonlinear core correction), not
of your code: reproduce the first and third columns. (iii) Spin--orbit
coupling changes these numbers little (Stage~4): at
$E_{\mathrm{cut}} = 35$\,Ha with a $6^3$ mesh, LDA gives
$4.095$\,\AA, $184$\,GPa without and $4.084$\,\AA, $190$\,GPa with SOC;
comparison with the first column also shows that the equation of state
is converged with respect to cutoff and mesh at the level of
$0.003$\,\AA{} and $1$\,GPa. The $d$-band position discrepancy versus
photoemission and optics is a known limitation of LDA and GGA ---
reproduce the DFT number, not experiment. The SCF for the metal must converge in
about ten iterations with Kerker--Pulay mixing; watch the difference
between the two energy expressions \eqref{eq:fdirect} and \eqref{eq:ftot}
close quadratically. A typical history (reference calculation,
$E_{\mathrm{cut}} = 40$\,Ha, $8^3$ mesh, SOC on):
\begin{center}\small
\begin{tabular}{@{}cccc@{}}
\toprule
iteration & $\|R\|$ & $F$ (Ha), Eq.~\eqref{eq:ftot} & Eq.~\eqref{eq:fdirect} $-$ Eq.~\eqref{eq:ftot}\\
\midrule
1 & $6.1\times10^{-1}$ & $-33.541764$ & $4.2\times10^{-1}$\\
2 & $2.3\times10^{-1}$ & $-33.330984$ & $1.3\times10^{-2}$\\
3 & $5.4\times10^{-2}$ & $-33.325191$ & $-5.9\times10^{-4}$\\
4 & $1.8\times10^{-3}$ & $-33.325816$ & $2.1\times10^{-6}$\\
5 & $5.9\times10^{-4}$ & $-33.325816$ & $6.0\times10^{-7}$\\
6 & $3.4\times10^{-5}$ & $-33.325815$ & $-6.2\times10^{-10}$\\
8 & $2.0\times10^{-6}$ & $-33.325815$ & $7\times10^{-12}$\\
10 & $4.7\times10^{-8}$ & $-33.325815$ & $-4\times10^{-13}$\\
\bottomrule
\end{tabular}
\end{center}

\section{Stage 4: gold with spin--orbit coupling --- the physics}
\begin{enumerate}[leftmargin=2em]
\item \textbf{Kramers pairs}: every eigenvalue doubly degenerate at
every $\kk$ (inversion $\times$ time reversal), to the accuracy of the
eigensolver (the reference calculation finds $4\times10^{-11}$\,Ha on the
SCF mesh after three LOBPCG iterations per step).
\item \textbf{Cubic degeneracies}: with a $\Gamma$-centered mesh, the
quartets at $\Gamma$ are exact and all symmetry-equivalent $\kk$-points
agree to solver accuracy.
\item \textbf{Level scheme at the high-symmetry points},
Table~\ref{tab:goldlevels}: without SOC two $d$ levels at $\Gamma$ (a
$t_{2g}$ triplet and an $e_g$ doublet, each $\times2$ for spin); with SOC
three, with degeneracies $4, 2, 4$ in ascending order. Published
fully-relativistic band structures of Au agree with these values to
within the differences between pseudopotentials and functionals
($\sim0.1$\,eV).
\item \textbf{Energetics barely move}: switching SOC on lowers the
free energy by about $0.02$\,Ha, reduces $a_0$ by $0.011$\,\AA{} and
raises $B_0$ by $6$\,GPa (reference equation-of-state fits, LDA,
$E_{\mathrm{cut}} = 35$\,Ha, $6^3$ mesh: $4.095\to4.084$\,\AA,
$184\to190$\,GPa) --- a check that SOC entered as a (strong) band
effect, not as an energy-scale bug.
\item \textbf{Switch test}: SOC on/off touches only the $k^{\ell}$ blocks;
with $k^{\ell} = 0$ the spinor code must return every eigenvalue exactly
twice (and, if you also wrote a scalar code with two electrons per
orbital, reproduce its eigenvalues and free energy to solver precision).
\end{enumerate}

\begin{table}[H]
\centering\small
\begin{tabular}{@{}lll@{}}
\toprule
Point & SOC & Levels (eV relative to $E_F$)\\
\midrule
$\Gamma$ & off & $-10.06$,\ \ $-4.82\,(\times3)$,\ \ $-3.28\,(\times2)$,\ \ $15.97$\\
         & on  & $-10.06$,\ \ $-5.45\,(\times2)$,\ \ $-4.24$,\ \ $-3.02\,(\times2)$,\ \ $15.97$\\
\midrule
$X$ & off & $-7.20$,\ \ $-6.88$,\ \ $-2.09$,\ \ $-1.70\,(\times2)$,\ \ $0.96$,\ \ $5.76$\\
    & on  & $-7.32$,\ \ $-7.03$,\ \ $-2.31$,\ \ $-2.09$,\ \ $-1.05$,\ \ $0.71$,\ \ $5.78$\\
\midrule
$L$ & off & $-7.47$,\ \ $-4.86\,(\times2)$,\ \ $-2.03\,(\times2)$,\ \ $-1.06$,\ \ $3.07$\\
    & on  & $-7.56$,\ \ $-5.57$,\ \ $-4.38$,\ \ $-2.36$,\ \ $-1.63$,\ \ $-1.22$,\ \ $3.13$\\
\midrule
$W$ & off & $-6.10$,\ \ $-5.28\,(\times2)$,\ \ $-3.42$,\ \ $-1.70$,\ \ $4.71$,\ \ $6.40\,(\times2)$\\
    & on  & $-6.31$,\ \ $-5.73$,\ \ $-5.15$,\ \ $-3.25$,\ \ $-1.63$,\ \ $4.70$,\ \ $5.60$,\ \ $6.91$\\
\midrule
$K$ & off & $-6.59$,\ \ $-5.97$,\ \ $-3.95$,\ \ $-2.98$,\ \ $-2.16$,\ \ $3.85$,\ \ $4.73$\\
    & on  & $-6.75$,\ \ $-6.17$,\ \ $-4.06$,\ \ $-3.01$,\ \ $-1.95$,\ \ $3.76$,\ \ $4.63$\\
\bottomrule
\end{tabular}
\caption{Kohn--Sham levels of gold in eV relative to $E_F$ from the
reference calculation (LDA, $q=11$, $a = 4.0782$\,\AA,
$E_{\mathrm{cut}} = 40$\,Ha, $\Gamma$-centered $8^3$ mesh, MP1
$\sigma = 0.01$\,Ha; $\mu = 0.58322$\,Ha, $F = -33.325815$\,Ha). ``SOC
off'' levels were evaluated in the same self-consistent potential with
$k^\ell = 0$. Multiplicities count spatial states (SOC off) or Kramers
pairs (SOC on). Compare with Fig.~\ref{fig:bands}.}
\label{tab:goldlevels}
\end{table}

\section{Stage 5: topology}
\begin{enumerate}[leftmargin=2em]
\item The sanity checks of Chapter~\ref{ch:topology} (empty-lattice
parities, the cosine-potential inversion, $\Pi^2 = \mathbf 1$,
$[\hat P,\hat H] = 0$, all-even at $\Gamma$, SOC-off parities,
Wilson-loop pinning, robustness).
\item Target: $\delta_\Gamma = +1$, $\delta_X = -1$, $\delta_L = -1$,
hence $(\nu_0;\nu_1\nu_2\nu_3) = (1;000)$ for the lowest six Kramers
pairs, with a strictly positive direct gap between spinor bands 12 and 13
(Section~\ref{sec:directgap}).
\item Wilson loops: the end-point pattern of Fig.~\ref{fig:wcc} (one
Kramers pair at $\theta=\pi$ only on the line through $\Gamma$), and,
with sufficient resolution, $z(k_3{=}0) = 1$, $z(k_3{=}\tfrac12) = 0$,
hence the same $\nu_0 = 1$, $\nu_3 = 0$.
\end{enumerate}

\section{Stage 6: transport and optics}
\begin{enumerate}[leftmargin=2em]
\item \textbf{Velocity operator}: the Hellmann--Feynman test of
Section~\ref{sec:velocity} at a generic $\kk$ (the reference uses
fractional $\kk = (0.13, 0.27, 0.41)$): finite-difference band slopes and
pair-averaged diagonal velocity elements agree to $10^{-7}$\,a.u.; the
velocity matrices are Hermitian; dropping the nonlocal term changes the
velocities by $0.1$--$0.3$\,a.u.
\item \textbf{Fermi-surface numbers} (fixed potential of Case~B,
irreducible wedge of a $\Gamma$-centered $40^3$ mesh, Gaussian width
$0.005$\,Ha): $g(E_F) = \GFb$ states/eV per atom,
$\hbar\omega_p = \WPb$\,eV, r.m.s.\ Fermi velocity
$\VFb\times10^{6}$\,m/s, effective number of carriers \NEFF. Expect
agreement to a few percent, not to all digits: these are Fermi-surface
integrals and depend on mesh and broadening.
\item \textbf{Colour chain}: $R\equiv1$ gives sRGB $(255,255,255)$;
$R\equiv0.5$ gives $(188,188,188)$; a mirror that reflects only for
$\lambda > 600$\,nm gives $(234,0,0)$; one that reflects only for
$\lambda < 480$\,nm gives $(67,0,255)$ (strict inequalities on the 1\,nm
grid; with $\ge$ and $\le$ the values become $(237,0,0)$ and
$(65,0,255)$).
\item \textbf{Spectrum}: interband $\varepsilon_2$ below 1 up to
$\approx1.2$\,eV and above 2 from $\ONSET$\,eV, reflectivity above $0.9$ in
the near infrared and falling through $0.7$ at $\EDGE$\,eV
(Fig.~\ref{fig:optics}); Kramers--Kronig consistency of
Eq.~\eqref{eq:epsinter}; sum rule of Chapter~\ref{ch:colour}.
\end{enumerate}

\section{Stage 7: convergence studies (the habits of a professional)}
For each observable, produce convergence curves in
(i) $E_{\mathrm{cut}}$ at fixed mesh, (ii) $\kk$-mesh at fixed
$E_{\mathrm{cut}}$, (iii) smearing $\sigma \to 0$ jointly with the mesh.
Quote every number in your write-up with the settings that converge it to a
stated tolerance. This discipline --- not any single number --- is what makes
a DFT calculation trustworthy.

\section*{Further reading}
\begin{itemize}[leftmargin=2em]
\item R.\ M.\ Martin, \emph{Electronic Structure: Basic Theory and Practical
Methods} (CUP) --- the encyclopedic companion to everything here.
\item M.\ C.\ Payne et al., Rev.\ Mod.\ Phys.\ \textbf{64}, 1045 (1992) ---
the classic plane-wave/iterative-methods review.
\item C.\ Hartwigsen, S.\ Goedecker, J.\ Hutter, Phys.\ Rev.\ B \textbf{58},
3641 (1998) --- the relativistic HGH pseudopotential tables (your gold
parameters, including SOC).
\item J.\ P.\ Perdew, K.\ Burke, M.\ Ernzerhof, Phys.\ Rev.\ Lett.\
\textbf{77}, 3865 (1996); J.\ P.\ Perdew, Y.\ Wang, Phys.\ Rev.\ B
\textbf{45}, 13244 (1992) --- PBE and PW92.
\item A.\ V.\ Knyazev, SIAM J.\ Sci.\ Comput.\ \textbf{23}, 517 (2001) ---
LOBPCG; U.\ Hetmaniuk, R.\ Lehoucq, J.\ Comput.\ Phys.\ \textbf{218}, 324
(2006), and J.\ A.\ Duersch et al., SIAM J.\ Sci.\ Comput.\ \textbf{40},
C655 (2018) --- why and how to keep its basis orthonormal.
\item G.\ Kresse, J.\ Furthm\"uller, Phys.\ Rev.\ B \textbf{54}, 11169
(1996) --- iterative diagonalization, Pulay mixing and Kerker
preconditioning for metals.
\item P.\ Pyykk\"o, Chem.\ Rev.\ \textbf{88}, 563 (1988) --- relativity in
chemistry; why gold is gold.
\item M.\ Krack, Theor.\ Chem.\ Acc.\ \textbf{114}, 145 (2005) --- the PBE
parameter sets of the same pseudopotential family.
\item D.\ Vanderbilt, \emph{Berry Phases in Electronic Structure Theory}
(CUP, 2018) --- the book-length version of Chapter~\ref{ch:topology}.
\item L.\ Fu, C.\ L.\ Kane, Phys.\ Rev.\ B \textbf{76}, 045302 (2007) ---
the parity formula for the $\mathbb{Z}_2$ invariant.
\item A.\ A.\ Soluyanov, D.\ Vanderbilt, Phys.\ Rev.\ B \textbf{83}, 235401
(2011) --- the Wilson-loop (Wannier charge center) method.
\item B.\ Yan et al., Nat.\ Commun.\ \textbf{6}, 10167 (2015) ---
topological states on the gold surface.
\item S.\ LaShell, B.\ A.\ McDougall, E.\ Jensen, Phys.\ Rev.\ Lett.\
\textbf{77}, 3419 (1996) --- the spin-split surface state of Au(111).
\item N.\ W.\ Ashcroft, N.\ D.\ Mermin, \emph{Solid State Physics},
chapters 1, 12 and 13 --- the Drude model and semiclassical transport
behind Chapter~\ref{ch:conductivity}.
\item P.\ B.\ Allen, Phys.\ Rev.\ B \textbf{17}, 3725 (1978) --- the
electron--phonon theory of the resistivity of metals
(Eq.~\eqref{eq:rhoT}); S.\ Y.\ Savrasov, D.\ Y.\ Savrasov, Phys.\ Rev.\ B
\textbf{54}, 16487 (1996) --- first-principles $\lambda_{\mathrm{tr}}$
and resistivities of eight metals; T.\ Gunst, A.\ Blom, K.\ Stokbro,
Proc.\ SISPAD 2018, p.~79 --- phonon-limited resistivity of Au, Ag, Cu.
\item T.\ Rangel et al., Phys.\ Rev.\ B \textbf{86}, 125125 (2012) ---
the band structure of gold from many-body perturbation theory, with a
table of measured band energies (Table~\ref{tab:dbandexp}).
\item A.\ Ambrosetti, P.\ L.\ Silvestrelli, Phys.\ Rev.\ B \textbf{94},
045124 (2016) --- lattice constants and bulk moduli of Cu, Ag, Au with
LDA, PBE and beyond.
\item P.\ B.\ Johnson, R.\ W.\ Christy, Phys.\ Rev.\ B \textbf{6}, 4370
(1972) --- the measured optical constants of the noble metals (the
tables used here are in \texttt{verification/reference\_data/}, from
the public-domain refractiveindex.info database); R.\ L.\ Olmon et al.,
Phys.\ Rev.\ B \textbf{86}, 235147 (2012) --- a modern remeasurement
for gold.
\item C.\ Wyman, P.-P.\ Sloan, P.\ Shirley, J.\ Comput.\ Graph.\ Tech.\
\textbf{2}(2), 1 (2013) --- the analytic colour-matching functions of
Eq.~\eqref{eq:cmf}.
\item The DFTK.jl and GPAW documentation --- two readable open-source
plane-wave codes to compare designs against (resist reading their source
until yours works).
\end{itemize}

\vspace{1em}
\begin{keybox}{You are ready}
Every equation your program needs is now on paper: the spinor Kohn--Sham
problem \eqref{eq:spinorks}, its plane-wave matrix elements
(Chapter~\ref{ch:hamiltonian}), the XC input (Chapter~\ref{ch:xc}), the
relativistic ionic data (Chapters~\ref{ch:pseudo}--\ref{ch:soc}), and the
solvers (Chapter~\ref{ch:algorithms}) in GPU-native form
(Chapter~\ref{ch:gpu}). Build the program in the order of this chapter:
components, empty lattice, one screened gold ion, scalar-relativistic
gold, the full spin--orbit band structure of the metal that is yellow
because of special relativity, its topological invariant, and finally
its Fermi-surface and optical properties.
\end{keybox}


\appendix
\chapter{Semicore (19-electron) Parameters for Gold}
\label{app:q19}

The 19-electron potentials treat $5s^25p^6$ as valence in addition to
$5d^{10}6s^1$. The formalism is unchanged; what changes is the size of the
tables: the local part has $r_{\mathrm{loc}} = 0.49$ (and, for LDA, a
second coefficient $C_2$), the $s$ and $p$ channels carry \emph{three}
projectors (use the $i=3$ lines of Eq.~\eqref{eq:hghprojG}), and the
coupling matrices are $3\times3$. There are then $3 + 3\cdot3 + 2\cdot5 =
22$ projector functions per atom. $N_e = 19$, so solve for at least
$\sim 30$ spinor bands; the lowest eight (four Kramers pairs: $5s$ and
$5p$) form narrow bands far below the $d$ bands, which shifts every band
index of Chapter~\ref{ch:topology} up by eight ($N = 10$ pairs below the
direct gap). Because the projector radii are much smaller, the cutoff must
be considerably higher than for $q = 11$: converge it explicitly, starting
from $E_{\mathrm{cut}}\approx 80$\,Ha.

Tables~\ref{tab:q19lda} and \ref{tab:q19pbe} come from the same database
as Table~\ref{tab:auhgh} (its entries \texttt{Au GTH-PADE-q19} and
\texttt{Au GTH-PBE-q19}). Each block lists the upper triangle of a
symmetric matrix.

\begin{table}[h]
\centering\small
\begin{tabular}{@{}llrrr@{}}
\toprule
\multicolumn{5}{@{}l}{\textbf{LDA}, $Z_{\mathrm{ion}} = 19$:\quad
$r_{\mathrm{loc}} = 0.49000000$,\ \ $C_1 = 9.07322639$,\ \ $C_2 = -0.34013367$}\\
\midrule
channel & & \multicolumn{1}{c}{$j=1$} & \multicolumn{1}{c}{$j=2$} & \multicolumn{1}{c}{$j=3$}\\
\midrule
$s$:\ $r_s = 0.28803780$ & $h_{1j}$ & $-5.90182040$ & $26.42537412$ & $-15.50837777$\\
 & $h_{2j}$ & & $-59.33938963$ & $40.04245923$\\
 & $h_{3j}$ & & & $-31.78271670$\\
\midrule
$p$:\ $r_p = 0.35586747$ & $h_{1j}$ & $-7.08858659$ & $8.19796427$ & $0$\\
 & $h_{2j}$ & & $-9.69996213$ & $0$\\
 & $h_{3j}$ & & & $0$\\
 & $k_{1j}$ & $0.62807493$ & $1.81668773$ & $-1.65141098$\\
 & $k_{2j}$ & & $-4.89888134$ & $3.90795165$\\
 & $k_{3j}$ & & & $-2.77740200$\\
\midrule
$d$:\ $r_d = 0.33386175$ & $h_{1j}$ & $-8.82572349$ & $9.56535285$ & \\
 & $h_{2j}$ & & $-10.84609065$ & \\
 & $k_{1j}$ & $0.18880114$ & $0.03175114$ & \\
 & $k_{2j}$ & & $-0.03600240$ & \\
\bottomrule
\end{tabular}
\caption{HGH parameters for Au with 19 valence electrons, LDA (Hartree
atomic units).}
\label{tab:q19lda}
\end{table}

\begin{table}[h]
\centering\small
\begin{tabular}{@{}llrrr@{}}
\toprule
\multicolumn{5}{@{}l}{\textbf{PBE}, $Z_{\mathrm{ion}} = 19$:\quad
$r_{\mathrm{loc}} = 0.49000000$,\ \ $C_1 = 8.36945885$}\\
\midrule
channel & & \multicolumn{1}{c}{$j=1$} & \multicolumn{1}{c}{$j=2$} & \multicolumn{1}{c}{$j=3$}\\
\midrule
$s$:\ $r_s = 0.29008402$ & $h_{1j}$ & $-5.98869002$ & $26.10502137$ & $-15.17621948$\\
 & $h_{2j}$ & & $-58.38556050$ & $39.18483021$\\
 & $h_{3j}$ & & & $-31.10199477$\\
\midrule
$p$:\ $r_p = 0.36150576$ & $h_{1j}$ & $-5.87695332$ & $4.73144606$ & $2.92710526$\\
 & $h_{2j}$ & & $-0.72513863$ & $-6.92679530$\\
 & $h_{3j}$ & & & $4.92291022$\\
 & $k_{1j}$ & $0.18370022$ & $2.64213942$ & $-2.01606276$\\
 & $k_{2j}$ & & $-6.48265855$ & $4.77087526$\\
 & $k_{3j}$ & & & $-3.39068639$\\
\midrule
$d$:\ $r_d = 0.34037313$ & $h_{1j}$ & $-8.80458432$ & $9.49191532$ & \\
 & $h_{2j}$ & & $-10.76282032$ & \\
 & $k_{1j}$ & $0.15609346$ & $0.06062427$ & \\
 & $k_{2j}$ & & $-0.06874145$ & \\
\bottomrule
\end{tabular}
\caption{HGH parameters for Au with 19 valence electrons, PBE (Hartree
atomic units).}
\label{tab:q19pbe}
\end{table}

\paragraph{How these tables were verified.} (i) A script compares every
number of Tables~\ref{tab:auhgh}, \ref{tab:q19lda} and \ref{tab:q19pbe}
with the source file, which is itself byte-identical to the current CP2K
distribution: all digits agree. (ii) The parameters were exercised in the
self-consistent radial atom solver of Chapter~\ref{ch:verify} (Stage~1),
extended to three projectors and, for PBE, to gradient corrections. If
the 11- and 19-electron potentials describe the same atom, the levels
they share must coincide --- and they do, to $0.15$\,mHa:
\begin{center}\small
\begin{tabular}{@{}llccccc@{}}
\toprule
 & & $6s$ & $5d_{3/2}$ & $5d_{5/2}$ & $6p_{1/2}$ & $6p_{3/2}$\\
\midrule
LDA & $q=11$ & $-0.22409$ & $-0.29705$ & $-0.24051$ & $-0.04859$ & $-0.02758$\\
    & $q=19$ & $-0.22424$ & $-0.29695$ & $-0.24044$ & $-0.04846$ & $-0.02748$\\
PBE & $q=11$ & $-0.21240$ & $-0.28794$ & $-0.23187$ & $-0.04235$ & $-0.02404$\\
    & $q=19$ & $-0.21250$ & $-0.28787$ & $-0.23182$ & $-0.04239$ & $-0.02405$\\
\bottomrule
\end{tabular}
\end{center}
(Eigenvalues in Ha of the neutral, spherical, non-spin-polarized
pseudo-atom, $5d^{10}6s^1$ with $6p$ empty; radial grid $h = 0.02$,
$r_{\max} = 16$\,Bohr.) The 19-electron atoms have in addition the
semicore levels $5s$ at $-3.988$ ($-3.973$)\,Ha and $5p_{1/2}$,
$5p_{3/2}$ at $-2.655$, $-2.042$ ($-2.639$, $-2.034$)\,Ha for LDA (PBE):
a spin--orbit splitting of the $5p$ shell of $16.7$ ($16.5$)\,eV. What
has \emph{not} been done with the 19-electron potentials is a
plane-wave calculation of the solid; the cutoff advice above is a
starting point, not a tested value.


\end{document}
